\documentclass[11pt]{article}

\usepackage[margin=1in]{geometry}
\usepackage{amsmath,amssymb,mathtools,bm}
\usepackage{booktabs,array}
\usepackage{graphicx}
\usepackage[dvipsnames]{xcolor}
\usepackage{microtype}
\usepackage{enumitem,needspace}
\usepackage{xurl}
\usepackage{hyperref}

\hypersetup{
  colorlinks=true,
  linkcolor=MidnightBlue,
  urlcolor=MidnightBlue,
  pdftitle={Space--Time Algebra},
  pdfauthor={Dr. Luke A. Wendt}
}

\setlength{\parindent}{0pt}
\setlength{\parskip}{0.4em}
\setlength{\abovedisplayskip}{0.55em}
\setlength{\belowdisplayskip}{0.55em}
\setlength{\abovedisplayshortskip}{0.35em}
\setlength{\belowdisplayshortskip}{0.35em}
\setlist{nosep,leftmargin=2em}
\numberwithin{equation}{section}
\setcounter{tocdepth}{2}
\makeatletter
\let\sta@tocsection\l@section
\renewcommand{\l@section}[2]{\Needspace{4\baselineskip}\sta@tocsection{#1}{#2}}
\makeatother

\newcommand{\R}{\mathbb{R}}
\newcommand{\C}{\mathbb{C}}
\newcommand{\ii}{\mathrm{i}}
\newcommand{\bsig}{\bm{\sigma}}
\newcommand{\bA}{\bm{A}}
\newcommand{\bB}{\bm{B}}
\newcommand{\bC}{\bm{C}}
\newcommand{\bV}{\bm{V}}
\newcommand{\br}{\bm{r}}
\newcommand{\bv}{\bm{v}}
\newcommand{\bu}{\bm{u}}
\newcommand{\bn}{\bm{n}}
\newcommand{\ba}{\bm{a}}
\newcommand{\bk}{\bm{k}}
\newcommand{\bg}{\bm{g}}
\newcommand{\bb}{\bm{b}}
\newcommand{\bq}{\bm{q}}
\newcommand{\bS}{\bm{S}}
\newcommand{\bF}{\bm{F}}
\newcommand{\bE}{\bm{E}}
\newcommand{\bBf}{\bm{B}}
\newcommand{\bJ}{\bm{J}}
\newcommand{\bp}{\bm{p}}
\newcommand{\bM}{\bm{M}}
\newcommand{\bN}{\bm{N}}
\newcommand{\bI}{\bm{I}}
\newcommand{\bPsi}{\bm{\Psi}}
\newcommand{\bpsi}{\bm{\psi}}
\newcommand{\bNm}{\bm{N}}
\newcommand{\bDm}{\bm{D}}
\newcommand{\pdr}{\partial_{\br}}
\newcommand{\rev}[1]{#1^{\mathsf H}}
\newcommand{\sconj}[1]{#1^{*}}
\newcommand{\sr}[1]{#1^{*\mathsf H}}
\newcommand{\adj}[1]{\operatorname{adj}\!\left(#1\right)}
\newcommand{\scpart}[1]{\operatorname{sc}(#1)}
\newcommand{\vecpart}[1]{\operatorname{vec}(#1)}
\newcommand{\repart}[1]{\operatorname{re}(#1)}
\newcommand{\impart}[1]{\operatorname{im}(#1)}
\newcommand{\Det}[1]{\det\left(#1\right)}
\newcommand{\Wop}{\bm{W}}
\newcommand{\dd}{\,\mathrm d}

\title{\textbf{Space--Time Algebra}\\[0.35em]
  \large A compact development of paravectors and their physical applications}
\author{Dr. Luke A. Wendt}
\date{}

\begin{document}
\maketitle

\begin{abstract}
This document develops a three-generator sigma algebra and its complex
paravectors, then applies them to relativity, electromagnetism, and quantum
mechanics.  Natural units $\hbar=c=1$, rationalized electromagnetic units,
and signature $(+,-,-,-)$ are used.
\end{abstract}

\tableofcontents
\clearpage

\section{The algebra at a glance}
\label{sec:algebra}

The generators $\sigma_n$, $n=1,2,3$, obey ordered multiplication rules.
Their products generally do not commute.  The panel gives their defining
relations and the general complex element developed below.

\begin{center}
\setlength{\fboxsep}{1.2em}
\fcolorbox{MidnightBlue}{MidnightBlue!4}{%
\begin{minipage}{0.88\textwidth}
\begin{align*}
 \sigma_n^2&=1,
 &\sigma_n\sigma_m&=-\sigma_m\sigma_n\quad(n\ne m),\\
 \ii&=\sigma_1\sigma_2\sigma_3,
 &\ii^2&=-1,\\
 \bsig&=\{\sigma_1,\sigma_2,\sigma_3\},
 &\ii\bsig&=\{\sigma_2\sigma_3,\sigma_3\sigma_1,\sigma_1\sigma_2\},\\
 Z&=(a+b\ii)+(\bA+\bB\ii)\cdot\bsig,
 &a,b&\in\R,\quad \bA,\bB\in\R^3.
\end{align*}
\end{minipage}}
\end{center}

Ordinary three-vector operations use
\begin{align*}
 \bA\cdot\bB&=A_1B_1+A_2B_2+A_3B_3,\\
 \bA\times\bB&=\{A_2B_3-A_3B_2,
 A_3B_1-A_1B_3,A_1B_2-A_2B_1\},\\
 \bA^2&:=\bA\cdot\bA.
\end{align*}
Thus the square of a vector always denotes its self dot product.
The elementary Cartesian vectors used below are the real vectors
\[
 \bm e_1=\{1,0,0\},\qquad
 \bm e_2=\{0,1,0\},\qquad
 \bm e_3=\{0,0,1\}.
\]
Thus $\bm e_k\cdot\bsig=\sigma_k$; the scalar basis element is $1$.
Unindexed $e$ denotes only Euler's number.

These eight elements form a complete basis, displayed by type as
\[
 \begin{array}{cccc}
 1&\sigma_1,\sigma_2,\sigma_3&\ii\sigma_1,\ii\sigma_2,\ii\sigma_3&\ii\\
 \text{scalar}&\text{vectors}&\text{bivectors}&\text{pseudoscalar}
 \end{array}
\]
Anticommutation gives
$\ii^2=-\sigma_1^2\sigma_2^2\sigma_3^2=-1$.  The pseudoscalar is central:
moving $\sigma_n$ through the other two generators produces two sign
changes, so
\begin{equation*}
 \ii\sigma_n=\sigma_n\ii.
\end{equation*}
Thus its order of multiplication does not matter.

For algebra elements $X,Y$, the commutator is
\begin{gather*}
 [X,Y]:=XY-YX\\
 [Y,X]=-[X,Y].
\end{gather*}
Thus $[X,Y]=0$ exactly when $X$ and $Y$ commute.
For $n,m,k=1,2,3$, define the Kronecker delta by
\begin{equation*}
 \delta_{nm}:=\begin{aligned}[t]
  &1\quad\text{if }n=m,\\
  &0\quad\text{if }n\ne m,
 \end{aligned}
\end{equation*}
and the Levi--Civita symbol by
\begin{equation*}
 \epsilon_{nmk}:=\begin{aligned}[t]
  &+1\quad\text{if }(n,m,k)\text{ is an even permutation of }(1,2,3),\\
  &-1\quad\text{if }(n,m,k)\text{ is an odd permutation of }(1,2,3),\\
  &\phantom{+}0\quad\text{if two or more indices are equal}.
 \end{aligned}
\end{equation*}
For $n,m=1,2,3$, the ordered multiplication rule and commutator are
\begin{align}
  \sigma_n\sigma_m=\delta_{nm}
  +\sum_{k=1}^{3}\epsilon_{nmk}\,\ii\sigma_k
 \label{eq:sigma-product}\\
  [\sigma_n,\sigma_m]
  =2\sum_{k=1}^{3}\epsilon_{nmk}\,\ii\sigma_k. \notag
\end{align}
\section{Hermitian and sigma conjugation}
\label{sec:conjugate-operators}

Define \emph{Hermitian conjugation} by \eqref{eq:hermitian-definition} and
\emph{sigma conjugation} by \eqref{eq:sigma-conjugation-definition}:
\begin{align}
 (\,\cdot\,)^{\mathsf H}
 &:=\text{reverse the order of every product of sigma generators},
 \label{eq:hermitian-definition}\\
 (\,\cdot\,)^*
 &:=\text{replace every $\sigma_n$ by $-\sigma_n$ without changing product order}.
 \label{eq:sigma-conjugation-definition}
\end{align}
Thus, for two generators,
\begin{align*}
 (\sigma_m\sigma_n)^{\mathsf H}&=\sigma_n\sigma_m,\\
 (\sigma_m\sigma_n)^*&=(-\sigma_m)(-\sigma_n).
\end{align*}
Both operations act term by term on sums and leave real coefficients fixed.
Their actions on the four basis types are
\begin{center}
\begin{tabular}{c c c c c}
  \toprule
  type & scalar & vector & bivector & pseudoscalar\\
  representative & $1$ & $\sigma_n$ & $\ii\sigma_n$ & $\ii$\\
  \midrule
  $(\cdot)^*$ & $+$ & $-$ & $+$ & $-$\\
  $(\cdot)^{\mathsf H}$ & $+$ & $+$ & $-$ & $-$\\
  $(\cdot)^{*\mathsf H}$ & $+$ & $-$ & $-$ & $+$\\
  \bottomrule
\end{tabular}
\end{center}
For example, on a bivector the signs follow directly:
\begin{align*}
 (\ii\sigma_n)^*&=\ii^*\sigma_n^*=\ii\sigma_n,\\
 (\ii\sigma_n)^{\mathsf H}
 &=\sigma_n^{\mathsf H}\ii^{\mathsf H}=-\ii\sigma_n.
\end{align*}
For arbitrary algebra elements $Z,Z'$, the definitions give the product
laws
\begin{align}
 (ZZ')^{\mathsf H}&=Z'^{\mathsf H}Z^{\mathsf H},
 \label{eq:hermitian-product}\\
 (ZZ')^*&=Z^*Z'^*.
 \label{eq:sigma-conjugate-product}
\end{align}
The operations commute with each other: $Z^{*\mathsf H}=Z^{\mathsf H*}$.
Appendix~\ref{sec:matrix} gives both operations in the Pauli-matrix
representation, where $^{\mathsf H}$ is the usual matrix Hermitian conjugation.

\section{Paravectors}
\label{sec:paravectors}

\subsection{Core operations}
\label{sec:algebra-operations}

Let $a,b\in\R$ and $\bA,\bB\in\R^3$.  The real paravectors
$X=a+\bA\cdot\bsig$ and $Y=b+\bB\cdot\bsig$ form the general complex
paravector
\begin{equation}
 Z=X+\ii Y=(a+b\ii)+(\bA+\ii\bB)\cdot\bsig
   =S+\bV\cdot\bsig.
 \label{eq:paravector}
\end{equation}
Here $S=a+b\ii\in\C$ and $\bV=\bA+\ii\bB\in\C^3$ are its
complex scalar and vector coefficients.  Define their extraction functions by
\begin{align}
 \scpart Z&:=S,
 \label{eq:scalar-extraction}\\
 \vecpart Z&:=\bV.
 \label{eq:vector-extraction}
\end{align}
This notation applies throughout the algebra unless a real specialization
is stated.  On the coefficients, both conjugations give
$S^*=S^{\mathsf H}=a-b\ii$ and
$\bV^*=\bV^{\mathsf H}=\bA-\ii\bB$; on the paravector, the two
conjugations differ in the vector sign:
\begin{align}
 Z^*&=S^*-\bV^*\cdot\bsig
 =(a-b\ii)+(-\bA+\ii\bB)\cdot\bsig,
 \label{eq:algebra-star}\\
 Z^{\mathsf H}&=S^*+\bV^*\cdot\bsig
 =(a-b\ii)+(\bA-\ii\bB)\cdot\bsig.
 \label{eq:algebra-H}
\end{align}
Their composition defines the adjugate
\begin{equation}
 \adj Z:=Z^{*\mathsf H}=S-\bV\cdot\bsig.
 \label{eq:adjugate}
\end{equation}
The trace is twice the scalar part:
\begin{equation}
 \operatorname{tr}(Z):=Z+\adj Z=2\scpart Z=2S.
 \label{eq:trace-scalar}
\end{equation}
The determinant is the complex scalar
\begin{equation}
 \det(Z):=\adj Z\,Z=Z\adj Z=S^2-\bV^2.
 \label{eq:det}
\end{equation}
For complex $Z'$, adjugation reverses product order, giving determinant
multiplicativity:
\begin{align}
 \adj{ZZ'}&=\adj{Z'}\adj Z,
 \label{eq:adjugate-product}\\
 \det(ZZ')&=\adj{Z'}(\adj Z\,Z)Z'
           =\det(Z)\det(Z').
 \label{eq:det-product}
\end{align}
Since the determinants are scalars, the determinant of any product is
unchanged by reordering its factors.  For a complex scalar $c$,
\begin{align}
 \adj{cZ}&=c\adj Z, \notag\\
 \det(cZ)&=c^2\det(Z).
 \label{eq:det-scalar-scaling}
\end{align}
When $\det(Z)\ne0$, the inverse is
\begin{equation}
 Z^{-1}=\frac{\adj Z}{\det(Z)}
       =\frac{S-\bV\cdot\bsig}{S^2-\bV^2}.
 \label{eq:inverse}
\end{equation}
For $Z'=S'+\bV'\cdot\bsig$, with complex coefficients, division is
order-sensitive.  An unambiguous quotient exists when $Z$ is invertible
and the factors commute:
\begin{equation}
 \frac{Z'}{Z}:=Z'Z^{-1}=Z^{-1}Z'
 \qquad\text{if }\bV\times\bV'=0.
 \label{eq:commuting-quotient}
\end{equation}
Differentiating $Z^{-1}Z=1$ gives the inverse rule
\begin{equation}
 \dd(Z^{-1})=-Z^{-1}(\dd Z)Z^{-1}.
 \label{eq:inverse-differential}
\end{equation}
For complex $Z,Z',Z''$, scalar part obeys the cyclic laws
\begin{align}
 \scpart{ZZ'}&=\scpart{Z'Z},
 \label{eq:scalar-cyclic}\\
 \scpart{ZZ'Z''}&=\scpart{Z'Z''Z}=\scpart{Z''ZZ'}.
 \label{eq:scalar-cyclic-three}
\end{align}
Multiplying these identities by $2$ gives trace cyclicity.  Their similarity
invariants are derived in Section~\ref{sec:preserved-quantities}.

The real and imaginary extractions use Hermitian conjugation:
\begin{align}
 \repart Z&:=\tfrac12(Z+Z^{\mathsf H})=X,
 \label{eq:algebra-re}\\
 \impart Z&:=\frac{Z-Z^{\mathsf H}}{2\ii}=Y.
 \label{eq:algebra-im}
\end{align}
Adjugation evaluates the scalar and vector extractions explicitly:
\begin{align}
 \scpart Z&=\tfrac12(Z+\adj Z)=S,
 \label{eq:scalar-from-adjugate}\\
 (\vecpart Z)\cdot\bsig&=\tfrac12(Z-\adj Z)=\bV\cdot\bsig.
 \label{eq:vector-from-adjugate}
\end{align}
The functions $\operatorname{re}$ and $\operatorname{im}$ also act on
complex coefficients.  Composing extractions recovers all four real parts:
\begin{align*}
 a&=\repart S=\scpart{\repart Z},\\
 b&=\impart S=\scpart{\impart Z},\\
 \bA&=\repart{\bV}=\vecpart{\repart Z},\\
 \bB&=\impart{\bV}=\vecpart{\impart Z}.
\end{align*}
Appendix~\ref{sec:matrix} reproduces these intrinsic operations with Pauli
matrices.

\subsection{Unit vectors and paravectors}
\label{sec:unit-forms}

Recall $Z=S+\bV\cdot\bsig$, with $S\in\C$ and $\bV\in\C^3$.
For $\bC\in\C^3$, distinguish the square from the positive norm:
\begin{align*}
 \bC^2&:=\bC\cdot\bC,\\
 \lVert\bC\rVert^2&:=\bC\cdot\bC^*.
\end{align*}
For $S\in\C$, set $\lVert S\rVert^2:=SS^*$.  The positive paravector norm is
\begin{equation}
 \lVert Z\rVert^2:=\lVert S\rVert^2+\lVert\bV\rVert^2.
 \label{eq:coefficient-norm}
\end{equation}
All norms denote nonnegative real square roots.
Here $\bu\in\C^3$ and $Z$ is a complex paravector.
Each name below specifies a condition; a letter alone implies none.

\Needspace{7\baselineskip}
\begin{itemize}
\item \textbf{Vector of square one.}
\begin{equation}
 \bu^2=\bu\cdot\bu=1.
 \label{eq:bilinear-unit-vector}
\end{equation}
For $\bC^2\ne0$, take $\bu=\bC/\sqrt{\bC^2}$ with either square root.
Its norm need not be one.  A nonzero vector with $\bC^2=0$
cannot be normalized this way.
\item \textbf{Vector of unit norm.}
\begin{equation}
 \lVert\bu\rVert^2=\bu\cdot\bu^*=1.
 \label{eq:unit-norm-vector}
\end{equation}
For any $\bC\ne0$, take $\bu=\bC/\lVert\bC\rVert$; then
$\bu^2=\bC^2/\lVert\bC\rVert^2$ need not be one.
\phantomsection\label{def:real-unit-vector}
For real vectors the two conditions coincide: a \emph{real unit vector}
has $\bu^2=\lVert\bu\rVert^2=1$.  For complex vectors, specify which holds.
\Needspace{14\baselineskip}
\item \textbf{Unit paravector.} $Z$ has an inverse on both sides:
\begin{equation}
 ZZ^{-1}=Z^{-1}Z=1.
 \label{eq:unit-paravector}
\end{equation}
\item \textbf{Unitary paravector.} Its Hermitian conjugate is its inverse:
\begin{equation}
 Z^{\mathsf H}Z=1.
 \label{eq:unitary-transform}
\end{equation}
Then $Z^{\mathsf H}=Z^{-1}$ and $ZZ^{\mathsf H}=1$.
\item \textbf{Paravector of determinant one.}
\begin{equation}
 \det(Z)=1.
 \label{eq:unit-determinant-paravector}
\end{equation}
It is a unit with $Z^{-1}=\adj Z$; it need not be unitary.
For real $Z$, $\adj Z=Z^*$, so this condition is also $Z^*Z=1$.
This equivalence need not hold for complex $Z$.
\item \textbf{Paravector of unit norm.}
\begin{equation}
 \lVert Z\rVert^2=1.
 \label{eq:unit-norm-paravector}
\end{equation}
For $Z\ne0$, $Z'=Z/\lVert Z\rVert$ has unit norm.
It need not have an inverse: $(1+\bu\cdot\bsig)/\sqrt2$ has norm one
and determinant zero for any real unit vector $\bu$.
Unitarity implies unit norm by \eqref{eq:hermitian-norm}.
\item \textbf{Self-inverse paravector.} $Z$ is its own inverse:
\begin{equation}
 Z^2=1.
 \label{eq:involutory-paravector}
\end{equation}
Then $Z^{-1}=Z$; it is unitary exactly when $Z^{\mathsf H}=Z$.
Its parts obey $S^2+\bV^2=1$ and $2S\bV=0$:
thus $Z=\pm1$ or $Z=\bu\cdot\bsig$ with $\bu^2=1$.
\end{itemize}

\Needspace{9\baselineskip}
Two complex vectors illustrate why the normalizations differ:
\begin{center}
\begin{tabular}{ccc}
\toprule
$\bC$ & $\bC^2$ & $\lVert\bC\rVert^2$\\
\midrule
$\{\sqrt2,\ii,0\}$ & $1$ & $3$\\
$\{1,\ii,0\}$ & $0$ & $2$\\
\bottomrule
\end{tabular}
\end{center}
Their vectors of unit norm have squares $1/3$ and $0$, respectively.
The physical roles of the paravector conditions are compared in
Section~\ref{sec:preserved-quantities}.

\subsection{Products and quadratic forms}

Here $\bA,\bB\in\R^3$ are real coefficient vectors.
The sigma rule~\eqref{eq:sigma-product} gives the dot--cross identity
\begin{equation}
 \begin{aligned}
 (\bA\cdot\bsig)(\bB\cdot\bsig)
 &=\sum_{n=1}^{3}\sum_{m=1}^{3}
   A_nB_m\sigma_n\sigma_m\\
 &=\sum_{n=1}^{3}\sum_{m=1}^{3}A_nB_m\delta_{nm}
   +\sum_{n=1}^{3}\sum_{m=1}^{3}\sum_{k=1}^{3}
   A_nB_m\epsilon_{nmk}\ii\sigma_k\\
 &=\bA\cdot\bB+\ii(\bA\times\bB)\cdot\bsig.
 \end{aligned}
 \label{eq:vector-product}
\end{equation}
Adding the product with $\bA$ and $\bB$ interchanged gives
\begin{equation*}
 2(\bA\cdot\bB)
 =(\bA\cdot\bsig)(\bB\cdot\bsig)
 +(\bB\cdot\bsig)(\bA\cdot\bsig).
\end{equation*}
Consequently, for every nonnegative integer $n$,
\begin{align*}
 (\bA\cdot\bsig)^2&=\bA^2,\\
 (\bA\cdot\bsig)^{2n}&=(\bA^2)^n,\\
 (\bA\cdot\bsig)^{2n+1}&=(\bA^2)^n\bA\cdot\bsig.
\end{align*}
Also,
\begin{equation*}
 (\bA\cdot\bsig)\bsig=\bA-\ii\bA\times\bsig.
\end{equation*}
Subtracting the two orders gives
\begin{equation}
 [\bA\cdot\bsig,\bB\cdot\bsig]
 =2\ii(\bA\times\bB)\cdot\bsig.
 \label{eq:vector-commutator}
\end{equation}
Hence
\begin{align*}
 [\bA\cdot\bsig,\bsig]&=-2\ii\bA\times\bsig,\\
 [\bA\cdot\ii\bsig,\bsig]&=2\bA\times\bsig,\\
 [\bA\cdot\bsig,\bsig]\cdot\bB
 &=2\ii(\bA\times\bB)\cdot\bsig.
\end{align*}

\subsection{Quadratic identity and eigenvalues}

Recall $Z=S+\bV\cdot\bsig$, with $S\in\C$ and $\bV\in\C^3$.
Since $Z-S=\bV\cdot\bsig$ and
$(\bV\cdot\bsig)^2=\bV^2$,
\begin{align*}
 (Z-S)^2&=\bV^2,\\
 Z^2-2SZ+(S^2-\bV^2)&=0.
\end{align*}
Using $\operatorname{tr}(Z)=2S$ and
$\det(Z)=S^2-\bV^2$ gives the trace--determinant identity
\begin{equation}
 Z^2-\operatorname{tr}(Z)Z+\det(Z)=0.
  \label{eq:minimal}
\end{equation}
Solving \eqref{eq:minimal} for a scalar eigenvalue gives
\begin{equation}
 \lambda_\pm
 =\frac{\operatorname{tr}(Z)
 \pm\sqrt{\operatorname{tr}(Z)^2-4\det(Z)}}{2}
 =S\pm\sqrt{\bV^2}.
 \label{eq:eigenvalues}
\end{equation}
Because $\bV$ is generally complex, $\bV^2=\bV\cdot\bV$ and its square
root may also be complex.
The quadratic identity therefore factors as
\begin{equation*}
 (Z-\lambda_-)(Z-\lambda_+)=0.
\end{equation*}

\subsection{Real products}

For real $X=a+\bA\cdot\bsig$ and $Y=b+\bB\cdot\bsig$, with
$a,b\in\R$ and $\bA,\bB\in\R^3$, \eqref{eq:vector-product} gives
\begin{align}
 XY&=(ab+\bA\cdot\bB)
 +(a\bB+b\bA+\ii\bA\times\bB)\cdot\bsig,
 \label{eq:real-paravector-product}\\
 [X,Y]&=2\ii(\bA\times\bB)\cdot\bsig, \notag\\
 X^2&=(a^2+\bA^2)+2a\bA\cdot\bsig. \notag
\end{align}
Their scalar and vector parts are read directly from these equations.
The real quadratic identity is \eqref{eq:minimal} with $S=a$.
If $\bA\times\bB=0$, then $XY=YX$; when $\det(X)\ne0$, also
$X^{-1}Y=YX^{-1}$.

\subsection{Complex products}

Let $Z=S+\bV\cdot\bsig$ and $Z'=S'+\bV'\cdot\bsig$, where
$S,S'\in\C$ and $\bV,\bV'\in\C^3$.  Complex bilinearity extends
\eqref{eq:vector-product} to the general product
\begin{equation}
 ZZ'=(SS'+\bV\cdot\bV')+
 \bigl(S\bV'+S'\bV+\ii\bV\times\bV'\bigr)\cdot\bsig.
 \label{eq:para-product}
\end{equation}
Thus
\begin{align*}
 \scpart{ZZ'}&=SS'+\bV\cdot\bV',\\
 \vecpart{ZZ'}&=S\bV'+S'\bV+\ii\bV\times\bV'.
\end{align*}
Equivalently, write $Z=X+\ii Y$ and $Z'=X'+\ii Y'$ with all four
factors real.  Their products need not be real, but distributivity gives
\begin{equation*}
 ZZ'=(XX'-YY')+\ii(XY'+YX').
\end{equation*}

Write $Z=(a+b\ii)+(\bA+\ii\bB)\cdot\bsig$ with real
$a,b,\bA,\bB$, and use primes for the second paravector's real components.  The four
real components of $ZZ'$ are
\begin{align*}
 \text{real scalar: }&aa'-bb'+\bA\cdot\bA'-\bB\cdot\bB',\\
 \text{imaginary scalar: }&ab'+ba'+\bA\cdot\bB'+\bB\cdot\bA',\\
 \text{real vector: }&a\bA'-b\bB'+a'\bA-b'\bB
                    -\bA\times\bB'-\bB\times\bA',\\
 \text{imaginary vector: }&a\bB'+b\bA'+a'\bB+b'\bA
                    +\bA\times\bA'-\bB\times\bB'.
\end{align*}
Specializing \eqref{eq:para-product} gives the quadratic products:
\begin{align}
 Z^2&=(a^2-b^2+\bA^2-\bB^2)
 +2\ii(ab+\bA\cdot\bB) \notag\\
 &\quad+2(a\bA-b\bB)\cdot\bsig
 +2\ii(a\bB+b\bA)\cdot\bsig, \notag\\
 ZZ^*&=(a^2+b^2-\bA^2-\bB^2)
 -2(\bA\times\bB)\cdot\bsig
 +2\ii(a\bB-b\bA)\cdot\bsig,
 \label{eq:star-product-left}\\
 Z^*Z&=(a^2+b^2-\bA^2-\bB^2)
 +2(\bA\times\bB)\cdot\bsig
 +2\ii(a\bB-b\bA)\cdot\bsig,
 \label{eq:star-product-right}\\
 ZZ^{\mathsf H}&=a^2+b^2+\bA^2+\bB^2
 +2(a\bA+b\bB+\bA\times\bB)\cdot\bsig,
 \label{eq:hermitian-product-left}\\
 Z^{\mathsf H}Z&=a^2+b^2+\bA^2+\bB^2
 +2(a\bA+b\bB-\bA\times\bB)\cdot\bsig,
 \label{eq:hermitian-product-right}\\
 ZZ^{*\mathsf H}&=\det(Z)
 =a^2-b^2-\bA^2+\bB^2
 +2\ii(ab-\bA\cdot\bB). \notag
\end{align}
Equations~\eqref{eq:star-product-left} and \eqref{eq:star-product-right}
show $ZZ^*=Z^*Z$ exactly when $\bA\times\bB=0$.  Equations
\eqref{eq:hermitian-product-left} and \eqref{eq:hermitian-product-right}
give the same condition for $ZZ^{\mathsf H}=Z^{\mathsf H}Z$.

\subsubsection{Norm and determinant}
For the complex coefficients $S=a+b\ii$ and $\bV=\bA+\ii\bB$, the norms
defined above give
\begin{align*}
 \lVert S\rVert^2&=a^2+b^2,\\
 \lVert\bV\rVert^2&=\bA^2+\bB^2.
\end{align*}
For the positive norm~\eqref{eq:coefficient-norm}, the quadratic products
and cyclicity give the Hermitian form
\begin{equation}
 \begin{aligned}
 \lVert Z\rVert^2&=\scpart{Z^{\mathsf H}Z}
 =\tfrac12\operatorname{tr}(Z^{\mathsf H}Z)
 =\scpart{ZZ^{\mathsf H}}\\
 &=\scpart{Z^*Z^{*\mathsf H}}
 =\scpart{Z^{*\mathsf H}Z^*}.
 \end{aligned}
 \label{eq:hermitian-norm}
\end{equation}
For complex coefficients, $S^2$ and $\bV^2$ generally differ from these
positive norms.  Also $\det(Z)$ may be complex.  For a complex scalar
$c=r e^{\ii\theta}$, with real $r\ge0$ and $\theta$, scalar scaling gives
\begin{align}
 \det(cZ)&=r^2e^{2\ii\theta}\det(Z), \notag\\
 \lVert cZ\rVert^2&=\lVert c\rVert^2\lVert Z\rVert^2
 =r^2\lVert Z\rVert^2.
 \label{eq:norm-scalar-scaling}
\end{align}
The determinant retains the scalar phase by \eqref{eq:det-scalar-scaling};
the norm loses it by \eqref{eq:norm-scalar-scaling}.

\subsection{Projection operators}
\label{sec:algebra-projectors}

For a \hyperref[def:real-unit-vector]{real unit vector} $\bu$, the
paravector $\bu\cdot\bsig$ is Hermitian and
\hyperref[eq:involutory-paravector]{self-inverse}.  Define its two orthogonal projectors
\begin{equation}
 \Pi_\pm:=\tfrac12(1\pm\bu\cdot\bsig).
 \label{eq:projectors}
\end{equation}
Since $(\bu\cdot\bsig)^2=1$,
\begin{align*}
 \Pi_\pm^2&=\Pi_\pm,\\
 \Pi_+\Pi_-&=0,\\
 \Pi_++\Pi_-&=1,\\
 \Pi_\pm^{\mathsf H}&=\Pi_\pm,\\
 \Pi_\pm^*&=\Pi_\mp,\\
 (\bu\cdot\bsig)\Pi_\pm&=\pm \Pi_\pm.
\end{align*}
Thus $\bu\cdot\bsig=\Pi_+-\Pi_-$.  For a second
\hyperref[def:real-unit-vector]{real unit vector} $\bu'$,
\begin{align*}
 (\bu\cdot\bsig)(\bu'\cdot\bsig)
 &=\bu\cdot\bu'+\ii(\bu\times\bu')\cdot\bsig,\\
 [\bu\cdot\bsig,\bu'\cdot\bsig]
 &=2\ii(\bu\times\bu')\cdot\bsig,\\
 (\bu\cdot\bsig)(\bu'\cdot\bsig)(\bu\cdot\bsig)
 &=\left(2(\bu\cdot\bu')\bu-\bu'\right)\cdot\bsig.
\end{align*}
Define $\bu'_{\parallel}=(\bu'\cdot\bu)\bu$ and
$\bu'_{\perp}=\bu'-\bu'_{\parallel}$.  Then
\begin{align*}
 \Pi_\pm(\bu'\cdot\bsig)\Pi_\pm
 &=\pm(\bu\cdot\bu')\Pi_\pm,\\
 \Pi_\pm(\bu'\cdot\bsig)\Pi_\mp
 &=\tfrac12(\bu'_{\perp}\pm\ii(\bu\times\bu'))\cdot\bsig\,
 \Pi_\mp.
\end{align*}
For $\Pi_\pm'=(1\pm\bu'\cdot\bsig)/2$, all four products are
\begin{align*}
 \Pi_\pm\Pi_\pm'&=\tfrac14(1+\bu\cdot\bu')
 +\tfrac14(\pm\bu\pm\bu'+\ii\bu\times\bu')\cdot\bsig,\\
 \Pi_\pm\Pi_\mp'&=\tfrac14(1-\bu\cdot\bu')
 +\tfrac14(\pm\bu\mp\bu'-\ii\bu\times\bu')\cdot\bsig.
\end{align*}

For arbitrary real vector increments $\dd\bu$ and $\dd\bu'$, the same
dot--cross decomposition gives the geometric determinant identity
\begin{equation}
 \Det{(\dd\bu\cdot\bsig)(\dd\bu'\cdot\bsig)}
 =(\dd\bu\cdot\dd\bu')^2
 +\frac14\left\lVert
 [\dd\bu\cdot\bsig,\dd\bu'\cdot\bsig]\right\rVert^2.
 \label{eq:lagrange-determinant}
\end{equation}
Dropping the nonnegative dot-product square in
\eqref{eq:lagrange-determinant} gives
\begin{equation}
 \sqrt{\Det{(\dd\bu\cdot\bsig)(\dd\bu'\cdot\bsig)}}
 \geq\frac12\left\lVert
 [\dd\bu\cdot\bsig,\dd\bu'\cdot\bsig]\right\rVert.
 \label{eq:lagrange-inequality}
\end{equation}
Equation~\eqref{eq:lagrange-determinant} is Lagrange's dot--cross identity;
\eqref{eq:lagrange-inequality} is a geometric bound, not a quantum uncertainty law.

\subsubsection{Ladder operators}
\label{sec:ladder-operators}

For arbitrary real unit vectors $\bu,\ba$ with $\bu\cdot\ba=0$,
define complex ladder paravectors and the real projectors of
\eqref{eq:projectors}:
\begin{align}
 Z_\pm&=\tfrac12\bigl(\ba\pm\ii\bu\times\ba\bigr)\cdot\bsig
 \label{eq:ladder-operators}\\
 \Pi_\pm&=\tfrac12(1\pm\bu\cdot\bsig). \notag
\end{align}
Then
\begin{align*}
 [\bu\cdot\bsig,Z_\pm]&=\pm2Z_\pm,\\
 Z_\pm^2&=0,\\
 Z_\pm Z_\mp&=\Pi_\pm,\\
 \ba\cdot\bsig&=Z_++Z_-,\\
 (\bu\times\ba)\cdot\bsig&=(Z_+-Z_-)/\ii,\\
 \bu\cdot\bsig&=\Pi_+-\Pi_-.
\end{align*}
The two signs give the maps between projected sectors:
\begin{align*}
 Z_\pm\Pi_\pm&=0,& Z_\pm\Pi_\mp&=Z_\pm,\\
 \Pi_\pm Z_\pm&=Z_\pm,& \Pi_\mp Z_\pm&=0.
\end{align*}

\subsection{Functions of paravectors}
\label{sec:spectral}

\subsubsection{Spectral rule}
Let $Z=S+\bV\cdot\bsig$, with $S\in\C$, $\bV\in\C^3$, and
$\bV^2\ne0$.  Define the complex scalars $w,\lambda_\pm$, complex vector $\bu$, and
paravector projectors $\Pi_\pm$ by
\begin{align*}
 w&:=+\sqrt{\bV^2},& \lambda_\pm&:=S\pm w,\\
 \bu&:=\bV/w,& \Pi_\pm&:=\tfrac12(1\pm\bu\cdot\bsig).
\end{align*}
The root is the principal scalar root.  Thus $\bu$ is a
\hyperref[eq:bilinear-unit-vector]{vector of square one};
$\bu\cdot\bsig$ is \hyperref[eq:involutory-paravector]{self-inverse}.
Its square gives
\begin{align*}
 \Pi_\pm^2&=\tfrac14(1\pm2\bu\cdot\bsig+1)=\Pi_\pm,\\
 \Pi_+\Pi_-&=\tfrac14(1-(\bu\cdot\bsig)^2)=0,\\
 \Pi_++\Pi_-&=1,\\
 Z&=S(\Pi_++\Pi_-)+w(\Pi_+-\Pi_-)
   =\lambda_+\Pi_++\lambda_-\Pi_-.
\end{align*}
These projectors need not be Hermitian.  The complementary-projector
identities give, for every
nonnegative integer $n$,
\[
 Z^n=\lambda_+^n\Pi_++\lambda_-^n\Pi_-
 =\tfrac12(\lambda_+^n+\lambda_-^n)
 +\tfrac12(\lambda_+^n-\lambda_-^n)\bu\cdot\bsig.
\]
For a convergent scalar power series $f(z)=\sum_{n=0}^\infty c_nz^n$,
with $c_n\in\C$, inserting these powers into $f(Z)=\sum_n c_nZ^n$
gives the spectral rule
\begin{equation}
 f(Z)=f(\lambda_+)\Pi_++f(\lambda_-)\Pi_-.
 \label{eq:complex-function-projectors}
\end{equation}
The same rule defines chosen analytic branches at the two eigenvalues.
Its proof uses only complementary idempotents and therefore applies to
any decomposition with these projector identities.

Adjugation exchanges the eigenvalues, so
\begin{align*}
 \adj Z&=\lambda_-\Pi_++\lambda_+\Pi_-,\\
 \det(Z)&=\lambda_+\lambda_-,\\
 Z^{-1}&=\frac{\adj Z}{\det(Z)}
       =\lambda_+^{-1}\Pi_++\lambda_-^{-1}\Pi_-.
\end{align*}
The inverse requires both eigenvalues to be nonzero.  Independently choosing
the signs of their principal square roots gives
\begin{equation}
 \sqrt Z=\pm\sqrt{\lambda_+}\Pi_+\pm\sqrt{\lambda_-}\Pi_-.
 \label{eq:paravector-square-root}
\end{equation}
There are four distinct roots when both eigenvalues are nonzero, and two
when one vanishes.  Changing $w$ to $-w$ exchanges the eigenvalue--projector
pairs and leaves the function unchanged when its branch choices are
exchanged with them.

\subsubsection{Real paravectors and direct square roots}
For real $X=a+\bA\cdot\bsig$, with $a\in\R$, $\bA\in\R^3$, and
$\bA\ne0$, the spectral quantities above reduce to
\[
 w=\lVert\bA\rVert,\qquad
 \bu=\bA/\lVert\bA\rVert,\qquad
 \lambda_\pm=a\pm\lVert\bA\rVert.
\]
The projectors are now Hermitian.  In particular, for a
\hyperref[def:real-unit-vector]{real unit vector} $\bu$,
\begin{align*}
 f(\bu\cdot\bsig)
 &=f(1)\Pi_++f(-1)\Pi_-\\
 &=\tfrac12(f(1)+f(-1))
  +\tfrac12(f(1)-f(-1))\bu\cdot\bsig,\\
 \sqrt{\bu\cdot\bsig}&=\pm\Pi_+\pm\ii\Pi_-.
\end{align*}
The two root signs are independent.  For a fixed analytic branch,
$f(X)\to f(a)$ as $\bA\to0$; independent root choices require the
scalar case below.

A direct component derivation also determines all \emph{real} roots.
Let $Y=b+\bB\cdot\bsig$ be real, with $b\in\R$, $\bB\in\R^3$,
and seek real $X=a+\bA\cdot\bsig$ such that $X^2=Y$.  Equating parts gives
\begin{align*}
 a^2+\bA^2&=b,\\
 2a\bA&=\bB.
\end{align*}
If $\bB\ne0$, then $a\ne0$ and $\bA=\bB/(2a)$.  Substitution gives
\begin{align}
 a^4-ba^2+\tfrac14\bB^2&=0, \notag\\
 a^2&=\tfrac12\left(b\pm\sqrt{b^2-\bB^2}\right).
 \label{eq:real-square-root-scalar}
\end{align}
For each nonnegative $a^2$ in \eqref{eq:real-square-root-scalar}, take both
scalar roots $a$ and set
$\bA=\bB/(2a)$; then $\bA^2=b-a^2$.  There are four real roots when
$b>\lVert\bB\rVert$, two when $b=\lVert\bB\rVert>0$, and none when
$b<\lVert\bB\rVert$.

If $\bB=0$, the equations instead give the scalar roots
$X=\pm\sqrt b$ and the pure-vector roots $X=\bA\cdot\bsig$ with
$\bA^2=b$ for $b>0$.  For $b=0$ only $X=0$ is real; for $b<0$ there
is no real root.

\subsubsection{Coalescing eigenvalues}
For complex $Z=S+\bV\cdot\bsig$ with $\bV^2=0$, define
$N=\bV\cdot\bsig$, so $N^2=0$.  A scalar function analytic at $S$
has the terminating Taylor expansion
\begin{equation}
 f(Z)=f(S)+f'(S)N.
 \label{eq:coalescing-function}
\end{equation}
For $S\ne0$, the chosen square root and logarithm are
\begin{align}
 \sqrt Z&=r+\frac{N}{2r},
 \label{eq:coalescing-square-root}\\
 \log Z&=\log S+\frac{N}{S}.
 \label{eq:coalescing-logarithm}
\end{align}
Here $r\in\C$ is a chosen scalar root of $S$ ($r^2=S$); use a consistent
logarithm branch.  A nonzero $N$ with $S=0$ has no square
root: such a root would have only zero eigenvalues, and its quadratic
identity would force its square to vanish.  It has no logarithm because
it is not invertible.  If $N=0$, a fixed analytic
branch gives $f(S)$; non-scalar square roots of a scalar must instead be
found from the component equations $a^2+\bA^2=S$, $2a\bA=0$, now allowing
$a$ and $\bA$ to be complex.  Independent branch choices need not remain
finite as distinct eigenvalues coalesce.

\subsubsection{Exponential: projection and series derivations}
For complex $Z=S+\bV\cdot\bsig$ with $\bV^2\ne0$, recall
$w=+\sqrt{\bV^2}$, the \hyperref[eq:bilinear-unit-vector]{vector of square one}
$\bu=\bV/w$, and $\lambda_\pm=S\pm w$.
The spectral rule~\eqref{eq:complex-function-projectors} gives
\begin{align*}
 \exp(Z)
 &=\tfrac12e^{S+w}(1+\bu\cdot\bsig)
  +\tfrac12e^{S-w}(1-\bu\cdot\bsig)\\
 &=e^S\left(\frac{e^w+e^{-w}}2
      +\frac{e^w-e^{-w}}2\bu\cdot\bsig\right).
\end{align*}
Independently, $S$ commutes with the vector part and
$(w\bu\cdot\bsig)^{2n}=w^{2n}$,
$(w\bu\cdot\bsig)^{2n+1}=w^{2n+1}\bu\cdot\bsig$.  Separating the
defining series therefore gives
\begin{align*}
 \exp(Z)&=e^S\sum_{n=0}^{\infty}
      \frac{(w\bu\cdot\bsig)^n}{n!}\\
 &=e^S\left(\sum_{n=0}^{\infty}\frac{w^{2n}}{(2n)!}
       +\sum_{n=0}^{\infty}\frac{w^{2n+1}}{(2n+1)!}\bu\cdot\bsig\right).
\end{align*}
Both derivations give
\begin{equation}
 \exp(Z)=e^S(\cosh w+\sinh w\,\bu\cdot\bsig).
 \label{eq:paravector-exponential}
\end{equation}
This is independent of the choice $w$ or $-w$: $\cosh$ is even, while
both $\sinh w$ and $\bu$ change sign.  When $\bV^2=0$,
\eqref{eq:coalescing-function} instead gives
\[
 \exp(S+\bV\cdot\bsig)=e^S(1+\bV\cdot\bsig).
\]
The defining series also gives
\begin{align*}
 \exp(Z^*)&=\exp(Z)^*,\\
 \exp(Z^{\mathsf H})&=\exp(Z)^{\mathsf H}.
\end{align*}

For the real components $a,b,\bA,\bB$ of
$Z=(a+b\ii)+(\bA+\ii\bB)\cdot\bsig$, write
\[
 w=+\sqrt{\bA^2-\bB^2+2\ii\bA\cdot\bB}.
\]
Substituting $\bu=(\bA+\ii\bB)/w$ into
\eqref{eq:paravector-exponential} gives the component form
\begin{equation}
 \exp(Z)=e^a(\cos b+\ii\sin b)
 \left(\cosh w+\frac{\sinh w}{w}(\bA+\ii\bB)\cdot\bsig\right).
 \label{eq:exponential-components}
\end{equation}
Its real and imaginary specializations are
\begin{align}
 \exp(a+\bA\cdot\bsig)
 &=e^a\left(\cosh\lVert\bA\rVert
  +\frac{\sinh\lVert\bA\rVert}{\lVert\bA\rVert}\bA\cdot\bsig\right),
 \label{eq:real-exponential}\\
 \exp\!\left(\ii(b+\bB\cdot\bsig)\right)
 &=(\cos b+\ii\sin b)\left(\cos\lVert\bB\rVert
  +\ii\frac{\sin\lVert\bB\rVert}{\lVert\bB\rVert}\bB\cdot\bsig\right).
 \label{eq:imaginary-exponential}
\end{align}
Setting $a=0$ in \eqref{eq:real-exponential} or $b=0$ in
\eqref{eq:imaginary-exponential} gives the pure-vector exponentials.
Every quotient uses its analytic limit at a zero argument.

\subsubsection{Rotation and boost factors}
\label{sec:algebra-rotation-boost}
Let $\theta\in\R$ and let $\bu\in\R^3$ be a
\hyperref[def:real-unit-vector]{real unit vector},
$\bu^2=1$.  The two special exponentials just derived define the complex
rotation factor $R$ and the real boost factor $L$:
\begin{equation}
 R(\theta,\bu):=\exp\!\left(\ii\frac\theta2\bu\cdot\bsig\right)
 =\cos\frac\theta2+\ii\sin\frac\theta2\bu\cdot\bsig.
 \label{eq:rotation-exponential}
\end{equation}
\begin{equation}
 L(\theta,\bu):=\exp\!\left(\frac\theta2\bu\cdot\bsig\right)
 =\cosh\frac\theta2+\sinh\frac\theta2\bu\cdot\bsig.
 \label{eq:boost-exponential}
\end{equation}
In \eqref{eq:rotation-exponential} and \eqref{eq:boost-exponential},
the parameter is a rotation angle for $R$ and a rapidity for $L$; both are
real dimensionless numbers.  Here these names specify algebraic objects.
Their actions on a general paravector are derived in
Section~\ref{sec:preserved-quantities}; their physical applications follow in
Section~\ref{sec:transformations}.

Because $(\bu\cdot\bsig)^2=1$, reversing the parameter gives the
inverses
\begin{align}
 R^{-1}&=R(-\theta,\bu)
 =\cos\frac\theta2-\ii\sin\frac\theta2\bu\cdot\bsig,
 \label{eq:rotation-inverse}\\
 L^{-1}&=L(-\theta,\bu)
 =\cosh\frac\theta2-\sinh\frac\theta2\bu\cdot\bsig.
 \label{eq:boost-inverse}
\end{align}
Sigma conjugation changes both $\ii$ and $\bsig$ in sign; Hermitian
conjugation changes $\ii$ but fixes a real vector.  Hence
\begin{align}
 R^*&=R, \label{eq:rotation-star}\\
 R^{\mathsf H}&=R^{-1}, \label{eq:rotation-H}\\
 L^*&=L^{-1}, \label{eq:boost-star}\\
 L^{\mathsf H}&=L. \label{eq:boost-H}
\end{align}
Equation~\eqref{eq:rotation-H} makes $R$
\hyperref[eq:unitary-transform]{unitary}; \eqref{eq:boost-H}
makes $L$ Hermitian, so $L^{\mathsf H}L=L^2$.  Their determinants and
positive norms follow from the circular and hyperbolic identities:
\begin{align}
 \det(R)&=\cos^2(\theta/2)+\sin^2(\theta/2)=1,
 \label{eq:rotation-determinant}\\
 \det(L)&=\cosh^2(\theta/2)-\sinh^2(\theta/2)=1,
 \label{eq:boost-determinant}\\
 \lVert R\rVert^2&=1, \label{eq:rotation-norm}\\
 \lVert L\rVert^2&=\cosh\theta. \label{eq:boost-norm}
\end{align}
Equations~\eqref{eq:rotation-determinant} and \eqref{eq:boost-determinant}
make both \hyperref[eq:unit-determinant-paravector]{paravectors of determinant one},
so their adjugates equal their inverses.  Equations~\eqref{eq:rotation-norm}
and \eqref{eq:boost-norm} distinguish their positive norms; only $R$ is
unitary for nonzero $\theta$.  The half-angle gives the sign and period:
\begin{align}
 R(\theta+2\pi,\bu)&=-R(\theta,\bu),
 \label{eq:rotation-two-pi}\\
 R(\theta+4\pi,\bu)&=R(\theta,\bu).
 \label{eq:rotation-four-pi}
\end{align}

\paragraph{Gamma form.}
Recall $L=\exp(\theta\bu\cdot\bsig/2)$ from
\eqref{eq:boost-exponential}.  Define its real velocity parameter and
gamma factor by
\begin{align}
 \beta&:=\tanh\theta,
 \label{eq:boost-beta}\\
 \gamma&:=\cosh\theta=\frac1{\sqrt{1-\beta^2}}.
 \label{eq:boost-gamma}
\end{align}
Here $-1<\beta<1$, $\gamma\ge1$, and $\gamma\beta=\sinh\theta$.
The half-angle identities give
\begin{align*}
 \cosh\frac\theta2&=\sqrt{\frac{\gamma+1}{2}},\\
 \sinh\frac\theta2&=\frac{\gamma\beta}{\sqrt{2(\gamma+1)}}.
\end{align*}
Hence the boost factor itself is
\begin{equation}
 L=\sqrt{\frac{\gamma+1}{2}}
 \left(1+\frac{\gamma\beta}{\gamma+1}\bu\cdot\bsig\right).
 \label{eq:boost-gamma-form}
\end{equation}
Squaring doubles the exponent, so the full-rapidity paravector is
\begin{equation}
 L^2=\exp(\theta\bu\cdot\bsig)
 =\gamma(1+\beta\bu\cdot\bsig).
 \label{eq:boost-square}
\end{equation}
Equation~\eqref{eq:boost-square} has prefactor $\gamma$; $L$ has the half-angle
prefactor in \eqref{eq:boost-gamma-form}.  Equivalently, for any real vector
$\bA\ne0$, put $\bv=\tanh\lVert\bA\rVert\,\bA/\lVert\bA\rVert$ and
$\gamma=\cosh\lVert\bA\rVert$.  Then
\[
 \exp(\bA\cdot\bsig)=\gamma(1+\bv\cdot\bsig).
\]
The zero-vector case follows by continuity.  This last real paravector has
\hyperref[eq:unit-determinant-paravector]{determinant one} and becomes proper velocity in
Section~\ref{sec:kinematics}.

\subsubsection{Cayley transform}
Let $Y=b+\bB\cdot\bsig$ be real, with $b\in\R$, $\bB\in\R^3$, and
$0<\lVert\bB\rVert<b$.  Set the \hyperref[def:real-unit-vector]{real unit vector}
$\bu=\bB/\lVert\bB\rVert$ and recall
$\Pi_\pm=(1\pm\bu\cdot\bsig)/2$.  Its sigma conjugate
$Y^*=b-\bB\cdot\bsig=\adj Y$ commutes with $Y$, so
\[
 \frac{Y}{Y^*}=\frac{Y^2}{\det(Y)}
 =\frac{b^2+\bB^2+2b\bB\cdot\bsig}{b^2-\bB^2}.
\]
Define the real parameter $\theta=2\tanh^{-1}(\lVert\bB\rVert/b)$.
The double-angle identities give
\begin{align*}
 \cosh\theta&=\frac{b^2+\bB^2}{b^2-\bB^2},\\
 \sinh\theta&=\frac{2b\lVert\bB\rVert}{b^2-\bB^2}.
\end{align*}
Thus the quotient is the full-rapidity exponential
\begin{equation}
 \begin{aligned}
 \frac{b+\bB\cdot\bsig}{b-\bB\cdot\bsig}
 &=\cosh\theta+\sinh\theta\,\bu\cdot\bsig\\
 &=\exp(\theta\bu\cdot\bsig)
 =e^\theta\Pi_++e^{-\theta}\Pi_-.
 \end{aligned}
 \label{eq:cayley-exponential}
\end{equation}
Equation~\eqref{eq:cayley-exponential} is $L(\theta,\bu)^2$ from
\eqref{eq:boost-square}.  Its real generator
vector is $\bA=\theta\bu$, while the Cayley vector obeys
\begin{equation}
 \bB=b\tanh\!\left(\frac{\lVert\bA\rVert}{2}\right)
       \frac{\bA}{\lVert\bA\rVert}.
 \label{eq:cayley-generator-map}
\end{equation}
Thus $\lVert\bB\rVert/b=\tanh(\theta/2)$, the half-rapidity
parameter.  The corresponding real gamma factor and vector are
\begin{align*}
 \gamma&=\frac{b^2+\bB^2}{b^2-\bB^2},\\
 \bv&=\frac{2b\bB}{b^2+\bB^2}.
\end{align*}
They give the equivalent forms
\[
 \exp(\bA\cdot\bsig)=\gamma(1+\bv\cdot\bsig)
 =\frac{b+\bB\cdot\bsig}{b-\bB\cdot\bsig}=\frac{Y}{Y^*}.
\]
Choosing $b=1$ gives the Cayley form with scalar part one; $\bB=0$ follows by
continuity.

\subsubsection{Composing the factors}
Recall $L(\theta,\bu)=\exp(\theta\bu\cdot\bsig/2)$ and
$R(\theta,\bu)=\exp(\ii\theta\bu\cdot\bsig/2)$ from
\eqref{eq:boost-exponential} and \eqref{eq:rotation-exponential}.
For a fixed \hyperref[def:real-unit-vector]{real unit direction} $\bu$
and real parameters $\theta,\theta'$, the
generators commute, so their parameters add:
\begin{equation}
 L(\theta,\bu)L(\theta',\bu)=L(\theta+\theta',\bu).
 \label{eq:boost-composition}
\end{equation}
\begin{align*}
 R(\theta,\bu)R(\theta',\bu)&=R(\theta+\theta',\bu),\\
 L(\theta,\bu)R(\theta',\bu)&=R(\theta',\bu)L(\theta,\bu).
\end{align*}
For the aligned real and imaginary generator, put
$w=(\theta+\ii\theta')/2\in\C$.  Its exponential is
\[
 \exp(w\bu\cdot\bsig)
 =L(\theta,\bu)R(\theta',\bu)
 =\cosh w+\sinh w\,\bu\cdot\bsig.
\]
Expanding the complex scalar coefficients gives
\begin{align*}
 \cosh w&=\cosh\frac\theta2\cos\frac{\theta'}2
 +\ii\sinh\frac\theta2\sin\frac{\theta'}2,\\
 \sinh w&=\sinh\frac\theta2\cos\frac{\theta'}2
 +\ii\cosh\frac\theta2\sin\frac{\theta'}2.
\end{align*}
This aligned factorization requires parallel real and imaginary generator
vectors.  For different \hyperref[def:real-unit-vector]{real unit directions}
$\bu,\bv$, the vector product
instead gives
\begin{equation}
 \left[\tfrac12\bu\cdot\bsig,\tfrac12\bv\cdot\bsig\right]
 =\tfrac\ii2(\bu\times\bv)\cdot\bsig.
 \label{eq:boost-generator-commutator}
\end{equation}
This commutator is a rotation generator:
boosts about different axes do not commute.

\subsubsection{Logarithm}
For complex $Z=S+\bV\cdot\bsig$ with $\bV^2\ne0$, recall
$w=+\sqrt{\bV^2}$, the \hyperref[eq:bilinear-unit-vector]{vector of square one}
$\bu=\bV/w$, $\lambda_\pm=S\pm w$, and
$\Pi_\pm=(1\pm\bu\cdot\bsig)/2$.  Assume both eigenvalues are nonzero.
For a nonzero complex scalar $\lambda$ and any integer $n$, the scalar branches are
\[
 \log\lambda=\log\lVert\lambda\rVert+\ii(\arg\lambda+2n\pi).
\]
Choose one branch for each eigenvalue.  The spectral rule gives
\begin{equation}
 \begin{aligned}
 \log Z&=\log(\lambda_+)\Pi_++\log(\lambda_-)\Pi_-\\
 &=\tfrac12\bigl(\log\lambda_++\log\lambda_-\bigr)
 +\tfrac12\bigl(\log\lambda_+-\log\lambda_-\bigr)\bu\cdot\bsig.
 \end{aligned}
 \label{eq:paravector-logarithm}
\end{equation}
Exchanging $w$ and $-w$ exchanges both eigenvalue--projector pairs and
preserves the same chosen logarithm.  Independent scalar branches give
the different logarithms.

For a \hyperref[def:real-unit-vector]{real unit vector} $\bu$,
\eqref{eq:paravector-logarithm} uses
eigenvalues $1$ and $-1$, giving, for integers $n,m$,
\begin{equation}
 \log(\bu\cdot\bsig)
 =2n\pi\ii\Pi_++(2m+1)\pi\ii\Pi_-.
 \label{eq:unit-vector-logarithm}
\end{equation}
For real $Y=b+\bB\cdot\bsig$, with $b\in\R$, $\bB\in\R^3$, and
$b>\lVert\bB\rVert$, both eigenvalues are positive.  Combining their real
logarithms yields
\begin{equation}
 \log Y=\tfrac12\log(b^2-\bB^2)
 +\tanh^{-1}\!\left(\frac{\lVert\bB\rVert}{b}\right)
       \frac{\bB}{\lVert\bB\rVert}\cdot\bsig.
 \label{eq:positive-real-logarithm}
\end{equation}
The vector quotient uses its limit at $\bB=0$; the general coalescing case
was treated above.

\section{Transformations and invariants}
\label{sec:preserved-quantities}

Here $Z=S+\bV\cdot\bsig$ is a complex paravector: $S\in\C$ and
$\bV\in\C^3$.  Transformation factors are constant complex paravectors.
Distinguish the scalar part $\scpart Z=S=\tfrac12\operatorname{tr}Z$
from \eqref{eq:trace-scalar}, the complex determinant $S^2-\bV^2$ from
\eqref{eq:det}, and the positive norm $\scpart{Z^{\mathsf H}Z}$ from
\eqref{eq:hermitian-norm}.  Scalar part and trace carry exactly the same
information; we use $\operatorname{sc}$ below.

\subsection{Similarity, determinant, and norm}

For a complex \hyperref[eq:unit-paravector]{unit paravector} $T$,
similarity means $Z'=T^{-1}ZT$.
Scalar-part cyclicity gives
\begin{equation}
 \scpart{Z'}=\scpart{T^{-1}ZT}
 =\scpart{ZTT^{-1}}=\scpart Z.
 \label{eq:scalar-similarity}
\end{equation}
Thus trace is also preserved.  Multiplicativity gives
\begin{equation}
 \det(Z')=\det(T^{-1})\det(Z)\det(T)=\det(Z).
 \label{eq:det-similarity}
\end{equation}
The opposite similarity $Z'=TZT^{-1}$ has the same invariants.
For independent complex factors $T,T'$, the law $Z'=T'ZT$ instead gives
\begin{equation}
 \det(Z')=\det(T')\det(T)\det(Z).
 \label{eq:two-factor-determinant}
\end{equation}
This preserves the determinant for
every $Z$ exactly when $\det(T')\det(T)=1$.  This alone does not preserve
scalar part, norm, or the real paravector subspace.

If the factor $T$ is unitary, then $T^{\mathsf H}=T^{-1}$ by
\eqref{eq:unitary-transform}.  For $Z'=T^{\mathsf H}ZT$, cancellation and
scalar-part cyclicity give
\begin{equation}
 \begin{aligned}
 \lVert Z'\rVert^2
 &=\scpart{T^{\mathsf H}Z^{\mathsf H}TT^{\mathsf H}ZT}\\
 &=\scpart{Z^{\mathsf H}ZTT^{\mathsf H}}=\lVert Z\rVert^2.
 \end{aligned}
 \label{eq:unitary-norm-invariance}
\end{equation}
The norm is also preserved by $Z'=TZ$, $Z'=ZT$, and
$Z'=TZT^{\mathsf H}$.  The two similarities additionally preserve scalar
part and determinant; a one-sided unitary product need not.

\subsubsection{Unit determinant is not Hermitian unitarity}

The word \emph{unit} in \eqref{eq:unit-paravector} means invertible;
\emph{unitary} means \eqref{eq:unitary-transform}.
A useful third condition is \hyperref[eq:unit-determinant-paravector]{determinant one}.
To see their difference,
let $\dd X=\dd t+\dd\br\cdot\bsig$ be a real paravector increment, with
$\dd t>\lVert\dd\br\rVert$, and define the positive real scalar
$\dd s=\sqrt{\dd t^2-\dd\br^2}$.  Reserve $U$ for the real proper
velocity $U=\dd X/\dd s$.  Since $\dd X^*\dd X=\dd s^2$,
\begin{equation}
 U^*U=\frac{\dd X^*\dd X}{\dd s^2}=1.
 \label{eq:proper-velocity-star-normalization}
\end{equation}
Thus $\det(U)=1$ and $U^{-1}=U^*$: $U$ is a unit paravector with
$UU^{-1}=U^{-1}U=1$.  Its letter alone implies no algebraic property.
The two conjugates are
\begin{align*}
 U^{\mathsf H}&=U=\frac{\dd t+\dd\br\cdot\bsig}{\dd s},\\
 U^*&=\frac{\dd t-\dd\br\cdot\bsig}{\dd s}.
\end{align*}
Hermitian conjugation leaves every real paravector fixed; sigma conjugation
reverses its vector part.  Therefore
\begin{align}
 U^{\mathsf H}U&=U^2
 =\frac{\dd t^2+\dd\br^2+2\dd t\,\dd\br\cdot\bsig}{\dd s^2},
 \label{eq:proper-velocity-hermitian-square}\\
 \lVert U\rVert^2&=\frac{\dd t^2+\dd\br^2}{\dd t^2-\dd\br^2}.
 \label{eq:proper-velocity-positive-norm}
\end{align}
Both expressions equal $1$ only when $\dd\br=0$, its rest frame.
Section~\ref{sec:kinematics} develops the kinematics of this proper velocity.

\subsection{Real paravectors and the Lorentz metric}

Let $X=a+\bA\cdot\bsig$ be real, with $a\in\R$ and
$\bA\in\R^3$.  Equation~\eqref{eq:det} gives
$\det(X)=a^2-\bA^2$, with Lorentz signature;
\eqref{eq:coefficient-norm} gives $\lVert X\rVert^2=a^2+\bA^2$,
with Euclidean signature.  A transformation of real spacetime paravectors must
preserve both reality and the determinant.  For a complex
\hyperref[eq:unit-paravector]{unit paravector} $T$ with $\det(T)=1$,
these requirements hold for the Hermitian congruence
\begin{equation}
 \boxed{X'=T^{-1}X(T^{-1})^{\mathsf H}.}
 \label{eq:general-lorentz-congruence}
\end{equation}
Reality and determinant are preserved because
\begin{align*}
 X'^{\mathsf H}&=T^{-1}X^{\mathsf H}(T^{-1})^{\mathsf H}=X',\\
 \det(X')&=\det(T^{-1})\det(X)\det((T^{-1})^{\mathsf H})
 =\det(X).
\end{align*}
These congruences give the proper, future-preserving Lorentz transformations.
Future $X$ has nonnegative eigenvalues $a\pm\lVert\bA\rVert$ and
a real square root.  Its image satisfies
$\scpart{X'}=\lVert T^{-1}\sqrt X\rVert^2\geq0$.
Together with determinant invariance, this preserves the future cone.  The factors
$T$ and $-T$ give the same transformation~\cite[Secs.~II and IV]{baylis2004}.

For real $\theta$ and a \hyperref[def:real-unit-vector]{real unit vector} $\bu$,
the algebraic definitions
$R=\exp(\ii\theta\bu\cdot\bsig/2)$ and
$L=\exp(\theta\bu\cdot\bsig/2)$ reduce this congruence to
$X'=R^{\mathsf H}XR$ and $X'=L^{-1}XL^{-1}$, respectively.
Their components are derived in the next two subsections.
For $X=1$, the invariants differ explicitly:
\begin{align}
 X'&=L^{-2}=\cosh\theta-\sinh\theta\bu\cdot\bsig,
 \label{eq:boost-time-axis}\\
 \det(X')&=1, \notag\\
 \scpart{X'}&=\cosh\theta, \notag\\
 \lVert X'\rVert^2&=\cosh(2\theta). \notag
\end{align}
Thus the determinant supplies the spacetime metric, while the object being
transformed determines whether to use congruence or similarity.

\subsection{Rotation action and Rodrigues' formula}
\label{sec:algebra-rotation-action}

Recall the unitary factor
$R=\exp(\ii\theta\bu\cdot\bsig/2)$ from
\eqref{eq:rotation-exponential}, with real angle $\theta$ and
\hyperref[def:real-unit-vector]{real unit axis} $\bu$.
For the complex paravector $Z=S+\bV\cdot\bsig$, define
\begin{equation}
 Z'=R^{\mathsf H}ZR=R^{-1}ZR.
 \label{eq:rotation-action}
\end{equation}
This keeps $S'=S$.  The vector product
\eqref{eq:vector-product} gives
\begin{align*}
 (\bu\cdot\bsig)(\bV\cdot\bsig)(\bu\cdot\bsig)
 &=2(\bu\cdot\bV)\bu\cdot\bsig-\bV\cdot\bsig,\\
 [\bV\cdot\bsig,\bu\cdot\bsig]
 &=-2\ii(\bu\times\bV)\cdot\bsig.
\end{align*}
Inserting the sine and cosine forms of $R$ and $R^{-1}$ therefore yields
\begin{align*}
 \bV'\cdot\bsig
 &={}\cos^2\frac\theta2\,\bV\cdot\bsig
 +2\sin\frac\theta2\cos\frac\theta2
       (\bu\times\bV)\cdot\bsig\\
 &\quad+\sin^2\frac\theta2
 \left(2(\bu\cdot\bV)\bu\cdot\bsig-\bV\cdot\bsig\right).
\end{align*}
Collecting terms gives Rodrigues' formula:
\begin{equation}
 \boxed{\bV'=\bV\cos\theta+(\bu\times\bV)\sin\theta
 +(\bu\cdot\bV)(1-\cos\theta)\bu.}
 \label{eq:rodrigues}
\end{equation}
This is a right-handed rotation through
$+\theta$ about $\bu$.
It preserves $\operatorname{sc}$, determinant, norm, and the real
paravector subspace.  For complex $\bV$, it rotates its real and imaginary
parts by the same real rotation.  The sign change in \eqref{eq:rotation-two-pi} affects both factors
in \eqref{eq:rotation-action}, so $Z'$ has period $2\pi$.

\subsection{Boost action and its sign convention}
\label{sec:algebra-boost-action}

Recall $L=\exp(\theta\bu\cdot\bsig/2)$ from
\eqref{eq:boost-exponential}; $\theta\in\R$, $\bu\in\R^3$, $\bu^2=1$,
$\beta=\tanh\theta$, and $\gamma=\cosh\theta$.
For a complex $Z=S+\bV\cdot\bsig$, choose the congruence
\begin{equation}
 Z'=L^{-1}ZL^{-1}.
 \label{eq:boost-action}
\end{equation}
The factors have the same sign because
$L^{\mathsf H}=L$ by \eqref{eq:boost-H}.
Write the complex vector as the sum of its parallel and perpendicular parts:
\begin{align*}
 \bV_{\parallel}&:=(\bu\cdot\bV)\bu,\\
 \bV_{\perp}&:=\bV-\bV_{\parallel}.
\end{align*}
The parallel term commutes with $\bu\cdot\bsig$; the perpendicular term
anticommutes with it.  Consequently,
\begin{align*}
 L^{-1}(S+\bV_{\parallel}\cdot\bsig)L^{-1}
 &=(S+\bV_{\parallel}\cdot\bsig)L^{-2},\\
 L^{-1}(\bV_{\perp}\cdot\bsig)L^{-1}
 &=\bV_{\perp}\cdot\bsig.
\end{align*}
Using $L^{-2}=\gamma(1-\beta\bu\cdot\bsig)$ from
\eqref{eq:boost-square}, multiplication gives
\begin{equation}
 \boxed{S'=\gamma(S-\beta\bu\cdot\bV).}
 \label{eq:boost-scalar-component}
\end{equation}
\begin{equation}
 \boxed{\bV'=\bV_{\perp}+\gamma(\bV_{\parallel}-\beta S\bu).}
 \label{eq:boost-vector-component}
\end{equation}
These preserve $S^2-\bV^2$ and reality; they generally
change $S$ and $\lVert Z\rVert^2$.

The inverse signs are a convention: \eqref{eq:boost-inverse}
gives $L^{-1}=\exp(-\theta\bu\cdot\bsig/2)=L(-\theta,\bu)$.
Replacing $\theta$ by $-\theta$ changes \eqref{eq:boost-action} to
$Z'=LZL$ and changes $\beta$ to $-\beta$ in the component equations.
The two conventions describe opposite boosts.  We keep the positive
exponential definition of $L$ throughout; in the spacetime application,
\eqref{eq:boost-action} gives components in the frame moving with velocity
$+\beta\bu$.

A boost similarity is a different action.  For $Z'=L^{-1}ZL$, the same
commutation rules instead give $S'=S$ and
\begin{equation}
 \bV'=\bV_{\parallel}
 +\gamma\bigl(\bV_{\perp}-\ii\beta\bu\times\bV\bigr).
 \label{eq:boost-similarity-vector}
\end{equation}
This similarity preserves scalar part and
determinant, but a real vector generally
acquires an imaginary component.  It will act on the complex electromagnetic
field, whose real and imaginary vectors represent two physical fields.

\section{Differentiated fields}
\label{sec:differentiated-fields}

\subsection{Paravector differential operators}

Let $t\in\R$ and $\br=\{r_1,r_2,r_3\}\in\R^3$ be real coordinates.
Define their scalar and vector derivatives by
\[
 \partial_t:=\frac{\partial}{\partial t},\qquad
 \pdr:=\left\{\frac{\partial}{\partial r_1},
              \frac{\partial}{\partial r_2},
              \frac{\partial}{\partial r_3}\right\}.
\]
Define the partial paravector and its sigma conjugate by
\begin{equation}
 \partial:=\partial_t+\bsig\cdot\pdr
 =\partial_t+\sum_{n=1}^{3}\sigma_n\frac{\partial}{\partial r_n},
 \label{eq:paravector-derivative}
\end{equation}
\begin{equation}
 \partial^*=\adj{\partial}=\partial_t-\bsig\cdot\pdr.
 \label{eq:conjugate-paravector-derivative}
\end{equation}
The sigma generators are constant.  Derivatives act from the left on
everything to their right unless parentheses specify their scope.
For fields with commuting second derivatives, anticommutation cancels
all mixed spatial terms:
\begin{align*}
 (\bsig\cdot\pdr)^2
 &=\sum_{n=1}^{3}\frac{\partial^2}{\partial r_n^2}
 +\sum_{n<m}(\sigma_n\sigma_m+\sigma_m\sigma_n)
           \frac{\partial^2}{\partial r_n\partial r_m}\\
 &=\pdr\cdot\pdr.
\end{align*}
The mixed time--space terms likewise cancel.  Both orders therefore give
the scalar d'Alembertian
\begin{equation}
 \partial\partial^*=\partial^*\partial
 =\partial_t^2-\pdr\cdot\pdr=:\Box.
 \label{eq:dalembertian}
\end{equation}
These are operator identities: $\partial\partial^*Z$ means
$\partial(\partial^*Z)$.

\subsection{Components and the product rule}

Let the complex paravector field be
\[
 Z(t,\br)=(a+b\ii)+(\bA+\ii\bB)\cdot\bsig,
\]
with real scalar functions $a,b:\R^4\to\R$ and real vector functions
$\bA,\bB:\R^4\to\R^3$.  Ordered vector multiplication gives its four
real components:
\begin{equation}
 \begin{aligned}
 \partial Z={}&(\partial_ta+\pdr\cdot\bA)
 +\ii(\partial_tb+\pdr\cdot\bB)\\
 &+(\pdr a+\partial_t\bA-\pdr\times\bB)\cdot\bsig\\
 &+\ii(\pdr b+\pdr\times\bA+\partial_t\bB)\cdot\bsig.
 \end{aligned}
 \label{eq:differentiated-components}
\end{equation}
Multiplication by $\ii$ exchanges real and imaginary parts, with the
former imaginary parts acquiring a minus sign.

For a second complex paravector field $Z'(t,\br)$, the prime labels the
field rather than a derivative.  Apply the scalar product rule before
reordering any factors:
\begin{align*}
 \partial(ZZ')
 ={}&(\partial_tZ)Z'+Z(\partial_tZ')\\
 &+\sum_{n=1}^{3}\sigma_n
   \left(\frac{\partial Z}{\partial r_n}Z'
          +Z\frac{\partial Z'}{\partial r_n}\right)\\
 ={}&(\partial Z)Z'+Z(\partial Z')
    +\sum_{n=1}^{3}[\sigma_n,Z]\frac{\partial Z'}{\partial r_n}.
\end{align*}
Define the last sum as
$[\bsig,Z]\cdot\pdr Z':=\sum_n[\sigma_n,Z]\,\partial Z'/\partial r_n$.
Changing the spatial sign gives the paired rules
\begin{align}
 \partial(ZZ')&=(\partial Z)Z'+Z(\partial Z')
                +[\bsig,Z]\cdot\pdr Z',
 \label{eq:paravector-product-rule}\\
 \partial^*(ZZ')&=(\partial^*Z)Z'+Z(\partial^*Z')
                -[\bsig,Z]\cdot\pdr Z'.
 \label{eq:conjugate-product-rule}
\end{align}
The contraction retains the displayed multiplied order.

For a noncommuting check, take $Z=\sigma_2$ and $Z'=r_1$.  Since
$[\sigma_1,\sigma_2]=2\ii\sigma_3$, the two right-hand sides give
\begin{align*}
 \partial(\sigma_2r_1)
 &=\sigma_2\sigma_1+[\sigma_1,\sigma_2]
  =-\ii\sigma_3+2\ii\sigma_3=\ii\sigma_3,\\
 \partial^*(\sigma_2r_1)
 &=-\sigma_2\sigma_1-[\sigma_1,\sigma_2]
  =\ii\sigma_3-2\ii\sigma_3=-\ii\sigma_3.
\end{align*}
These agree with direct differentiation, $\pm\sigma_1\sigma_2$.
The ordinary Leibniz rule holds when the left factor commutes with every
$\sigma_n$, in particular for a complex scalar.  Otherwise the commutator
term must be retained.

\section{Space--time kinematics}
\label{sec:kinematics}

Let $X=t+\br\cdot\bsig$ be a real event, with time $t\in\R$ and
position $\br\in\R^3$.  A massive particle follows a future-directed
time-like path $\br(t)$; $s$ denotes its proper time.

\subsection{Intervals, velocity, and momentum}

Recall the real proper velocity $U=\dd X/\dd s$, a
\hyperref[eq:unit-paravector]{unit paravector} with $UU^{-1}=1$.
The differential event and its invariant interval are
\[
 \dd X=\dd t+\dd\br\cdot\bsig.
\]
\begin{equation}
 \dd s^2=\det(\dd X)=\dd X\,\dd X^*
 =\dd t^2-\dd\br^2.
 \label{eq:interval}
\end{equation}
By \eqref{eq:general-lorentz-congruence}, a constant complex
\hyperref[eq:unit-paravector]{unit paravector} $T$ with
\hyperref[eq:unit-determinant-paravector]{$\det T=1$} gives
\[
 \dd X'=T^{-1}\dd X(T^{-1})^{\mathsf H},
 \qquad \det(\dd X')=\dd s^2.
\]
Positive, zero, and negative values of \eqref{eq:interval} mean time-like, null, and space-like
separation.  Massive particles follow time-like paths; light follows null
paths; causal signals cannot connect space-like separated events.
For two real intervals,
\begin{align*}
 [\dd X,\dd X']&=2\ii(\dd\br\times\dd\br')\cdot\bsig,\\
 (\dd X)^{-1}&=\frac{\dd X^*}{\dd s^2}\qquad(\dd s^2\ne0).
\end{align*}
A null interval has no inverse.

Along the massive path, define the real coordinate velocity
$\bv=\dd\br/\dd t$ and choose $\dd s>0$.  Then
$\dd s^2=(1-\bv^2)\dd t^2$ gives
\begin{align}
 \gamma&:=\frac1{\sqrt{1-\bv^2}}=\frac{\dd t}{\dd s},
 \label{eq:lorentz-factor}\\
 \frac{\dd}{\dd s}&=\gamma\frac{\dd}{\dd t}.
 \label{eq:proper-time-derivative}
\end{align}
Applying \eqref{eq:proper-time-derivative} with \eqref{eq:lorentz-factor}
gives the real proper velocity
\begin{equation}
 U:=\frac{\dd X}{\dd s}=\gamma(1+\bv\cdot\bsig).
 \label{eq:proper-velocity}
\end{equation}
The normalization \eqref{eq:proper-velocity-star-normalization} gives
$U^*U=1$ and $U^{-1}=U^*=\gamma(1-\bv\cdot\bsig)$.
Since $U^{\mathsf H}=U$, this is
\hyperref[eq:unit-determinant-paravector]{determinant one}, not
Hermitian unitarity except at rest.
For a \hyperref[def:real-unit-vector]{real unit direction} $\bu$, write $\bv=\beta\bu$ and
$\theta=\tanh^{-1}\beta$.  Equation~\eqref{eq:boost-square} gives $U=L^2$ for
$L=\exp(\theta\bu\cdot\bsig/2)$.

For constant real mass $m>0$, define the real four-momentum
\begin{equation}
 P:=mU=E+\bp\cdot\bsig.
 \label{eq:four-momentum}
\end{equation}
Its real coefficients are energy and momentum:
\begin{align*}
 E&=\scpart P=m\gamma,\\
 \bp&=\vecpart P=m\gamma\bv,\\
 P^{-1}&=\frac{E-\bp\cdot\bsig}{m^2}.
\end{align*}
Since $\det U=1$, the determinant gives the mass shell
\begin{equation}
 E^2-\bp^2=\det P=m^2.
 \label{eq:mass-shell}
\end{equation}
For a massive particle of finite energy,
\[
 \bv^2=\frac{\bp^2}{E^2}=1-\frac{m^2}{E^2}<1.
\]
Approaching light speed requires unbounded energy.
A massless path has no proper-time velocity.  Define its momentum directly
as a future null paravector $P=E+\bp\cdot\bsig$, with $E>0$ and
$E^2=\bp^2$; it has no inverse.

\subsection{Proper acceleration}
\label{sec:proper-acceleration}

Recall the real proper velocity $U=\dd X/\dd s$, a
\hyperref[eq:unit-paravector]{unit paravector} with $UU^{-1}=1$,
and the real coordinate velocity $\bv(t)$ and Lorentz factor $\gamma$.
Define the real coordinate acceleration
$\ba=\dd\bv/\dd t$.  Differentiating gives
\begin{align*}
 \frac{\dd\gamma}{\dd t}
 &=-\tfrac12(1-\bv^2)^{-3/2}(-2\bv\cdot\ba)
 =\gamma^3\bv\cdot\ba,\\
 \frac{\dd^2}{\dd s^2}
 &=\gamma\frac{\dd}{\dd t}\left(\gamma\frac{\dd}{\dd t}\right)
 =\gamma^2\frac{\dd^2}{\dd t^2}
  +\gamma^4(\bv\cdot\ba)\frac{\dd}{\dd t}.
\end{align*}
Applying this to $X$ yields the real proper-acceleration paravector:
\begin{equation}
 \frac{\dd U}{\dd s}
 =\gamma^4(\bv\cdot\ba)
 +\left(\gamma^2\ba+\gamma^4(\bv\cdot\ba)\bv\right)\cdot\bsig.
 \label{eq:proper-acceleration-components}
\end{equation}
Differentiating $U^*U=1$ gives
\[
 \left(\frac{\dd U}{\dd s}\right)^*U
 +U^*\frac{\dd U}{\dd s}=0.
\]
The two terms have equal real scalar parts, so
\begin{equation}
 \scpart{U^*\dd U/\dd s}=0.
 \label{eq:proper-acceleration-orthogonal}
\end{equation}
Thus four-velocity and four-acceleration are Minkowski-orthogonal.
Multiplication on the right by $U^*$ also gives
\begin{align*}
 \adj{\frac{\dd U}{\dd s}}
 &=-U^*\frac{\dd U}{\dd s}U^*.
\end{align*}
\begin{equation}
 U^*\frac{\dd U}{\dd s}
 =\gamma^3(\ba-\ii\bv\times\ba)\cdot\bsig.
 \label{eq:velocity-acceleration-product}
\end{equation}
The real vector records acceleration and the imaginary vector records
turning of the velocity.  Since $\ba\cdot(\bv\times\ba)=0$,
\begin{equation}
 \begin{aligned}
 \det\left(\frac{\dd U}{\dd s}\right)
 &=-\left(U^*\frac{\dd U}{\dd s}\right)^2\\
 &=-\gamma^6\left(\ba^2-(\bv\times\ba)^2\right)\\
 &=-\gamma^4\left(\ba^2+\gamma^2(\bv\cdot\ba)^2\right).
 \end{aligned}
 \label{eq:proper-acceleration-norm}
\end{equation}
The last equality uses
$(\bv\times\ba)^2=\bv^2\ba^2-(\bv\cdot\ba)^2$ and
$1-\bv^2=\gamma^{-2}$.  Nonzero proper acceleration is therefore
space-like; its invariant squared magnitude is the negative determinant.
For $\bv\ne0$, define
$\ba_{\parallel}=(\bv\cdot\ba)\bv/\bv^2$ and
$\ba_{\perp}=\ba-\ba_{\parallel}$.  Equivalently,
\begin{equation}
 -\det\left(\frac{\dd U}{\dd s}\right)
 =\gamma^4\ba_{\perp}^2+\gamma^6\ba_{\parallel}^2.
 \label{eq:proper-acceleration-split}
\end{equation}
At rest this reduces to $\ba^2$.

\paragraph{Energy and inertia.}
Recall $E=m\gamma$ and $\bp=m\gamma\bv$ for invariant real mass $m>0$.
Differentiating, using the derivative of $\gamma$ above, gives
\[
 \frac{\dd\bp}{\dd t}
 =m\gamma\ba+m\gamma^3(\bv\cdot\ba)\bv.
\]
For $\bv\ne0$, resolving along and perpendicular to $\bv$ yields
\begin{equation}
 \boxed{\frac{\dd\bp}{\dd t}
 =m\gamma\ba_{\perp}+m\gamma^3\ba_{\parallel}.}
 \label{eq:relativistic-inertia}
\end{equation}
Taking its dot product with $\bv$ gives the work rate
\begin{equation}
 \frac{\dd E}{\dd t}=m\gamma^3\bv\cdot\ba
 =\bv\cdot\frac{\dd\bp}{\dd t}.
 \label{eq:relativistic-work}
\end{equation}
At rest the force is $m\ba$.  Restoring $c$, $E=\gamma mc^2$;
$E/c^2=\gamma m$ is sometimes called \emph{relativistic mass}.
We retain invariant $m$: the longitudinal and transverse inertia in
\eqref{eq:relativistic-inertia} differ
\cite[pp.~136--137, 140]{tong-relativity}.

\subsection{World lines, clocks, and light cones}
\label{sec:worldlines}

\subsubsection{Integrating proper velocity}
Recall the real path $X(s)=t(s)+\br(s)\cdot\bsig$ and its tangent
$U=\dd X/\dd s$, a \hyperref[eq:unit-paravector]{unit paravector}
with $UU^{-1}=1$.  For real proper times $s_0<s_1$, integration gives
\begin{equation}
 \boxed{X(s_1)-X(s_0)=\int_{s_0}^{s_1}U(s)\dd s.}
 \label{eq:worldline-integral}
\end{equation}
Its real components are
\begin{align*}
 t(s_1)-t(s_0)&=\int_{s_0}^{s_1}\gamma(s)\dd s,\\
 \br(s_1)-\br(s_0)&=\int_{s_0}^{s_1}\gamma(s)\bv(s)\dd s.
\end{align*}
Thus between coordinate times $t_0$ and $t$ the clock records
\begin{equation}
 \boxed{s(t)-s(t_0)=\int_{t_0}^{t}\frac{\dd t'}{\gamma(t')}
 =\int_{t_0}^{t}\sqrt{1-\bv(t')^2}\dd t'.}
 \label{eq:elapsed-proper-time}
\end{equation}
This is the standard world-line proper-time integral
\cite[pp.~133--134, eq.~(7.19)]{tong-relativity}.
For constant $U$, \eqref{eq:worldline-integral} is the straight line
$X(s)=X(s_0)+(s-s_0)U$.  For varying $U$, integrate the tangent;
an endpoint boost alone does not specify the path.

\subsubsection{Clock comparison}
Let $X_0=X(s_0)$ and $X_1=X(s_1)$ be fixed, time-like separated events.
Write $X_k=t_k+\br_k\cdot\bsig$ for $k=0,1$.
In their separation's rest frame, $\br_1=\br_0$ and
$t_1-t_0=\sqrt{\det(X_1-X_0)}$.  Since $\sqrt{1-\bv^2}\le1$,
\eqref{eq:elapsed-proper-time} gives
\begin{equation}
 s_1-s_0\le\sqrt{\det(X_1-X_0)}.
 \label{eq:maximum-proper-time}
\end{equation}
Equality requires the straight inertial path.  The determinant of a finite
displacement therefore does not generally equal the squared elapsed
proper time of an accelerated clock.

For a return trip lasting real coordinate time $T>0$, take velocities
$+\beta\bu$ and $-\beta\bu$ for $T/2$ each, with
\hyperref[def:real-unit-vector]{real unit vector} $\bu$ and
$0<\beta<1$.  The ideal instantaneous reversal is the limit of a short
smooth turn.  The traveling clock records
\[
 s(T)-s(0)=T\sqrt{1-\beta^2},
\]
while the clock at rest records $T$.  This is the twin-clock comparison
\cite[pp.~120--122]{tong-relativity}.

\subsubsection{Constant proper acceleration}
Recall the real proper velocity $U=\dd X/\dd s$, a
\hyperref[eq:unit-paravector]{unit paravector} with $UU^{-1}=1$.
Let $a>0$ be a constant real proper-acceleration magnitude and $\bu$ a fixed
\hyperref[def:real-unit-vector]{real unit axis}.  Starting at $X(0)=0$ and $U(0)=1$, take rapidity
$\theta(s)=as$.  Then
\[
 U(s)=L(as,\bu)^2=\cosh(as)+\sinh(as)\bu\cdot\bsig.
\]
Differentiation verifies the stated acceleration:
\[
 -\det\left(\frac{\dd U}{\dd s}\right)
 =a^2\bigl(\cosh^2(as)-\sinh^2(as)\bigr)=a^2.
\]
Integrating \eqref{eq:worldline-integral} yields
\begin{equation}
 \boxed{X(s)=\frac{\sinh(as)}{a}
       +\frac{\cosh(as)-1}{a}\bu\cdot\bsig.}
 \label{eq:hyperbolic-worldline}
\end{equation}
Its components give
\begin{align*}
 \bv(s)&=\frac{\dd\br/\dd s}{\dd t/\dd s}=\tanh(as)\bu,\\
 \bigl(1+a\bu\cdot\br\bigr)^2-(at)^2&=1.
\end{align*}
Thus \eqref{eq:hyperbolic-worldline} is the usual hyperbolic motion, with
speed below light speed at every finite $s$
\cite[pp.~142--143]{tong-relativity}.

\subsubsection{Light cones}
For a real event separation $\Delta X=\Delta t+\Delta\br\cdot\bsig$,
the eigenvalues are $\Delta t\pm\lVert\Delta\br\rVert$.
Both are nonnegative precisely in the future causal cone:
\begin{equation}
 \Delta t\ge\lVert\Delta\br\rVert.
 \label{eq:future-light-cone}
\end{equation}
The boundary has $\det(\Delta X)=0$; its nonzero future rays are
\[
 \Delta X=\Delta t(1+\bu\cdot\bsig),
 \qquad \Delta t>0,\quad\bu^2=1.
\]
The Hermitian congruence~\eqref{eq:general-lorentz-congruence} preserves
nonnegativity of the eigenvalues and the determinant, hence the cone and its
boundary.  For a future time-like or null path,
\[
 \lVert\Delta\br\rVert
 =\left\lVert\int\bv\dd t\right\rVert
 \le\int\lVert\bv\rVert\dd t\le\Delta t.
\]
Its endpoints therefore satisfy \eqref{eq:future-light-cone}; a space-like
separation cannot carry a causal signal
\cite[pp.~117--118]{tong-relativity}.

\subsection{Free-particle action}

Recall real $m>0$, $\bv=\dd\br/\dd t$, and $\dd s=\dd t/\gamma$.
Here $L(t,\br,\bv)$ is a real scalar Lagrangian, not the boost factor.
The free action is the physical choice
\begin{equation}
 \int L\dd t=-m\int\dd s.
 \label{eq:free-particle-action}
\end{equation}
Using $\dd s=\dd t/\gamma$ gives the Lagrangian
\begin{equation}
 L=-m\sqrt{1-\bv^2}=-m/\gamma.
 \label{eq:free-particle-lagrangian}
\end{equation}
For endpoint-fixed real path variations $\delta\br(t)$, varying the
Lagrangian and integrating by parts gives
\[
 \delta\int L\dd t
 =\int\left(\pdr L-\frac{\dd}{\dd t}\partial_{\bv}L\right)
 \cdot\delta\br\dd t.
\]
Stationarity for every variation yields the Euler--Lagrange equation.
With real Hamiltonian $H:=\bv\cdot\partial_{\bv}L-L$, the chain rule gives
\begin{align}
 \frac{\dd}{\dd t}\partial_{\bv}L&=\pdr L, \label{eq:euler-lagrange}\\
 \frac{\dd H}{\dd t}
 &=\bv\cdot\left(\frac{\dd}{\dd t}\partial_{\bv}L-\pdr L\right)
   -\partial_tL=-\partial_tL. \label{eq:hamiltonian-energy-balance}
\end{align}
By \eqref{eq:euler-lagrange} and \eqref{eq:hamiltonian-energy-balance},
spatial and temporal translation invariance conserve canonical
momentum and energy.  For the free action,
\begin{align*}
 \partial_{\bv}L&=m\gamma\bv=\bp,\\
 H&=m\gamma\bv^2+m/\gamma=m\gamma=E.
\end{align*}
Writing $v=\lVert\bv\rVert\ll1$, the binomial series gives
\begin{align*}
 \sqrt{1-v^2}&=1-\tfrac12v^2+O(v^4),\\
 \gamma&=1+\tfrac12v^2+O(v^4),\\
 L&=-m+\tfrac12mv^2+O(mv^4),\\
 H&=m+\tfrac12mv^2+O(mv^4).
\end{align*}
The exact kinetic energy above rest is
\begin{equation}
 K:=H-m=m(\gamma-1).
 \label{eq:kinetic-energy}
\end{equation}
This and both expansions recover $mv^2/2$
at low speed after removal of the constant rest term.

\subsection{Flat spacetime in generalized coordinates}
\label{sec:generalized}

Here $(t,\br)$ denote generalized real chart coordinates.  Smooth functions
$f:\R^4\to\R$ and $\bg:\R^4\to\R^3$ specify inertial coordinates
through the real paravector
\[
 X=f(t,\br)+\bg(t,\br)\cdot\bsig.
\]
Along a future timelike path with proper time $s$, recall the real
$U=\dd X/\dd s$, a \hyperref[eq:unit-paravector]{unit paravector}
with $UU^{-1}=1$.
Require the Jacobian of $(f,\bg)$ to be nonsingular.  Differentiating gives
\begin{align*}
 \dd X&=(\partial_tf+\partial_t\bg\cdot\bsig)\dd t
 +\sum_{k=1}^3(\partial_{r_k}f+\partial_{r_k}\bg\cdot\bsig)\dd r_k,\\
 \dd s^2&=\dd X\,\dd X^*=\dd f^2-\dd\bg^2.
\end{align*}
With $x_0=t$, $x_k=r_k$ and $\partial_k=\partial/\partial x_k$, the
interval is $\dd s^2=\sum_{j,k=0}^3 g_{jk}\dd x_j\dd x_k$, with real coefficients
\begin{equation}
 g_{jk}=\partial_jf\,\partial_kf-\partial_j\bg\cdot\partial_k\bg,
 \qquad j,k=0,1,2,3.
 \label{eq:coordinate-metric}
\end{equation}
This is the flat Minkowski metric in changed coordinates; varying
coefficients do not imply curvature.

Along this future timelike curve, $\dd f>0$.
Define its real inertial three-velocity $\bv=\dd\bg/\dd f$ and
$\gamma=(1-\bv^2)^{-1/2}$, with $\bv^2<1$.  Then
\begin{align*}
 \dd s&=\sqrt{1-\bv^2}\,\dd f,\\
 U=\frac{\dd X}{\dd s}&=\gamma(1+\bv\cdot\bsig).
\end{align*}
Apply \eqref{eq:proper-acceleration-components} with inertial time $f$
and coordinate acceleration $\dd\bv/\dd f$.
For inertial motion, the scalar part gives $\bv\cdot\dd\bv/\dd f=0$,
and its vector part then gives $\dd\bv/\dd f=0$.  Therefore
\begin{equation}
 \bg=\bC_1f+\bC_2.
 \label{eq:inertial-coordinate-path}
\end{equation}
The real constants $\bC_1,\bC_2\in\R^3$ are velocity and initial position.
If primes denote total derivatives along this curve with respect to $t$,
then, wherever $f'\ne0$,
\begin{align*}
 \frac{\dd\bv}{\dd t}&=\frac{\bg''f'-\bg'f''}{(f')^2},\\
 [U,\dd U/\dd s]&=2\ii\gamma^2\left(\bv\times\frac{\dd\bv}{\dd s}\right)\cdot\bsig.
\end{align*}

\subsection{Varying frames and local bases}
\label{sec:varying-transformations}

Use a fixed coordinate chart and let the local frame vary along a curve.
Let $T(s)$ be a smooth complex \hyperref[eq:unit-paravector]{unit paravector}
of \hyperref[eq:unit-determinant-paravector]{determinant one} along a real parameter $s$.  When $s$ is proper time,
recall the real tangent $U=\dd X/\dd s$, a unit paravector with $UU^{-1}=1$.
Define the complex paravector
$\Omega(s):=T^{-1}\dd T/\dd s$.  The inverse rule gives
\[
 \frac{\dd T^{-1}}{\dd s}=-\Omega T^{-1}.
\]
For a complex paravector $Z(s)$ transformed by similarity,
$Z'=T^{-1}ZT$, product differentiation yields
\begin{equation}
 \frac{\dd Z'}{\dd s}
 =T^{-1}\frac{\dd Z}{\dd s}T+[Z',\Omega].
 \label{eq:varying-similarity}
\end{equation}
Differentiating again and replacing
$T^{-1}(\dd Z/\dd s)T$ by its first-derivative expression gives
\begin{equation}
 \begin{aligned}
 \frac{\dd^2 Z'}{\dd s^2}={}&T^{-1}\frac{\dd^2Z}{\dd s^2}T
 +2\left[\frac{\dd Z'}{\dd s},\Omega\right]\\
 &+\left[Z',\frac{\dd\Omega}{\dd s}\right]
 +[\Omega,[Z',\Omega]].
 \end{aligned}
 \label{eq:varying-similarity-second}
\end{equation}
The last term follows from
$-[ [Z',\Omega],\Omega]=[\Omega,[Z',\Omega]]$.
For a real paravector transformed instead by congruence,
$Z'=T^{-1}Z(T^{-1})^{\mathsf H}$, the same rule gives
\begin{equation}
 \frac{\dd Z'}{\dd s}
 =T^{-1}\frac{\dd Z}{\dd s}(T^{-1})^{\mathsf H}
 -\Omega Z'-Z'\Omega^{\mathsf H}.
 \label{eq:varying-congruence}
\end{equation}
Equations~\eqref{eq:varying-similarity} and \eqref{eq:varying-congruence}
include the frame terms omitted by a constant-factor transformation.
For a varying rotation $T=R=\exp(\ii\theta(s)\bu(s)\cdot\bsig/2)$,
with real $\theta(s)$ and \hyperref[def:real-unit-vector]{real unit vector} $\bu(s)$, unitarity gives
$\Omega^{\mathsf H}=-\Omega$, so both formulas agree.

When $s$ is proper time, applying the appropriate formula to the pointwise
transformed tangent $U$ gives its acceleration law.  This differs from
first transforming the event $X$ and then differentiating: a varying event
map already contributes frame terms to its first derivative.
The similarity formula also gives the local algebra-basis law
$\sigma_k'=T^{-1}\sigma_kT$, with
\begin{equation}
 \frac{\dd\sigma_k'}{\dd s}=[\sigma_k',\Omega],\qquad k=1,2,3.
 \label{eq:moving-algebra-basis}
\end{equation}
This evolution preserves all generator products.  A general complex local basis
is nevertheless distinct from a real spacetime frame, whose vectors require
the congruence law.

\section{Electromagnetism}
\label{sec:electromagnetism}

All coefficient fields in this section are real functions of real inertial
coordinates $(t,\br)$.  Define
\begin{align}
 F&=(\bE+\ii\bBf)\cdot\bsig
 &&\text{electric and magnetic fields}, \label{eq:em-field-definition}\\
 \Phi&=V+\bA\cdot\bsig
 &&\text{scalar and vector potentials}, \notag\\
 D&=\rho+\bJ\cdot\bsig
 &&\text{charge and current densities}. \notag
\end{align}
Thus $\Phi,D$ are real paravectors, while $F$ is a complex pure vector.
Recall $\partial=\partial_t+\bsig\cdot\pdr$ and
$\partial^*\partial=\Box=\partial_t^2-\pdr^2$.

The dot/cross product~\eqref{eq:vector-product} combines metric and
orientation.  With a spatial derivative it gives divergence and curl,
while the six real components of $F$ represent an antisymmetric spacetime
field tensor.  This algebra of physical space is isomorphic to the even
subalgebra of spacetime algebra; its real paravectors represent spacetime
vectors relative to an observer~\cite[Secs.~IV and V]{baylis2004}.
The product explains the compact form; Maxwell dynamics remains physical
input, derived from an action in Section~\ref{sec:maxwell-action}.

\subsection{Maxwell's equations and wave propagation}

Maxwell's four equations are the single field law
\begin{equation}
 \boxed{\partial F=D^*=\rho-\bJ\cdot\bsig.}
 \label{eq:maxwell}
\end{equation}
The product rule expands its left side into
\begin{align*}
 \partial F={}&\underset{\text{real scalar}}{\pdr\cdot\bE}
 +\underset{\text{imaginary scalar}}{\ii\pdr\cdot\bBf}\\
 &+\underset{\text{real vector}}{(\partial_t\bE-\pdr\times\bBf)\cdot\bsig}\\
 &+\underset{\text{imaginary vector}}{\ii(\partial_t\bBf+\pdr\times\bE)\cdot\bsig}.
\end{align*}
Equating the four parts gives
\begin{align}
 \pdr\cdot\bE&=\rho, \label{eq:gauss-electric}\\
 \pdr\cdot\bBf&=0, \label{eq:gauss-magnetic}\\
 \pdr\times\bBf-\partial_t\bE&=\bJ, \label{eq:ampere-maxwell}\\
 \pdr\times\bE+\partial_t\bBf&=0. \label{eq:faraday}
\end{align}
Equations~\eqref{eq:gauss-electric}--\eqref{eq:faraday} are Gauss'
electric and magnetic laws, Amp\`ere--Maxwell's law, and Faraday's law~\cite[p.~6]{tong-em}.
Applying $\partial^*$ to \eqref{eq:maxwell} gives
\begin{equation}
 \begin{aligned}
 \Box F&=\partial^*D^*=(\partial D)^*\\
 &=(\partial_t\rho+\pdr\cdot\bJ)
 +(-\partial_t\bJ-\pdr\rho+\ii\pdr\times\bJ)\cdot\bsig.
 \end{aligned}
 \label{eq:maxwell-wave}
\end{equation}
The scalar part vanishes because $F$ is scalar-free.  Thus
\begin{align}
 \partial_t\rho+\pdr\cdot\bJ&=0, \label{eq:charge-continuity}\\
 \Box\bE&=-\partial_t\bJ-\pdr\rho, \label{eq:electric-wave}\\
 \Box\bBf&=\pdr\times\bJ. \label{eq:magnetic-wave}
\end{align}
Equation~\eqref{eq:charge-continuity} is charge conservation;
\eqref{eq:electric-wave} and \eqref{eq:magnetic-wave} are the sourced field waves.
In vacuum, $D=0$ and both fields obey the free wave equation.

\subsection{Potentials and gauge freedom}

Recall the real potential $\Phi=V+\bA\cdot\bsig$.  Its derivative is
\begin{align*}
 \partial\Phi={}&(\partial_tV+\pdr\cdot\bA)
 +(\partial_t\bA+\pdr V)\cdot\bsig
 +\ii(\pdr\times\bA)\cdot\bsig.
\end{align*}
Define the real scalar $S:=\partial_tV+\pdr\cdot\bA$.  Identifying the
remaining components with $F^*$ gives
\begin{align}
 \partial\Phi&=S+F^*, \label{eq:potential-field}\\
 \bE&=-\partial_t\bA-\pdr V, \label{eq:electric-potential}\\
 \bBf&=\pdr\times\bA. \label{eq:magnetic-potential}
\end{align}
Equation~\eqref{eq:potential-field} separates the gauge scalar from the field;
\eqref{eq:electric-potential} and \eqref{eq:magnetic-potential} satisfy
\eqref{eq:gauss-magnetic} and \eqref{eq:faraday} identically.
Since $F=\partial^*\Phi^*-S$,
\begin{align*}
 \partial F&=\Box\Phi^*-\partial S\\
 &=(\Box V-\partial_tS)-(\Box\bA+\pdr S)\cdot\bsig.
\end{align*}
Comparison with \eqref{eq:maxwell} yields
\begin{align}
 \rho&=\Box V-\partial_tS, \label{eq:potential-charge}\\
 \bJ&=\Box\bA+\pdr S. \label{eq:potential-current}
\end{align}
These source equations hold before gauge fixing.
For a real scalar gauge function $\lambda(t,\br)$, let
\begin{equation}
 \Phi'=\Phi+\partial^*\lambda.
 \label{eq:gauge-transformation}
\end{equation}
Differentiating gives
$\partial\Phi'=\partial\Phi+\Box\lambda$.  Only the scalar changes: $S'=S+\Box\lambda$, while $F'=F$.
In the Lorenz gauge $S=0$, the source equations reduce to
$\Box V=\rho$ and $\Box\bA=\bJ$.

\subsection{The electromagnetic field action}
\label{sec:maxwell-action}

For real scalar potential $V(t,\br)$ and vector potential
$\bA(t,\br)$, write $\Phi=V+\bA\cdot\bsig$.  With prescribed $D$,
Maxwell dynamics follows from the real Lagrangian density
\begin{equation}
 L=\tfrac12\repart{F^2}-\scpart{\Phi^*D}
  =\tfrac12(\bE^2-\bBf^2)-\rho V+\bJ\cdot\bA.
 \label{eq:maxwell-lagrangian}
\end{equation}
In varying the action, boundary-vanishing variations satisfy
\begin{align*}
 \delta\bE&=-\pdr\delta V-\partial_t\delta\bA,\\
 \delta\bBf&=\pdr\times\delta\bA.
\end{align*}
Integration by parts gives
\begin{equation}
 \delta\int L\dd t\dd^3\br
 =\int\left((\pdr\cdot\bE-\rho)\delta V
 +(\partial_t\bE-\pdr\times\bBf+\bJ)\cdot\delta\bA\right)\dd t\dd^3\br.
 \label{eq:maxwell-action-variation}
\end{equation}
Independent variations give
\eqref{eq:gauss-electric} and \eqref{eq:ampere-maxwell}; the potential
identities supply \eqref{eq:gauss-magnetic} and \eqref{eq:faraday}.
Together they recover \eqref{eq:maxwell}.  The symbol $\delta$ denotes variation.

\subsection{Field invariants}

Since $\scpart F=0$, its adjugate is $-F$.  Squaring gives
\begin{equation}
 F^2=-\det F=\bE^2-\bBf^2+2\ii\bE\cdot\bBf.
 \label{eq:field-invariants}
\end{equation}
For a complex \hyperref[eq:unit-paravector]{unit paravector} $T$, the field similarity
$F'=T^{-1}FT$ preserves this scalar by \eqref{eq:det-similarity}.
Physical rotations and boosts use
$R=\exp(\ii\theta\bu\cdot\bsig/2)$ and
$L=\exp(\theta\bu\cdot\bsig/2)$, respectively, with
\hyperref[def:real-unit-vector]{real unit vector} $\bu$
and real angle or rapidity $\theta$; their field laws are derived in
Section~\ref{sec:transformations}.
The real and imaginary parts of \eqref{eq:field-invariants} are
therefore observer-independent.
For a nonzero field, they give the classification
\begin{center}
\begin{tabular}{ll}
\toprule
Condition & Consequence\\
\midrule
$\bE\cdot\bBf=0$, $\bE^2>\bBf^2$ & Some frame has $\bBf'=0$.\\
$\bE\cdot\bBf=0$, $\bE^2<\bBf^2$ & Some frame has $\bE'=0$.\\
$\bE\cdot\bBf\ne0$ & Neither field can be eliminated.\\
$\bE\cdot\bBf=0$, $\bE^2=\bBf^2$ & Null field: equal, perpendicular fields.\\
\bottomrule
\end{tabular}
\end{center}

\subsection{Lorentz force from the action}

Recall the real proper velocity $U=\dd X/\dd s=\gamma(1+\bv\cdot\bsig)$,
a \hyperref[eq:unit-paravector]{unit paravector} with $UU^{-1}=1$.
Let $q\in\R$ be charge and $m>0$ constant mass.  The real momentum is
$P=mU=E+\bp\cdot\bsig$, with $E=m\gamma$ and $\bp=m\gamma\bv$.
The real particle Lagrangian $L$ is fixed by the action
\begin{equation}
 \int L\dd t=-m\int\dd s-q\int\scpart{\Phi^*U}\dd s.
 \label{eq:charged-particle-action}
\end{equation}
Since $\scpart{\Phi^*U}=\gamma(V-\bv\cdot\bA)$, the Lagrangian is
\begin{equation}
 L=-m/\gamma-qV+q\bv\cdot\bA.
 \label{eq:charged-particle-lagrangian}
\end{equation}
Its Euler--Lagrange terms are
\begin{align*}
 \partial_{\bv}L&=\bp+q\bA,\\
 \frac{\dd}{\dd t}\partial_{\bv}L
 &=\frac{\dd\bp}{\dd t}+q\partial_t\bA+q(\bv\cdot\pdr)\bA,\\
 \pdr L&=-q\pdr V+q\pdr(\bv\cdot\bA).
\end{align*}
Here spatial differentiation holds the independent velocity argument fixed.
Using
$\pdr(\bv\cdot\bA)-(\bv\cdot\pdr)\bA
=\bv\times(\pdr\times\bA)$ gives
\begin{equation}
 \frac{\dd\bp}{\dd t}=q(\bE+\bv\times\bBf).
 \label{eq:lorentz-momentum}
\end{equation}
Here \eqref{eq:electric-potential} and \eqref{eq:magnetic-potential}
turn the potential derivatives into the fields.
The Hamiltonian and its balance law are
\begin{equation}
 H=\bv\cdot(\bp+q\bA)-L=E+qV.
 \label{eq:charged-hamiltonian}
\end{equation}
Using \eqref{eq:hamiltonian-energy-balance} in \eqref{eq:charged-hamiltonian},
\begin{align*}
 \frac{\dd E}{\dd t}+q\partial_tV+q\bv\cdot\pdr V
 &=-\partial_tL=q\partial_tV-q\bv\cdot\partial_t\bA.
\end{align*}
\begin{equation}
 \frac{\dd E}{\dd t}=q\bv\cdot\bE.
 \label{eq:lorentz-energy}
\end{equation}
This is the mechanical work rate~\cite[pp.~115--116]{tong-em}.  Thus $\bp+q\bA$ and $H$ are canonical momentum and energy; $\bp,E$
are mechanical momentum and energy, with $E$ including rest energy.
Canonical energy is conserved for static potentials in the chosen gauge.
A time-dependent gauge transformation can change $H$ while leaving $E$
and the force unchanged.

Now expand the complex product of the field and proper velocity:
\begin{equation}
 \begin{aligned}
 \frac{FU}{\gamma}={}&\bv\cdot\bE
   +(\bE+\bv\times\bBf)\cdot\bsig\\
 &+\ii\left(\bv\cdot\bBf
   +(\bBf-\bv\times\bE)\cdot\bsig\right).
 \end{aligned}
 \label{eq:force-product}
\end{equation}
The real scalar and vector parts reproduce
\eqref{eq:lorentz-energy} and \eqref{eq:lorentz-momentum}.
Since $P$ is real and $\dd/\dd s=\gamma\dd/\dd t$, they combine into
\begin{equation}
 \boxed{\frac{\dd P}{\dd s}
 =q\repart{FU}=\frac q2(FU+UF^{\mathsf H}).}
 \label{eq:lorentz}
\end{equation}
Hermitian extraction enforces real mechanical
momentum and reproduces the action-derived force.  The imaginary part
need not vanish for an electric charge.

\subsection{Acceleration and the mass shell}
\label{sec:lorentz-proper-acceleration}

Recall the real proper velocity $U=\dd X/\dd s$, a
\hyperref[eq:unit-paravector]{unit paravector} with $UU^{-1}=1$ and
$U^{-1}=U^*$, real momentum $P=mU$, and coordinate acceleration
$\ba=\dd\bv/\dd t$.  Scalar-part cyclicity proves tangency to the mass shell:
\begin{align*}
 \scpart{U^*\dd P/\dd s}
 &=\frac q2\scpart{U^*FU+F^{\mathsf H}}\\
 &=\frac q2\left(\scpart F+\scpart{F^{\mathsf H}}\right)=0.
\end{align*}
This agrees with $\scpart{U^*\dd U/\dd s}=0$ in
Section~\ref{sec:proper-acceleration}.
Combining \eqref{eq:proper-acceleration-components} with the force components
in \eqref{eq:force-product} gives
\begin{align*}
 m\gamma^3\bv\cdot\ba&=q\bv\cdot\bE,\\
 m\gamma\left(\ba+\gamma^2(\bv\cdot\ba)\bv\right)
 &=q(\bE+\bv\times\bBf).
\end{align*}
Eliminating $\bv\cdot\ba$ yields the coordinate acceleration
\begin{equation}
 \boxed{\ba=\frac{q}{m\gamma}
 \left(\bE+\bv\times\bBf-(\bv\cdot\bE)\bv\right).}
 \label{eq:lorentz-coordinate-acceleration}
\end{equation}
The magnetic term changes momentum without doing work.
Taking the determinant of the real force gives
\begin{equation}
 \begin{aligned}
 -\det\left(\frac{\dd P}{\dd s}\right)
 &=-m^2\det\left(\frac{\dd U}{\dd s}\right)\\
 &=q^2\gamma^2\left((\bE+\bv\times\bBf)^2-(\bv\cdot\bE)^2\right).
 \end{aligned}
 \label{eq:proper-force-magnitude}
\end{equation}
Dividing by $m^2$ gives the invariant squared proper acceleration.
The factor $\gamma^2((\bE+\bv\times\bBf)^2-(\bv\cdot\bE)^2)$
is the squared electric field in the particle's instantaneous rest frame.

\subsection{Energy and momentum flux}

Recall $F=(\bE+\ii\bBf)\cdot\bsig$.  Its Hermitian product is
\begin{equation}
 \tfrac12FF^{\mathsf H}
 =\tfrac12(\bE^2+\bBf^2)+(\bE\times\bBf)\cdot\bsig.
 \label{eq:field-energy-flux}
\end{equation}
Define the real energy density $u=(\bE^2+\bBf^2)/2$ and real Poynting
vector $\bS=\bE\times\bBf$.  These coefficients are frame-dependent,
whereas $F^2$ contains the Lorentz invariants.
The full stress--energy tensor is represented by the real linear map
\begin{equation}
 T(N):=\tfrac12FNF^{\mathsf H},
 \label{eq:electromagnetic-stress-map}
\end{equation}
with real paravector input $N$.  Its time-direction value is
$T(1)=u+\bS\cdot\bsig$.  It is future causal because $u\geq0$ and
\begin{equation}
 \det(T(1))=u^2-\bS^2=\tfrac14\lVert\det F\rVert^2\geq0.
 \label{eq:stress-causal-bound}
\end{equation}
Equality holds for null fields.  Reversing the order gives
$F^{\mathsf H}F=(FF^{\mathsf H})^*$: the same energy density and opposite
flux.  The transformation of the full map is derived in
Section~\ref{sec:transformations}.

Using Maxwell's vector equations,
\begin{align*}
 \partial_tu
 &=\bE\cdot(\pdr\times\bBf-\bJ)-\bBf\cdot(\pdr\times\bE)\\
 &=-\pdr\cdot\bS-\bJ\cdot\bE,
\end{align*}
where
$\pdr\cdot(\bE\times\bBf)
=\bBf\cdot(\pdr\times\bE)-\bE\cdot(\pdr\times\bBf)$.
Thus Poynting's energy-balance law is
\begin{equation}
 \partial_tu+\pdr\cdot\bS=-\bJ\cdot\bE.
 \label{eq:poynting}
\end{equation}
Equivalently, the left side is
$\tfrac12\partial_t\scpart{FF^{\mathsf H}}
+\tfrac12\pdr\cdot\vecpart{FF^{\mathsf H}}$.
The right side is the field's loss of energy to matter per unit volume.

\section{Relativistic wave equations}
\label{sec:waves}

Here $X=t+\br\cdot\bsig$ has real coordinates $(t,\br)\in\R^4$,
$m>0$ is mass, $q\in\R$ is charge, and $\hbar=c=1$.
Amplitudes are complex functions of these real coordinates;
their scalar or paravector type is specified below.

\subsection{Amplitudes and their pairing}
\label{sec:quantum-amplitudes}
\label{sec:paravector-inner-products}
\label{sec:axis-ideal-basis}

For complex paravectors $\psi,\psi'$ and a left-acting paravector $Z$,
define the complex scalar pairings
\begin{align}
 \langle\psi'\mid\psi\rangle
 &:=\scpart{\psi'^{\mathsf H}\psi},
 \label{eq:paravector-inner-product}\\
 \langle\psi'\mid Z\mid\psi\rangle
 &:=\scpart{\psi'^{\mathsf H}Z\psi}.
 \label{eq:paravector-matrix-element}
\end{align}
Their real norm density is
\begin{equation}
 \rho:=\scpart{\psi^{\mathsf H}\psi}=\lVert\psi\rVert^2\geq0.
 \label{eq:probability-density}
\end{equation}
For a normalized Pauli wavefunction, $\int\rho\,\dd^3\br=1$;
Dirac uses the sum over its two entries.  This probability interpretation
does not apply to an arbitrary scalar KG field.

For $\rho>0$ and a real observable $Z^{\mathsf H}=Z$, its mean is
\begin{equation}
 \langle Z\rangle_\psi
 :=\frac{\scpart{\psi^{\mathsf H}Z\psi}}{\rho}.
 \label{eq:paravector-expectation}
\end{equation}
For a spatial wavefunction, integrate numerator and denominator.
For differential observables, derivatives act to the right and boundary
terms must vanish when taking the Hermitian conjugate.

For arbitrary constant perpendicular \hyperref[def:real-unit-vector]{real unit vectors}
$\bu,\ba$, choose a constant complex paravector $\psi_+$ of unit norm
with $(\bu\cdot\bsig)\psi_+=\psi_+$, and set $\psi_-=(\ba\cdot\bsig)\psi_+$.
Then
\begin{equation}
 (\bu\cdot\bsig)\psi_\pm=\pm\psi_\pm.
 \label{eq:axis-ideal-basis}
\end{equation}
The pair is orthonormal: $\langle\psi_a\mid\psi_b\rangle=\delta_{ab}$
for $a,b=+,-$.  A Pauli amplitude or either Dirac entry is
\begin{equation}
 \psi(t,\br)=c_+(t,\br)\psi_++c_-(t,\br)\psi_-.
 \label{eq:axis-ideal-components}
\end{equation}
Here $c_+,c_-:\R^4\to\C$ and
$\rho=\lVert c_+\rVert^2+\lVert c_-\rVert^2$.
Both Dirac entries use the same fixed pair.  Its construction and matrix
verification are in Appendix~\ref{sec:matrix-spinors}.

\Needspace{10\baselineskip}
\subsection{Energy and momentum as differential operators}
\label{sec:quantum-energy-momentum}

Recall $X=t+\br\cdot\bsig$, with real coordinates $(t,\br)\in\R^4$.
Temporarily restore Planck's constant $h>0$ and $\hbar=h/(2\pi)$,
keeping $c=1$.  For a free wave, let $\omega\in\R$ be angular frequency
($2\pi$ times cycles per unit time), and let $\bk\in\R^3$ be its
wavenumber vector in cycles per unit length.  For $\bk\ne0$, the wavelength
$\lambda>0$ satisfies $\lVert\bk\rVert=1/\lambda$.

The physical input is the Planck--Einstein energy relation and de Broglie
momentum relation~\cite[pp.~248--249]{debroglie1929}:
\begin{align}
 E&=\hbar\omega,
 \label{eq:planck-einstein}\\
 \bp&=h\bk=2\pi\hbar\bk.
 \label{eq:debroglie-momentum}
\end{align}
Here $E\in\R$ is total energy, including rest energy, and
$\bp\in\R^3$ is momentum; its magnitude is $h/\lambda$.
Photoelectric measurements test the energy relation~\cite[pp.~61--62]{millikan1924};
electron diffraction tests the momentum relation~\cite[p.~255]{debroglie1929}.

For the constant real four-momentum $P=E+\bp\cdot\bsig$, the
complex scalar plane wave is
\[
 \begin{aligned}
 f(X,P)&:=\exp\!\left(-\ii(\omega t-2\pi\bk\cdot\br)\right)\\
       &=\exp\!\left(-\frac{\ii}{\hbar}(Et-\bp\cdot\br)\right).
 \end{aligned}
\]
The second line substitutes \eqref{eq:planck-einstein} and
\eqref{eq:debroglie-momentum}.  Differentiating gives
\begin{align*}
 \ii\hbar\partial_t f&=Ef,\\
 -\ii\hbar\pdr f&=\bp f.
\end{align*}
Thus $\ii\hbar\partial_t$ and $-\ii\hbar\pdr$ represent energy and
momentum on each mode, and extend to superpositions by linearity.
Returning to $\hbar=1$ means $h=2\pi$, giving
\begin{equation}
 f(X,P)=\exp\!\left(-\ii(Et-\bp\cdot\br)\right).
 \label{eq:plane-wave}
\end{equation}
The two eigenvalue equations combine into
\begin{equation}
 \ii\partial^*f=Pf.
 \label{eq:momentum-operator-eigenvalues}
\end{equation}

In a field, recall the real potential paravector
$\Phi(t,\br)=V(t,\br)+\bA(t,\br)\cdot\bsig$:
$V$ is a real scalar potential and $\bA$ a real vector potential.
The real charge is $q$.  The particle action gives canonical energy
$E+qV$ and momentum $\bp+q\bA$; identifying them with translation
operators defines the gauge-covariant mechanical operators
\begin{align}
 E&:=\ii\partial_t-qV,
 \label{eq:mechanical-energy-operator}\\
 \bp&:=-\ii\pdr-q\bA.
 \label{eq:mechanical-momentum-operator}
\end{align}
Here $E$ includes rest energy.  Below, $E,\bp$ denote these differential
operators unless real mode eigenvalues are specified.  They act to the
right and combine as
\begin{align}
 E-\bp\cdot\bsig&=\ii\partial-q\Phi^*,
 \label{eq:covariant-momentum-adjugate}\\
 E+\bp\cdot\bsig&=\ii\partial^*-q\Phi.
 \label{eq:covariant-momentum-paravectors}
\end{align}
Operator order matters in a field; an ordinary paravector determinant
cannot replace their product.  For free commuting derivatives,
$\Box=\partial_t^2-\pdr^2$ gives
\begin{equation}
 (\ii\partial)(\ii\partial^*)=-\Box.
 \label{eq:free-momentum-product}
\end{equation}
Reversing the order gives the same result.

\subsubsection{Field strength from noncommuting derivatives}
Recall the real electric and magnetic fields from
\eqref{eq:electric-potential} and \eqref{eq:magnetic-potential}, and
$S=\partial_tV+\pdr\cdot\bA$, the real scalar in
\eqref{eq:potential-field}.  The real potential is $\Phi=V+\bA\cdot\bsig$;
$F=(\bE+\ii\bBf)\cdot\bsig$ is a complex vector paravector.

Apply the operators to a differentiable complex amplitude.
The product rule cancels all derivatives of the amplitude in the
commutators; for Cartesian indices $j,k=1,2,3$,
\begin{align*}
 [E,p_j]&=-\ii q(\partial_tA_j+\partial_jV)=\ii qE_j,\\
 [p_j,p_k]&=\ii q(\partial_jA_k-\partial_kA_j),\\
 \bp\times\bp&=\ii q\bBf.
\end{align*}
Sigma multiplication and the energy commutator consequently give
\begin{equation}
 (\bp\cdot\bsig)^2=\bp^2-q\bBf\cdot\bsig.
 \label{eq:kinetic-sigma-square}
\end{equation}
\begin{equation}
 [E,\bp\cdot\bsig]=\ii q\bE\cdot\bsig.
 \label{eq:kinetic-electric-commutator}
\end{equation}
Here $E_j$ is an electric-field component.  Using
\eqref{eq:kinetic-sigma-square} and \eqref{eq:kinetic-electric-commutator}
in the ordered products gives
\begin{align}
 (E+\bp\cdot\bsig)(E-\bp\cdot\bsig)
 &=E^2-(\bp\cdot\bsig)^2-[E,\bp\cdot\bsig]\notag\\
 &=E^2-\bp^2-\ii qF,
 \label{eq:covariant-product-plus}\\[0.35em]
 (E-\bp\cdot\bsig)(E+\bp\cdot\bsig)
 &=E^2-(\bp\cdot\bsig)^2+[E,\bp\cdot\bsig]\notag\\
 &=E^2-\bp^2-\ii qF^*.
 \label{eq:covariant-product-minus}
\end{align}
Their common scalar part gives KG; the vector terms give the Dirac
spin coupling.


\subsection{Klein--Gordon equation}

Let $\psi:\R^4\to\C$ be a scalar field and $m>0$ its real mass.
For a free real four-momentum, $\det P=m^2$ means
$E^2-\bp^2=m^2$.  Replacing its coefficients by the translation
operators in \eqref{eq:momentum-operator-eigenvalues} gives
\[
 \bigl((\ii\partial)(\ii\partial^*)-m^2\bigr)\psi=0.
\]
Using \eqref{eq:free-momentum-product}, this is the Klein--Gordon equation~\cite[p.~9]{tong-qft}:
\begin{equation}
 \boxed{(\Box+m^2)\psi=0.}
 \label{eq:kg}
\end{equation}

\subsubsection{Positive and negative frequency}
Recall the complex scalar field $\psi(t,\br)$ and real mass $m>0$.
Define the positive scalar operator
\[
 \Omega:=\sqrt{m^2-\pdr^2},\qquad
 \Omega^2=\det(m+\pdr\cdot\bsig).
\]
Since $\Omega$ commutes with $\partial_t$, KG factors as
\begin{equation}
 \Box+m^2=\partial_t^2+\Omega^2
        =(\partial_t-\ii\Omega)(\partial_t+\ii\Omega).
 \label{eq:kg-frequency-factorization}
\end{equation}
The reverse product is the same.  Its two frequency branches obey
\[
 (\partial_t\pm\ii\Omega)\psi_\pm=0,
 \qquad
 \psi_\pm(t,\br)=e^{\mp\ii\Omega t}\psi_\pm(0,\br).
\]
Their sum gives the solution for arbitrary complex initial branch fields
$\psi_\pm(0,\br)$:
\begin{equation}
 \boxed{\psi(t,\br)=e^{-\ii\Omega t}\psi_+(0,\br)
                         +e^{\ii\Omega t}\psi_-(0,\br).}
 \label{eq:kg-general-solution}
\end{equation}

\subsubsection{Square roots from Fourier series}
For spatially periodic fields, write the discrete Fourier series
\[
 \psi=\sum_n c_n f_n,\qquad
 f_n(t,\br)=\exp\!\left(\ii(\omega_n t+2\pi\bk_n\cdot\br)\right).
\]
Here $n$ labels modes, $c_n\in\C$ are constant coefficients,
$\omega_n\in\R$ are signed angular phase frequencies, and $\bk_n\in\R^3$
are wavenumber vectors in cycles per unit length.  For $\bk_n\ne0$,
its real wavelength $\lambda_n>0$ satisfies $\lVert\bk_n\rVert=1/\lambda_n$.
By \eqref{eq:debroglie-momentum}, the mode momentum is $2\pi\bk_n$
in our units.  The $+\omega_n t$ convention gives energy
eigenvalue $-\omega_n$.
Differentiation gives
\begin{align*}
 \Box\psi&=\sum_n(4\pi^2\bk_n^2-\omega_n^2)c_n f_n,\\
 (\Box+m^2)\psi&=\sum_n(4\pi^2\bk_n^2+m^2-\omega_n^2)c_n f_n.
\end{align*}
Thus each nonzero KG coefficient requires
\[
 \omega_n^2=4\pi^2\bk_n^2+m^2,\qquad
 \omega_n=\pm\sqrt{4\pi^2\bk_n^2+m^2}.
\]
Both signs are independent modes~\cite[pp.~22--23]{tong-qft}; no common time period is assumed.
A square-root operator acts on the same coefficients by
\[
 \sqrt\Box\,\psi=\sum_n\sqrt{4\pi^2\bk_n^2-\omega_n^2}\,c_n f_n.
\]
Choose one square-root branch consistently.  Squaring restores $\Box$, giving
\begin{equation}
 \Box+m^2=(m-\ii\sqrt\Box)(m+\ii\sqrt\Box).
 \label{eq:kg-root-factorization}
\end{equation}
For the positive spatial root, the same series gives
\begin{align*}
 \Omega^2\psi&=\sum_n(m^2+4\pi^2\bk_n^2)c_n f_n,\\
 \Omega\psi&=\sum_n\sqrt{m^2+4\pi^2\bk_n^2}\,c_n f_n.
\end{align*}
Thus $e^{\mp\ii\Omega t}$ multiplies each spatial mode by
$e^{\mp\ii\sqrt{m^2+4\pi^2\bk_n^2}\,t}$.
Unlike $\Omega$, $\sqrt\Box$ has imaginary eigenvalues on KG modes;
its branch does not select the frequency sign.

\subsubsection{Potential}
Recall the real potentials $V,\bA$, charge $q$, and kinetic operators
$E=\ii\partial_t-qV$, $\bp=-\ii\pdr-q\bA$.
The scalar equation in a field is
\begin{equation}
 \boxed{(-E^2+\bp^2+m^2)\psi=0.}
 \label{eq:kg-field}
\end{equation}
The mass-shell equation uses the scalar part $E^2-\bp^2$ of
\eqref{eq:covariant-product-plus} or \eqref{eq:covariant-product-minus}.
The product rule gives
\begin{align}
 E^2&=-\partial_t^2-\ii q(\partial_tV)
       -2\ii qV\partial_t+q^2V^2,
 \label{eq:energy-operator-square}\\
 \bp^2&=-\pdr^2+\ii q(\pdr\cdot\bA)
       +2\ii q\bA\cdot\pdr+q^2\bA^2.
 \label{eq:momentum-operator-square}
\end{align}
Substituting \eqref{eq:energy-operator-square} and
\eqref{eq:momentum-operator-square} into \eqref{eq:kg-field} gives
\begin{equation}
 \boxed{\left(\Box+\ii qS-q^2\det\Phi+m^2
 +2\ii q(V\partial_t+\bA\cdot\pdr)\right)\psi=0.}
 \label{eq:kg-potential-expanded}
\end{equation}

All potential derivatives are retained~\cite[p.~139]{tong-qft}.

\subsection{Dirac equation}
\label{sec:dirac-development}

\subsubsection{Block form and free factorization}

Recall $m>0$ and real space--time coordinates $(t,\br)$.
A Dirac field is the two-entry block
\begin{equation}
 \bPsi(t,\br)=\begin{Bmatrix}\Psi_1(t,\br)\\\Psi_2(t,\br)\end{Bmatrix}.
 \label{eq:dirac-block}
\end{equation}
Each entry is a complex paravector amplitude of \eqref{eq:axis-ideal-components};
the block therefore has four independent complex components.
On blocks, Hermitian conjugation also transposes entries; sigma
conjugation acts entrywise.  The local pairing is
\begin{equation}
 \langle\bPsi'\mid\bPsi\rangle
 :=\sum_{k=1}^{2}\scpart{\Psi_k'^{\mathsf H}\Psi_k}.
 \label{eq:dirac-inner-product}
\end{equation}
This gives the positive real density
\begin{equation}
 \lVert\bPsi\rVert^2
 =\scpart{\Psi_1^{\mathsf H}\Psi_1}
 +\scpart{\Psi_2^{\mathsf H}\Psi_2}.
 \label{eq:dirac-norm}
\end{equation}

For a complex paravector $Z$, define the complex-linear block map
\begin{equation}
 \Wop(Z):=\begin{Bmatrix}0&Z\\\adj Z&0\end{Bmatrix}.
 \label{eq:dirac-block-map}
\end{equation}
Here $\bI$ is the two-entry block identity.  Applying $\Wop(Z)$ twice
multiplies the first entry by $Z\adj Z$ and the second by $\adj Z\,Z$.
Both products are the paravector determinant, so
\begin{equation}
 \Wop(Z)^2=\det(Z)\bI.
 \label{eq:dirac-block-square}
\end{equation}
Sigma conjugation on each block entry gives $\Wop(Z)^*=\Wop(Z^*)$;
Hermitian conjugation also transposes the block:
\begin{equation}
 \Wop(Z)^{\mathsf H}=\Wop(Z^*).
 \label{eq:dirac-block-adjoint}
\end{equation}
Thus $\Wop(Z)$ is not generally Hermitian.  For the free field define
\begin{align*}
 \bDm&:=\Wop(\ii\partial)-m\bI,\\
 \bDm'&:=\Wop(\ii\partial)+m\bI.
\end{align*}
The Dirac equation is~\cite[p.~90]{tong-qft}
\begin{equation}
 \boxed{\bDm\bPsi=0.}
 \label{eq:dirac}
\end{equation}
Since $\Wop(\ii\partial)^2=-\Box\bI$, its mass conjugate gives
\begin{equation}
 \bDm'\bDm=\Wop(\ii\partial)^2-m^2\bI
            =-(\Box+m^2)\bI.
 \label{eq:dirac-free-factorization}
\end{equation}
Applying \eqref{eq:dirac-free-factorization} to \eqref{eq:dirac} gives
$(\Box+m^2)\bPsi=0$: every Dirac component satisfies KG.  The original first-order
constraint is the coupled pair
\begin{align}
 \ii\partial\Psi_2&=m\Psi_1,\label{eq:dirac-free-first}\\
 \ii\partial^*\Psi_1&=m\Psi_2.\label{eq:dirac-free-second}
\end{align}

\subsubsection{Reconstructing free solutions from Klein--Gordon}
Let $\psi':\R^4\to\C$ solve KG and recall the constant basis amplitude
$\psi_+$.  For constant complex scalars $a,b$, a two-constant construction is
\begin{equation}
 \bPsi=\frac1m
 \begin{Bmatrix}
  ma+\ii b\partial\\
  \ii a\partial^*+mb
 \end{Bmatrix}\psi'\psi_+.
 \label{eq:kg-dirac-general-lift}
\end{equation}
Substituting \eqref{eq:kg-dirac-general-lift} into
\eqref{eq:dirac-free-first} and \eqref{eq:dirac-free-second} gives
\begin{align*}
 \ii\partial\Psi_2-m\Psi_1
    &=-\frac{a}{m}(\Box+m^2)\psi'\psi_+=0,\\
 \ii\partial^*\Psi_1-m\Psi_2
    &=-\frac{b}{m}(\Box+m^2)\psi'\psi_+=0.
\end{align*}
Taking $a=1,b=0$ gives one scalar-seeded solution.  To span every
Dirac field, instead let $\psi$ have the form \eqref{eq:axis-ideal-components},
with both scalar coefficients solving KG.  Then
\begin{equation}
 \boxed{\bPsi=\begin{Bmatrix}\psi\\\ii\partial^*\psi/m\end{Bmatrix}.}
 \label{eq:kg-dirac-complete-lift}
\end{equation}
The second Dirac equation holds by construction; the first is KG.
Conversely, every massive free solution has this form with $\psi=\Psi_1$.
Equation~\eqref{eq:kg-general-solution} evolves its two coefficients.

On one Fourier mode, let $s\in\C$ satisfy $\bDm\bPsi=s\bPsi$
for a nonzero amplitude.  Here $\Box$ denotes its scalar mode eigenvalue.
Since $\bDm+m\bI=\Wop(\ii\partial)$, \eqref{eq:dirac-free-factorization} gives
\[
 (\bDm+m\bI)^2\bPsi=-\Box\bPsi
 \quad\Longrightarrow\quad
 \bigl((s+m)^2+\Box\bigr)\bPsi=0.
\]
Thus the eigenvalues are
\[
 s=-m\pm\ii\sqrt\Box.
\]
When distinct, each allows two independent complex coefficients in
$\Psi_1$, with $\Psi_2=\ii\partial^*\Psi_1/(s+m)$.
Each root therefore occurs twice, giving the scalar characteristic equation
\begin{equation}
 \boxed{\bigl((s+m)^2+\Box\bigr)^2=0.}
 \label{eq:dirac-characteristic}
\end{equation}
At $\Box=0$ the roots coalesce with total algebraic multiplicity four.
For KG modes, $\Box=-m^2$, so $s=0,-2m$; the Dirac equation selects $s=0$.
The matrix verification is \eqref{eq:matrix-dirac-characteristic}.

\subsubsection{Potential and paravector components}
Recall the real $\Phi=V+\bA\cdot\bsig$, real charge $q$, and
$E=\ii\partial_t-qV$, $\bp=-\ii\pdr-q\bA$.  Minimal coupling gives
\begin{equation}
 \boxed{\bigl(\Wop(\ii\partial-q\Phi^*)-m\bI\bigr)\bPsi=0.}
 \label{eq:dirac-field}
\end{equation}
Its entries are
\begin{align}
 (E-\bp\cdot\bsig)\Psi_2&=m\Psi_1,\label{eq:dirac-entry-one}\\
 (E+\bp\cdot\bsig)\Psi_1&=m\Psi_2.\label{eq:dirac-entry-two}
\end{align}
Isolating the time derivative gives
$\ii\partial_t\bPsi=H\bPsi$, where the Hamiltonian is
\begin{equation}
 H:=qV\bI+\Wop(1)\bigl(m\bI+\Wop(\bp\cdot\bsig)\bigr).
 \label{eq:dirac-hamiltonian}
\end{equation}
Real potentials and vanishing boundary terms make $H$ Hermitian in the
integrated pairing, so it conserves the norm \eqref{eq:dirac-norm}.

To expose each coefficient, write
$\Psi_k=a_k+\bb_k\cdot\bsig$ for $k=1,2$, with complex scalar
functions $a_k:\R^4\to\C$ and complex vector functions
$\bb_k:\R^4\to\C^3$.  Sigma multiplication gives
\begin{align}
 Ea_2-\bp\cdot\bb_2&=ma_1,\label{eq:dirac-scalar-one}\\
 E\bb_2-\bp a_2-\ii\bp\times\bb_2&=m\bb_1,\label{eq:dirac-vector-one}\\
 Ea_1+\bp\cdot\bb_1&=ma_2,\label{eq:dirac-scalar-two}\\
 E\bb_1+\bp a_1+\ii\bp\times\bb_1&=m\bb_2.\label{eq:dirac-vector-two}
\end{align}
Every operator acts on the following coefficient.  Each entry still
has the two independent complex components of \eqref{eq:axis-ideal-components}.
Appendix~\ref{sec:matrix-spinors} converts the system to
the conventional four-component Dirac equation~\cite[pp.~101, 137--139]{tong-qft}.

\subsubsection{Interacting factorization and reconstruction}
\label{sec:interacting-factorization}

Recall the real potential $\Phi=V+\bA\cdot\bsig$, the complex field
$F=(\bE+\ii\bBf)\cdot\bsig$, and
$S=\partial_tV+\pdr\cdot\bA$.
For any complex paravector field $\psi(t,\br)$, the product rule gives
\begin{align}
 \partial^*(\Phi^*\psi)+\Phi\partial\psi
 &=(S+F)\psi+2(V\partial_t+\bA\cdot\pdr)\psi,
 \label{eq:potential-product-one}\\
 \partial(\Phi\psi)+\Phi^*\partial^*\psi
 &=(S+F^*)\psi+2(V\partial_t+\bA\cdot\pdr)\psi.
 \label{eq:potential-product-two}
\end{align}
The difference of \eqref{eq:potential-product-one} and
\eqref{eq:potential-product-two} is $(F-F^*)\psi=2(\bE\cdot\bsig)\psi$.

Applying the block map $\Wop$ twice gives
\begin{align*}
 \Wop(\partial)\Wop(\Phi^*)\bPsi
   &=\begin{Bmatrix}\partial(\Phi\Psi_1)\\\partial^*(\Phi^*\Psi_2)\end{Bmatrix},\\
 \Wop(\Phi^*)\Wop(\partial)\bPsi
   &=\begin{Bmatrix}\Phi^*\partial^*\Psi_1\\\Phi\partial\Psi_2\end{Bmatrix},\\
 \Wop(\Phi^*)^2&=\det\Phi\,\bI.
\end{align*}
Differentiating a parenthesized product includes the potential derivative.
Define the interacting mass-conjugate blocks
\begin{align*}
 \bDm&:=\Wop(\ii\partial-q\Phi^*)-m\bI,\\
 \bDm'&:=\Wop(\ii\partial-q\Phi^*)+m\bI.
\end{align*}
Multiplying and using the preceding mixed-derivative identities gives
\begin{equation}
 \bDm'\bDm\bPsi=
 -\begin{Bmatrix}
  (-E^2+\bp^2+m^2+\ii qF^*)\Psi_1\\
  (-E^2+\bp^2+m^2+\ii qF)\Psi_2
 \end{Bmatrix}.
 \label{eq:dirac-field-factorization}
\end{equation}
This recovers
\eqref{eq:dirac-free-factorization} at $q=0$.  With coupling, its entries obey
\begin{align*}
 \left((\partial+\ii q\Phi^*)(\partial^*+\ii q\Phi)+m^2\right)\Psi_1&=0,\\
 \left((\partial^*+\ii q\Phi)(\partial+\ii q\Phi^*)+m^2\right)\Psi_2&=0.
\end{align*}
Expanding gives the exact second-order equations
\begin{align}
 (-E^2+\bp^2+m^2+\ii qF^*)\Psi_1&=0,
 \label{eq:dirac-squared-first}\\
 (-E^2+\bp^2+m^2+\ii qF)\Psi_2&=0.
 \label{eq:dirac-squared-field}
\end{align}
For the spin-term components, write $\bC=\bE+\ii\bBf$ and
$\psi=a+\bb\cdot\bsig$, with complex functions
$a:\R^4\to\C$ and $\bb,\bC:\R^4\to\C^3$.
Here $\bC^*$ conjugates its scalar coefficients.  Sigma multiplication gives
\begin{align}
 \ii qF\psi
 &=\ii q\bC\cdot\bb+
   (\ii qa\bC-q\bC\times\bb)\cdot\bsig,
 \label{eq:spin-field-components}\\
 \ii qF^*\psi
 &=-\ii q\bC^*\cdot\bb+
   (-\ii qa\bC^*+q\bC^*\times\bb)\cdot\bsig.
 \label{eq:spin-field-star-components}
\end{align}
To construct solutions, take a complex paravector amplitude $\psi$ of
\eqref{eq:axis-ideal-components} satisfying
$(-E^2+\bp^2+m^2+\ii qF)\psi=0$.  Set
\begin{equation}
 \boxed{\bPsi=\begin{Bmatrix}(E-\bp\cdot\bsig)\psi/m\\\psi\end{Bmatrix}.}
 \label{eq:dirac-field-reconstruction}
\end{equation}
The first Dirac equation holds by construction; the second is the seed
equation.  Every massive solution has this form with $\psi=\Psi_2$.
One must solve that seed equation for the chosen potential: a scalar KG
seed lacks the spin term, and independently solved squared entries need
not satisfy the first-order coupling.

\subsubsection{Gauge phase}
Recall the real scalar $\lambda:\R^4\to\R$ and gauge transformation
$\Phi'=\Phi+\partial^*\lambda$, or
$V'=V+\partial_t\lambda$, $\bA'=\bA-\pdr\lambda$.
The extra derivatives of $\lambda$ are canceled by differentiating a
scalar phase.  Write $\bPsi=e^{\ii q\lambda}\bPsi'$; entrywise,
\begin{align*}
 (\ii\partial-q\Phi^*)(e^{\ii q\lambda}\Psi_k')
 &=e^{\ii q\lambda}
   (\ii\partial-q\partial\lambda-q\Phi^*)\Psi_k'\\
 &=e^{\ii q\lambda}(\ii\partial-q\Phi'^*)\Psi_k'.
\end{align*}
The same product rule for $\ii\partial^*-q\Phi$ gives the other block entry.
Thus the new amplitude is
\begin{equation}
 \bPsi'=e^{-\ii q\lambda}\bPsi.
 \label{eq:dirac-gauge-amplitude}
\end{equation}
With $E'=\ii\partial_t-qV'$ and $\bp'=-\ii\pdr-q\bA'$, the same
product rule gives
\begin{align*}
 E'\bPsi'&=e^{-\ii q\lambda}E\bPsi,\\
 \bp'\bPsi'&=e^{-\ii q\lambda}\bp\bPsi.
\end{align*}
Consequently
\begin{equation}
 \begin{aligned}
 &\bigl(\Wop(\ii\partial-q\Phi'^*)-m\bI\bigr)\bPsi'\\
 &\qquad=e^{-\ii q\lambda}
       \bigl(\Wop(\ii\partial-q\Phi^*)-m\bI\bigr)\bPsi.
 \end{aligned}
 \label{eq:gauge-operator-covariance}
\end{equation}
Equation~\eqref{eq:gauge-operator-covariance} preserves the Dirac equation;
the phase sign is fixed by the potential shift~\cite[p.~137]{tong-qft}.
Its modulus is one, so the norm is unchanged.  The same product rules
preserve the KG and Pauli equations.

\section{Low-energy limits}
\label{sec:low-energy-limits}

Recall the real mass $m>0$, charge $q$, and potentials
$V:\R^4\to\R$, $\bA:\R^4\to\R^3$.
The mechanical operators are $E=\ii\partial_t-qV$ and
$\bp=-\ii\pdr-q\bA$; the fields are
$\bE=-\pdr V-\partial_t\bA$ and $\bBf=\pdr\times\bA$.
Remove the rest oscillation, then iterate the residual energy.
This gives Pauli from Dirac and Schr\"odinger from KG at leading order.

\subsection{Exact reduction and the energy recursion}
\label{sec:low-energy-recursion}

The Dirac entries $\Psi_1,\Psi_2$ are complex paravector amplitudes of
\eqref{eq:axis-ideal-components}, depending on real $(t,\br)$.
At free positive-energy rest, the time equation is
\[
 \ii\partial_t\Psi_k=m\Psi_k
 \quad\Longrightarrow\quad
 \Psi_k(t)=e^{-\ii mt}\Psi_k(0).
\]
A free mode of real energy $\omega$ near $m$ therefore factors as
\[
 e^{-\ii\omega t}=e^{-\ii mt}e^{-\ii(\omega-m)t}.
\]
The envelope has the smaller frequency $\omega-m$.
Write $\Psi_k=e^{-\ii mt}\Psi_k'$, or equivalently
\begin{equation}
 \Psi_k':=e^{\ii mt}\Psi_k,\qquad k=1,2.
 \label{eq:dirac-rest-phase}
\end{equation}
The positive sign removes the rest oscillation by its inverse.
The product rule gives
\[
 (\ii\partial_t-qV)(e^{-\ii mt}\Psi_k')
 =e^{-\ii mt}(m+\ii\partial_t-qV)\Psi_k'.
\]
Hence
\begin{align*}
 E\Psi_k&=e^{-\ii mt}(E+m)\Psi_k',\\
 \bp\Psi_k&=e^{-\ii mt}\bp\Psi_k'.
\end{align*}
The potentials are unchanged; $E$ on the envelope measures mechanical
energy above rest.  Phase removal is exact and preserves the norm and
both branches.  The approximation below selects slow positive-energy
envelopes: a free mode near $-m$ instead has residual frequency near $-2m$.

The coupled equations \eqref{eq:dirac-entry-one} and
\eqref{eq:dirac-entry-two}, after division by $e^{-\ii mt}$, become
\begin{align*}
 (E+m-\bp\cdot\bsig)\Psi_2'&=m\Psi_1',\\
 (E+m+\bp\cdot\bsig)\Psi_1'&=m\Psi_2'.
\end{align*}
Eliminate $\Psi_1'$ before expanding the ordered product:
\[
 \left((E+m+\bp\cdot\bsig)(E+m-\bp\cdot\bsig)-m^2\right)\Psi_2'=0.
\]
Using \eqref{eq:kinetic-sigma-square} and
$[\bp\cdot\bsig,E]=-\ii q\bE\cdot\bsig$ gives
\begin{equation}
 \left(E^2+2mE-\bp^2+q\bBf\cdot\bsig
                   -\ii q\bE\cdot\bsig\right)\Psi_2'=0.
 \label{eq:dirac-residual-squared}
\end{equation}
Rearranging gives the residual-energy recursion, with the
operator order retained:
\begin{equation}
 E\Psi_2'=\frac{\bp^2-E^2-q\bBf\cdot\bsig
                      -[\bp\cdot\bsig,E]}{2m}\Psi_2'.
 \label{eq:dirac-residual-recursion}
\end{equation}
The imaginary electric term in \eqref{eq:dirac-residual-recursion}
belongs to one chiral entry, whose norm alone is not the Dirac
probability.  To derive a probability-preserving Pauli equation,
keep both entries.

At free positive-energy rest, $\Psi_1'=\Psi_2'$: their sum survives and their
difference vanishes.  Define the corresponding complex paravector amplitudes
\begin{align}
 \phi&:=\frac{\Psi_1'+\Psi_2'}{\sqrt2},
 \label{eq:dirac-large-amplitude}\\
 \chi&:=\frac{\Psi_2'-\Psi_1'}{\sqrt2}.
 \label{eq:dirac-small-amplitude}
\end{align}
Adding and subtracting the coupled equations gives the exact pair
\begin{align}
 E\phi&=(\bp\cdot\bsig)\chi,
 \label{eq:dirac-large-component}\\
 (E+2m)\chi&=(\bp\cdot\bsig)\phi.
 \label{eq:dirac-large-small}
\end{align}
Their conserved norm is
$\int(\lVert\phi\rVert^2+\lVert\chi\rVert^2)\dd^3\br$.
Rearrange the second equation for iteration:
\begin{equation}
 \chi=\frac{(\bp\cdot\bsig)\phi}{2m}-\frac{E\chi}{2m}.
 \label{eq:dirac-small-recursion}
\end{equation}
Start with the first term and substitute repeatedly into
\eqref{eq:dirac-small-recursion}, replacing $E\phi$ with
\eqref{eq:dirac-large-component}.  Keep the operator order.
For a real small parameter $0<v\ll1$, assume
$\bp/m=O(v)$, $E/m=O(v^2)$, $qV/m=O(v^2)$,
$q\bBf/m^2=O(v^2)$, and $q\bE/m^2=O(v^3)$, with smooth fields
varying on scales $1/(mv)$ in space and $1/(mv^2)$ in time or longer.
Then $\chi=O(v\phi)$ and each iteration adds a relative $O(v^2)$.

\Needspace{15\baselineskip}
\subsection{Leading order: Pauli}
\label{sec:pauli-limit}

Recall $\phi,\chi$ from \eqref{eq:dirac-large-amplitude} and
\eqref{eq:dirac-small-amplitude}.  Keeping the first iterate gives
\begin{align*}
 \chi&=\frac{\bp\cdot\bsig}{2m}\phi+O(v^3\phi),\\
 E\phi&=\frac{(\bp\cdot\bsig)^2}{2m}\phi+O(mv^4\phi).
\end{align*}
The product identity
$(\bp\cdot\bsig)^2=\bp^2-q\bBf\cdot\bsig$
from \eqref{eq:kinetic-sigma-square} gives Pauli's equation at this order~\cite[pp.~35--36]{foldy-wouthuysen}:
\begin{equation}
 \boxed{\ii\partial_t\phi=\left(
  \frac{(-\ii\pdr-q\bA)^2}{2m}+qV
  -\frac{q}{2m}\bBf\cdot\bsig\right)\phi.}
 \label{eq:pauli}
\end{equation}
Here $\int\scpart{\phi^{\mathsf H}\phi}\dd^3\br=1$ at leading order.
Its two complex components are the Pauli amplitudes; see
Appendix~\ref{sec:matrix-spinors}.  It is the $\psi$ used in the
spin discussion.
Equation~\eqref{eq:spin-half-generator} identifies spin as $\bsig/2$, so the magnetic energy is
$-(q/m)\bBf\cdot(\bsig/2)$: magnetic moment equals $q/m$ times spin,
so the gyromagnetic factor is $2$ before radiative corrections.

\subsubsection{Components along a general field direction}

For a nonzero static uniform field, let
$\bu=\bBf/\lVert\bBf\rVert$ be its
\hyperref[def:real-unit-vector]{real unit direction} and
$\Pi_\pm=(1\pm\bu\cdot\bsig)/2$ its spin projectors.
The complex paravector fields $\phi_\pm=\Pi_\pm\phi$ obey
\begin{equation}
 \ii\partial_t\phi_\pm=\left(
 \frac{\bp^2\mp q\lVert\bBf\rVert}{2m}+qV\right)\phi_\pm.
 \label{eq:pauli-fixed-spin-sectors}
\end{equation}
For a varying field direction, the projectors depend on $(t,\br)$.
Let $H$ denote the Pauli Hamiltonian in \eqref{eq:pauli}.
The product rule gives
\begin{equation}
 \begin{aligned}
 \ii\partial_t(\Pi_\pm\phi)
 &=\ii(\partial_t\Pi_\pm)\phi+\Pi_\pm H\phi\\
 &=H(\Pi_\pm\phi)
   +\left(\ii\partial_t\Pi_\pm
          +\frac{[\Pi_\pm,\bp^2]}{2m}\right)\phi.
 \end{aligned}
 \label{eq:pauli-varying-spin-sectors}
\end{equation}
Only the projector derivatives and kinetic commutator couple the
sectors; the magnetic term commutes with its own projectors.

\subsection{The same leading order: Schr\"odinger}
\label{sec:schrodinger-limit}

When $\bBf=0$, each complex coefficient of the Pauli amplitude obeys
the scalar Schr\"odinger equation.  KG gives the same result directly.
Keep the operators $E,\bp$ defined above.  For the full scalar KG
field $\psi':\R^4\to\C$, remove the rest oscillation as in
\eqref{eq:dirac-rest-phase}, defining its scalar envelope
$\psi:\R^4\to\C$ by
\begin{equation}
 \psi'(t,\br)=e^{-\ii mt}\psi(t,\br).
 \label{eq:kg-rest-phase}
\end{equation}
The product rule gives $E\psi'=e^{-\ii mt}(E+m)\psi$.
Since $m$ is constant, applying $E$ again gives
$E^2\psi'=e^{-\ii mt}(E+m)^2\psi$; spatial derivatives leave the
phase untouched.  Substitution into \eqref{eq:kg-field} therefore gives
\[
 \begin{aligned}
 0&=(-E^2+\bp^2+m^2)\psi'\\
  &=e^{-\ii mt}\bigl(-(E+m)^2+\bp^2+m^2\bigr)\psi\\
  &=e^{-\ii mt}(-E^2-2mE+\bp^2)\psi.
 \end{aligned}
\]
The rest terms cancel.  Divide by the nonzero phase to obtain
\begin{equation}
 E\psi=\frac{\bp^2-E^2}{2m}\psi.
 \label{eq:kg-residual-recursion}
\end{equation}
On a slow positive-energy envelope, $E/m=O(v^2)$, so $E^2\psi$ is
smaller than $2mE\psi$ by $O(v^2)$.  Dropping it at leading order
gives~\cite[pp.~41--42, 139]{tong-qft}
\begin{equation}
 \boxed{\ii\partial_t\psi=\left(
 \frac{(-\ii\pdr-q\bA)^2}{2m}+qV\right)\psi.}
 \label{eq:schrodinger}
\end{equation}
With $\bA=0$, its kinetic operator is $-\pdr^2/(2m)$.
Schr\"odinger and Pauli therefore occur at the same leading order
$1/m$; retaining spin gives the magnetic term in \eqref{eq:pauli}.

\subsection{The next free order}
\label{sec:free-energy-expansion}

For the free complex scalar envelope $\psi(t,\br)$, set $qV=0$ and
$q\bA=0$, so $E$ and $\bp$ commute.
On the positive-frequency KG branch, solving
\eqref{eq:kg-residual-recursion} gives
\begin{equation}
 E\psi=\left(\sqrt{m^2+\bp^2}-m\right)\psi.
 \label{eq:kg-positive-residual-energy}
\end{equation}
The same energy acts on either normalized positive-energy Dirac spin
component.  Iteration starts with $E\psi=\bp^2\psi/(2m)$;
its square supplies the next term:
\begin{equation}
 E\psi=\left(\frac{\bp^2}{2m}-\frac{\bp^4}{8m^3}
             +O\!\left(\frac{\bp^6}{m^5}\right)\right)\psi.
 \label{eq:kg-low-energy-expansion}
\end{equation}

\section{Spin and probability}
\label{sec:spin-and-probability}
\label{sec:paravector-spin-content}

Spin specifies the intrinsic rotation of an amplitude, comparing
corresponding rotated points of real spacetime.  A scalar KG amplitude
stays unchanged and has spin zero; a two-component amplitude has the
one-sided law below.  The wave equation alone does not fix spin:
free Dirac components also satisfy KG.

\subsection{Spin one-half and rotations}
\label{sec:spin-half}

Recall the complex paravector amplitude $\psi(t,\br)$ of
\eqref{eq:axis-ideal-components}.  For a
\hyperref[def:real-unit-vector]{real unit axis} $\bu$ and real angle
$\theta$, use the unitary rotor
$R=\exp(\ii\theta\bu\cdot\bsig/2)$ of
\eqref{eq:rotation-exponential}.  The spinor law is
\begin{equation}
 \psi'=R^{\mathsf H}\psi.
 \label{eq:spin-half-law}
\end{equation}
Writing its infinitesimal action as
$\psi'=\psi-\ii\theta S(\bu)\psi+O(\theta^2)$ defines the spin
operator along $\bu$:
\begin{equation}
 S(\bu)\psi=\tfrac12(\bu\cdot\bsig)\psi.
 \label{eq:spin-half-generator}
\end{equation}
Since $(\bu\cdot\bsig)^2=1$, its two spin projections are
$\pm1/2$ in units $\hbar=1$.  A $2\pi$ rotation gives $R=-1$,
reversing the amplitude's sign; a $4\pi$ rotation restores it.
Probabilities are unchanged by the overall sign, while a relative sign
can be detected by interference~\cite[Sec.~6--3]{feynmanspin}.
Each Dirac entry obeys this same rotation law.

For a real paravector observable $Z$, set $Z'=R^{\mathsf H}ZR$.
Recall $\rho=\lVert\psi\rVert^2>0$ and the scalar expectation
\eqref{eq:paravector-expectation}.  Since $RR^{\mathsf H}=1$,
\begin{equation}
 \langle Z'\rangle_{\psi'}
 =\frac{\scpart{\psi^{\mathsf H}RR^{\mathsf H}ZRR^{\mathsf H}\psi}}{\rho}
 =\langle Z\rangle_\psi.
 \label{eq:rotation-expectation}
\end{equation}
Thus rotating both amplitude and observable preserves the measurement;
$Z=\bu'\cdot\bsig$ recovers the vector rotation of
\eqref{eq:rodrigues} for any real unit vector $\bu'$.

\subsection{Projectors and probabilities}
\label{sec:physical-spinors}

For a real unit measurement axis $\bu$, recall the real projectors
\begin{equation}
 \Pi_\pm=\tfrac12(1\pm\bu\cdot\bsig).
 \label{eq:spin-projectors}
\end{equation}
Their algebra in Section~\ref{sec:algebra-projectors} gives
\begin{align*}
 \psi&=\Pi_+\psi+\Pi_-\psi,\\
 S(\bu)(\Pi_\pm\psi)&=\pm\tfrac12\Pi_\pm\psi,\\
 \langle\Pi_+\psi\mid\Pi_-\psi\rangle&=0,\\
 \lVert\Pi_\pm\psi\rVert^2&=\scpart{\psi^{\mathsf H}\Pi_\pm\psi}.
\end{align*}
The Born probability is the projected squared norm divided by
$\rho=\lVert\psi\rVert^2$.  In the general-axis orthonormal basis
\eqref{eq:axis-ideal-basis}, $\psi=c_+\psi_++c_-\psi_-$ with
complex scalars $c_\pm$, so
$\rho=\lVert c_+\rVert^2+\lVert c_-\rVert^2$ and the probabilities
are $\lVert c_\pm\rVert^2/\rho$.
The same-side and cross-side projector products are already derived in
Section~\ref{sec:algebra-projectors}; the ladder paravectors in
\eqref{eq:ladder-operators} exchange the two sectors.

To evaluate the probabilities directly from paravector components, define
\begin{equation}
 B:=\frac{\psi\psi^{\mathsf H}}{\rho}=1+\bm{s}\cdot\bsig.
 \label{eq:spin-state}
\end{equation}
Here $B$ is real and $\bm{s}\in\R^3$ is the polarization.
The basis construction gives $\det\psi=0$ (Appendix~\ref{sec:matrix-spinors}),
hence $0=\det B=1-\bm{s}^2$: the polarization is a real unit vector.
Cyclicity of the scalar part gives
\[
 \lVert\Pi_\pm\psi\rVert^2
 =\scpart{\psi^{\mathsf H}\Pi_\pm\psi}
 =\rho\,\scpart{\Pi_\pm B}
 =\tfrac12\rho(1\pm\bu\cdot\bm{s}).
\]
Therefore
\begin{equation}
 \boxed{p_\pm(\bu)=\frac{\lVert\Pi_\pm\psi\rVert^2}{\rho}
                    =\tfrac12(1\pm\bu\cdot\bm{s}).}
 \label{eq:spin-probability-intrinsic}
\end{equation}
The mean spin along $\bu$ is
$(p_+-p_-)/2=\bu\cdot\bm{s}/2$.
For a spatial wavefunction, \eqref{eq:spin-probability-intrinsic}
is conditional on position; integrate numerator and denominator to
obtain the whole-packet probability.
If the angle between $\bu$ and $\bm{s}$ is the real angle $\theta$,
\[
 p_+=\cos^2(\theta/2),\qquad p_-=\sin^2(\theta/2).
\]
For $R=\exp(\ii\theta\ba\cdot\bsig/2)$, with real unit rotation
axis $\ba$ and real angle $\theta$, the amplitude law \eqref{eq:spin-half-law} gives
\begin{equation}
 B'=R^{\mathsf H}BR.
 \label{eq:spin-state-rotation}
\end{equation}
Thus the polarization rotates as a vector although its amplitude has
spin one-half.  Rodrigues' formula gives the original-axis overlap:
\[
 \bm{s}\cdot\bm{s}'=(\ba\cdot\bm{s})^2
   +\left(1-(\ba\cdot\bm{s})^2\right)\cos\theta.
\]
For $\ba\cdot\bm{s}=0$, the probability to retain the initial spin
is $\cos^2(\theta/2)$.

\subsection{Spin uncertainty}

Recall the real polarization $\bm{s}$ in \eqref{eq:spin-state}.  For
real unit axes $\bu,\bu'$, set $A=\bu\cdot\bsig/2$ and
$A'=\bu'\cdot\bsig/2$.  Define the real variance by
$\operatorname{var}(A)=\langle A^2\rangle-\langle A\rangle^2$.
The sigma product gives
\begin{align*}
 [A,A']&=\frac{\ii}{2}(\bu\times\bu')\cdot\bsig,\\
 \operatorname{var}(A)&=\frac{1-(\bu\cdot\bm{s})^2}{4}.
\end{align*}
For a normalized amplitude, positivity of
$\lVert(A-\langle A\rangle)\psi+\ii\lambda(A'-\langle A'\rangle)\psi\rVert^2$
for every real scalar $\lambda$ makes its quadratic discriminant
nonpositive.  This yields the Robertson bound
\begin{equation}
 \boxed{\operatorname{var}(A)\;\operatorname{var}(A')
 \geq\tfrac14\lVert\langle[A,A']\rangle\rVert^2
 =\frac{((\bu\times\bu')\cdot\bm{s})^2}{16}.}
 \label{eq:spin-uncertainty}
\end{equation}

\section{Rotations and Lorentz boosts}
\label{sec:transformations}

We now apply the algebraic actions of Section~\ref{sec:preserved-quantities}
to the physical variables already defined.  Primed fields are evaluated at
$X'$ and unprimed fields at $X$.  The frame factors here are constant;
varying factors are treated in Section~\ref{sec:varying-transformations}.

\subsection{Covariance of the physical equations}

Let $T$ be a complex \hyperref[eq:unit-paravector]{unit paravector}
with \hyperref[eq:unit-determinant-paravector]{$\det T=1$}.  It may be any ordered
product of $R=\exp(\ii\theta\bu\cdot\bsig/2)$ and
$L=\exp(\theta\bu\cdot\bsig/2)$, with
\hyperref[def:real-unit-vector]{real unit vector} $\bu$ and real angle
or rapidity $\theta$.  Recall the real event $X$, momentum $P$, potential
$\Phi$, and electric current $D$.
Each transforms by
\begin{equation}
 Z'=T^{-1}Z(T^{-1})^{\mathsf H},
 \qquad Z=X,P,\Phi,D.
 \label{eq:physical-lorentz-congruence}
\end{equation}
The derivative and complex field $F=(\bE+\ii\bBf)\cdot\bsig$ use the corresponding laws
\begin{align}
 \partial'&=T^{\mathsf H}\partial T, \label{eq:derivative-transformation}\\
 F'&=T^{-1}FT. \label{eq:field-transformation}
\end{align}
Equation~\eqref{eq:physical-lorentz-congruence} acts on real four-vectors;
\eqref{eq:derivative-transformation} is the chain rule.  For a
differentiable scalar $f:\R^4\to\C$ with $f'(X')=f(X)$, cyclicity gives
\[
 \dd f=\scpart{(\partial f)\dd X}
 =\scpart{\bigl(T^{\mathsf H}(\partial f)T\bigr)\dd X'}.
\]
Since this holds for every real increment $\dd X'$, the primed derivative
is $\partial'f'=T^{\mathsf H}(\partial f)T$
\cite[p.~101]{tong-em}.

The potential identity~\eqref{eq:potential-field}, with real scalar
$S=\scpart{\partial\Phi}$, then gives
\[
 \partial'\Phi'=T^{\mathsf H}(S+F^*)(T^{\mathsf H})^{-1}.
\]
The scalar part remains $S$; sigma conjugation of the vector part yields
the field similarity~\eqref{eq:field-transformation}.  Since
$T^*=(T^{\mathsf H})^{-1}$, Maxwell covariance follows by cancellation:
\begin{equation}
 \partial'F'=T^{\mathsf H}(\partial F)T
 =T^{\mathsf H}D^*T=D'^*.
 \label{eq:maxwell-covariance}
\end{equation}
This recovers \eqref{eq:maxwell} in the primed frame.
Adjugation of $P'$ and scalar-part cyclicity give
\begin{equation}
 \scpart{\adj{P'}X'}
 =\scpart{T^{\mathsf H}\adj P\,X(T^{\mathsf H})^{-1}}
 =\scpart{\adj P\,X}.
 \label{eq:phase-invariance}
\end{equation}
Thus the plane wave~\eqref{eq:plane-wave} satisfies $f(X',P')=f(X,P)$.

\subsection{Moving frames, clock rates, and lengths}
\label{sec:boosts}

Recall $L=\exp(\theta\bu\cdot\bsig/2)$, with real rapidity $\theta$,
\hyperref[def:real-unit-vector]{real unit vector} $\bu$, $\beta=\tanh\theta$, and $\gamma=\cosh\theta$.
Here $\gamma$ is the factor for frame velocity $\beta\bu$.
The explicit gamma form is \eqref{eq:boost-gamma-form}.
For real $X=t+\br\cdot\bsig$, take coincident inertial origins at
$t=t'=0$.  The passive boost
$X'=L^{-1}XL^{-1}$ means that the primed frame moves at $+\beta\bu$.
Equations~\eqref{eq:boost-scalar-component}--\eqref{eq:boost-vector-component}
become
\begin{align}
 t'&=\gamma(t-\beta\bu\cdot\br), \label{eq:lorentz-time}\\
 \br'&=\br_{\perp}+\gamma(\br_{\parallel}-\beta t\bu),
 \label{eq:lorentz-position}
\end{align}
where $\br_{\parallel}=(\bu\cdot\br)\bu$ and
$\br_{\perp}=\br-\br_{\parallel}$
\cite[p.~96, eqs.~(5.1)--(5.2)]{tong-em}.
Reversing $\theta$ gives $X=LX'L$.

\subsubsection{Spacetime increments and velocity}
Recall the real proper velocity $U=\dd X/\dd s$, a
\hyperref[eq:unit-paravector]{unit paravector} with $UU^{-1}=1$,
along a real world line $X(s)$.  For constant $L$, differentiation gives
\begin{equation}
 \dd X'=L^{-1}\dd X L^{-1}.
 \label{eq:boost-worldline-step}
\end{equation}
Because $\dd s$ is invariant, $U'=L^{-1}UL^{-1}$.  At a point where
$U=L^2$, this gives $U'=1$ and $\dd X'=\dd s$: the boost takes the
particle's instantaneous tangent to rest.  For varying velocity, a
different $L$ is needed at each point; Section~\ref{sec:varying-transformations}
supplies its derivative terms.

Recall the particle velocity $\bv=\dd\br/\dd t$ and resolve it along
$\bu$.  Dividing the spatial increment by
$\dd t'=\gamma(1-\beta\bu\cdot\bv)\dd t$ gives
\begin{equation}
 \boxed{\bv'=\frac{\bv_{\perp}/\gamma+\bv_{\parallel}-\beta\bu}
                       {1-\beta\bu\cdot\bv}.}
 \label{eq:velocity-boost}
\end{equation}
Interval invariance also gives
\begin{equation}
 1-\bv'^2=\frac{1-\bv^2}{\gamma^2(1-\beta\bu\cdot\bv)^2}.
 \label{eq:velocity-causal-bound}
\end{equation}
Thus \eqref{eq:velocity-boost} preserves subluminal motion and light speed
\cite[pp.~135--136, eqs.~(7.22)--(7.23)]{tong-relativity}.
The constant-factor law~\eqref{eq:boost-worldline-step} also commutes with
the world-line integral~\eqref{eq:worldline-integral}.

\subsubsection{Clock rates and measured lengths}
For a clock at rest in the primed frame, \eqref{eq:lorentz-time} uses $\dd\br=\beta\bu\dd t$, hence
\begin{equation}
 \dd t'=\gamma(1-\beta^2)\dd t=\dd t/\gamma.
 \label{eq:time-dilation}
\end{equation}
Here $\dd t'=\dd s$; finite elapsed time follows from
\eqref{eq:elapsed-proper-time} \cite[pp.~118--120]{tong-relativity}.
For a rod at rest in the unprimed frame, let $\Delta\br$ be its endpoint
separation.  A simultaneous measurement in the primed frame requires
$\Delta t'=0$, so $\Delta t=\beta\bu\cdot\Delta\br$.  Substitution gives
\begin{equation}
 \Delta\br'_{\parallel}=\Delta\br_{\parallel}/\gamma.
 \label{eq:length-contraction}
\end{equation}
The transverse separation remains unchanged by \eqref{eq:lorentz-position}.
Thus longitudinal contraction is measured at equal primed times
\cite[pp.~122--123]{tong-relativity}.
Conversely, equal unprimed times give
$\Delta t'=-\gamma\beta\bu\cdot\Delta\br$; simultaneity is frame-dependent.

\subsubsection{Composing frames}
Collinear boosts add rapidities by \eqref{eq:boost-composition}.
For a second real rapidity $\theta'$, set $\beta'=\tanh\theta'$ and
$\gamma'=\cosh\theta'$.  Their combined parameters are
\begin{align}
 \beta''&=\tanh(\theta+\theta')
 =\frac{\beta+\beta'}{1+\beta\beta'}, \label{eq:velocity-addition}\\
 \gamma''&=\gamma\gamma'(1+\beta\beta'). \label{eq:gamma-composition}
\end{align}
These are the collinear velocity and clock-factor composition laws.
Noncollinear boosts instead have the rotation commutator
\eqref{eq:boost-generator-commutator}, producing the Thomas--Wigner rotation.

\subsubsection{Energy and momentum}
For the real momentum $P=E+\bp\cdot\bsig$, the same congruence gives
\begin{align}
 E'&=\gamma(E-\beta\bu\cdot\bp),
 \label{eq:energy-boost}\\
 \bp'&=\bp_{\perp}+\gamma(\bp_{\parallel}-\beta E\bu).
 \label{eq:momentum-boost}
\end{align}
Here the projections are along the boost axis $\bu$.
Both coefficients change, but $E'^2-\bp'^2=m^2$ leaves rest mass unchanged
\cite[pp.~136--138]{tong-relativity}.
The energy increase and force response are derived in
\eqref{eq:relativistic-inertia}--\eqref{eq:relativistic-work}.

\subsection{Boosted fields and energy flux}

For $L=\exp(\theta\bu\cdot\bsig/2)$ with the parameters above,
both the field coefficients and their spacetime arguments transform:
\begin{equation}
 F'(X')=L^{-1}F(LX'L)L.
 \label{eq:boosted-field-arguments}
\end{equation}
Thus $F'(X')=L^{-1}F(X)L$ at corresponding events $X=LX'L$.
Apply \eqref{eq:boost-similarity-vector} to
$\bE+\ii\bBf$ and separate its real and imaginary parts:
\begin{align}
 \bE'&=\bE_{\parallel}+\gamma(\bE_{\perp}+\beta\bu\times\bBf),
 \label{eq:electric-boost}\\
 \bBf'&=\bBf_{\parallel}+\gamma(\bBf_{\perp}-\beta\bu\times\bE).
 \label{eq:magnetic-boost}
\end{align}
Parallel components are unchanged; the projections use the same
\hyperref[def:real-unit-vector]{real unit axis} as
\eqref{eq:lorentz-position}~\cite[p.~108]{tong-em}.
Recall the stress map $T(N)=FNF^{\mathsf H}/2$, with real paravector
input $N$.  Since $N'=L^{-1}NL^{-1}$,
\begin{equation}
 T'(N')=\tfrac12F'N'F'^{\mathsf H}=L^{-1}T(N)L^{-1}.
 \label{eq:stress-covariance}
\end{equation}
Here $T'$ denotes the map formed from $F'$.  Each energy--momentum
flux has the real four-vector law; $T(1)$ alone is only the flux through
the chosen time direction.  The same cancellation in
\eqref{eq:lorentz} gives the transformed proper force.

\newcommand{\referencesummary}{%
\section*{Equation reference}
\addcontentsline{toc}{section}{Equation reference}
\label{sec:reference-summary}
\begingroup\allowdisplaybreaks[1]

Each displayed formula repeats a numbered result; its right-hand number
links to the definition or derivation.  Natural units are $\hbar=c=1$.
The linked text supplies domains, component expansions, and approximation
conditions.

\Needspace{8\baselineskip}
\paragraph{Algebra.}
Recall complex paravectors $Z=S+\bV\cdot\bsig$ and $Z'=S'+\bV'\cdot\bsig$,
with complex scalar $S,S'$ and complex vectors $\bV,\bV'$.
The prime here labels a second element.
\begin{equation*}
 \begin{aligned}
 ZZ'={}&(SS'+\bV\cdot\bV')\\
 &+(S\bV'+S'\bV+\ii\bV\times\bV')\cdot\bsig.
 \end{aligned}\tag*{\eqref{eq:para-product}}
\end{equation*}
\begin{align*}
 Z^{\mathsf H}&=S^*+\bV^*\cdot\bsig
   \tag*{\eqref{eq:algebra-H}}\\
 Z^*&=S^*-\bV^*\cdot\bsig
   \tag*{\eqref{eq:algebra-star}}\\
 \adj Z&=Z^{*\mathsf H}=S-\bV\cdot\bsig
   \tag*{\eqref{eq:adjugate}}\\
 \operatorname{tr}Z&=2\scpart Z
   \tag*{\eqref{eq:trace-scalar}}\\
 \repart Z&=\tfrac12(Z+Z^{\mathsf H})
   \tag*{\eqref{eq:algebra-re}}\\
 \impart Z&=(Z-Z^{\mathsf H})/(2\ii)
   \tag*{\eqref{eq:algebra-im}}\\
 \det Z&=\adj Z\,Z=S^2-\bV^2
   \tag*{\eqref{eq:det}}\\
 Z^{-1}&=\adj Z/\det Z\qquad(\det Z\ne0)
   \tag*{\eqref{eq:inverse}}\\
 \lVert Z\rVert^2&=\scpart{Z^{\mathsf H}Z}
   \tag*{\eqref{eq:hermitian-norm}}.
\end{align*}
The coefficient form is \eqref{eq:coefficient-norm}; product-order and
cyclicity rules are \eqref{eq:hermitian-product},
\eqref{eq:sigma-conjugate-product}, and \eqref{eq:scalar-cyclic-three}.
Projectors and the spectral rule are \eqref{eq:projectors} and
\eqref{eq:complex-function-projectors}.  Matrix and quaternion images are
\eqref{eq:matrix} and \eqref{eq:quaternion-matrix}.

\Needspace{8\baselineskip}
\paragraph{Rotations and boosts.}
Here $\bu$ is a \hyperref[def:real-unit-vector]{real unit axis} and $\theta$ a real angle or rapidity;
$\beta=\tanh\theta$ and $\gamma=\cosh\theta$ in the boost formulas.
\begin{align*}
 R&=\exp(\ii\theta\bu\cdot\bsig/2)
   =\cos(\theta/2)+\ii\sin(\theta/2)\bu\cdot\bsig
   \tag*{\eqref{eq:rotation-exponential}}\\
 L&=\exp(\theta\bu\cdot\bsig/2)
   =\cosh(\theta/2)+\sinh(\theta/2)\bu\cdot\bsig
   \tag*{\eqref{eq:boost-exponential}}\\
 L&=\sqrt{\frac{\gamma+1}{2}}
     \left(1+\frac{\gamma\beta}{\gamma+1}\bu\cdot\bsig\right)
   \tag*{\eqref{eq:boost-gamma-form}}\\
 L^2&=\gamma(1+\beta\bu\cdot\bsig)
   \tag*{\eqref{eq:boost-square}}\\
 Z'&=R^{\mathsf H}ZR
   \tag*{\eqref{eq:rotation-action}}\\
 Z'&=L^{-1}ZL^{-1}
   \tag*{\eqref{eq:boost-action}}.
\end{align*}
The last line is the passive boost of a real four-vector.
Rodrigues' rotation is \eqref{eq:rodrigues}; boost components are
\eqref{eq:boost-scalar-component}--\eqref{eq:boost-vector-component}.
The general determinant-one congruence is
\eqref{eq:general-lorentz-congruence}; physical field laws are
in Section~\ref{sec:transformations}.

\Needspace{8\baselineskip}
\paragraph{Spacetime and particles.}
Recall real event $X$, proper time $s$, coordinate velocity $\bv$, and
constant mass $m>0$.  The real proper velocity $U=\dd X/\dd s$ is a
\hyperref[eq:unit-paravector]{unit paravector} with $UU^{-1}=1$.
The real four-momentum $P$ has mechanical energy
$E$ (including rest energy) and spatial momentum $\bp$.
\begin{align*}
 \partial&=\partial_t+\bsig\cdot\pdr
   \tag*{\eqref{eq:paravector-derivative}}\\
 \Box&=\partial^*\partial=\partial_t^2-\pdr^2
   \tag*{\eqref{eq:dalembertian}}\\
 \dd s^2&=\det(\dd X)=\dd t^2-\dd\br^2
   \tag*{\eqref{eq:interval}}\\
 U&=\dd X/\dd s=\gamma(1+\bv\cdot\bsig)
   \tag*{\eqref{eq:proper-velocity}}\\
 U^*U&=1
   \tag*{\eqref{eq:proper-velocity-star-normalization}}\\
 P&=mU=E+\bp\cdot\bsig
   \tag*{\eqref{eq:four-momentum}}\\
 E^2-\bp^2&=m^2
   \tag*{\eqref{eq:mass-shell}}.
\end{align*}
Here $\gamma=\dd t/\dd s=(1-\bv^2)^{-1/2}$.  Proper velocity is real,
with \hyperref[eq:unit-determinant-paravector]{determinant one}; Hermitian unitarity is the separate condition
\eqref{eq:unitary-transform}.
For real proper times $s_0<s_1$ and coordinate times $t_0<t$,
\begin{align*}
 X(s_1)-X(s_0)&=\int_{s_0}^{s_1}U(s)\dd s
   \tag*{\eqref{eq:worldline-integral}}\\
 s(t)-s(t_0)&=\int_{t_0}^{t}\frac{\dd t'}{\gamma(t')}
 =\int_{t_0}^{t}\sqrt{1-\bv(t')^2}\dd t'
   \tag*{\eqref{eq:elapsed-proper-time}}.
\end{align*}
Clock comparison is \eqref{eq:maximum-proper-time}; constant proper
acceleration gives \eqref{eq:hyperbolic-worldline}; the future cone is
\eqref{eq:future-light-cone}.  Proper acceleration, inertia, and particle
action are \eqref{eq:proper-acceleration-components},
\eqref{eq:relativistic-inertia}, and \eqref{eq:free-particle-action}.
Moving-frame velocity and its causal bound are
\eqref{eq:velocity-boost}--\eqref{eq:velocity-causal-bound}; clock dilation
and measured contraction are \eqref{eq:time-dilation} and
\eqref{eq:length-contraction}; energy and momentum transform by
\eqref{eq:energy-boost}--\eqref{eq:momentum-boost}.

\Needspace{8\baselineskip}
\paragraph{Electromagnetism.}
Recall the real proper velocity $U=\dd X/\dd s$, a
\hyperref[eq:unit-paravector]{unit paravector} with $UU^{-1}=1$.
Recall real potentials $V,\bA$, real electric and magnetic fields
$\bE,\bBf$, the complex field $F$, real potential paravector $\Phi$,
and real charge-current paravector $D$.  The real gauge scalar is $S$;
$q$ is charge.  These fields are functions of real spacetime.
Their boosted arguments and components are
\eqref{eq:boosted-field-arguments} and
\eqref{eq:electric-boost}--\eqref{eq:magnetic-boost}.
\begin{align*}
 F&=(\bE+\ii\bBf)\cdot\bsig
   \tag*{\eqref{eq:em-field-definition}}\\
 \partial F&=D^*
   \tag*{\eqref{eq:maxwell}}\\
 \Box F&=\partial^*D^*
   \tag*{\eqref{eq:maxwell-wave}}\\
 \partial\Phi&=S+F^*
   \tag*{\eqref{eq:potential-field}}\\
 F^2&=\bE^2-\bBf^2+2\ii\bE\cdot\bBf
   \tag*{\eqref{eq:field-invariants}}\\
 \frac{\dd P}{\dd s}&=q\repart{FU}
   \tag*{\eqref{eq:lorentz}}\\
 T(N)&=\tfrac12FNF^{\mathsf H}
   \tag*{\eqref{eq:electromagnetic-stress-map}}\\
 \partial_tu+\pdr\cdot\bS&=-\bJ\cdot\bE
   \tag*{\eqref{eq:poynting}}.
\end{align*}
Here $N$ is a real paravector argument of the stress map, $u$ field energy
density, $\bS$ energy flux, and $\bJ$ electric current density.
Potential components are \eqref{eq:electric-potential}--\eqref{eq:magnetic-potential};
force components are \eqref{eq:lorentz-momentum} and \eqref{eq:lorentz-energy}.
The action derivation uses \eqref{eq:maxwell-lagrangian} and
\eqref{eq:charged-particle-action}.

\Needspace{8\baselineskip}
\paragraph{Relativistic waves.}
Recall real mass $m>0$, charge $q$, and potentials $V,\bA$.
Their real potential paravector is $\Phi=V+\bA\cdot\bsig$.
Here $E,\bp$ are mechanical differential operators, $\psi$ a complex
scalar KG field, and $\bPsi$ a two-entry complex paravector Dirac field.
\begin{align*}
 E&=\ii\partial_t-qV
   \tag*{\eqref{eq:mechanical-energy-operator}}\\
 \bp&=-\ii\pdr-q\bA
   \tag*{\eqref{eq:mechanical-momentum-operator}}\\
 (\bp\cdot\bsig)^2&=\bp^2-q\bBf\cdot\bsig
   \tag*{\eqref{eq:kinetic-sigma-square}}\\
 [E,\bp\cdot\bsig]&=\ii q\bE\cdot\bsig
   \tag*{\eqref{eq:kinetic-electric-commutator}}\\
 (-E^2+\bp^2+m^2)\psi&=0
   \tag*{\eqref{eq:kg-field}}\\
 \Wop(Z)&=\begin{Bmatrix}0&Z\\\adj Z&0\end{Bmatrix}
   \tag*{\eqref{eq:dirac-block-map}}\\
 (\Wop(\ii\partial-q\Phi^*)-m\bI)\bPsi&=0.
   \tag*{\eqref{eq:dirac-field}}
\end{align*}
Here $Z$ is a complex paravector and $\bI$ the block identity.
The free equations are \eqref{eq:kg} and \eqref{eq:dirac};
the two-branch solution is \eqref{eq:kg-general-solution}.
Massive Dirac reconstruction is given by \eqref{eq:kg-dirac-complete-lift}
and \eqref{eq:dirac-field-reconstruction}; the squared interacting entries
are \eqref{eq:dirac-squared-first} and \eqref{eq:dirac-squared-field}.

\Needspace{12\baselineskip}
\paragraph{Low-energy limits.}
After removing rest energy by \eqref{eq:dirac-rest-phase} and
\eqref{eq:kg-rest-phase}, the complex scalar envelope $\psi$ obeys
Schr\"odinger's equation to leading order:
\begin{equation*}
 \ii\partial_t\psi=\left(qV+\frac{\bp^2}{2m}\right)\psi.
 \tag*{\eqref{eq:schrodinger}}
\end{equation*}
The complex paravector envelope $\phi$ obeys Pauli's equation:
\begin{equation*}
 \ii\partial_t\phi=\left(qV+
 \frac{\bp^2-q\bBf\cdot\bsig}{2m}\right)\phi.
 \tag*{\eqref{eq:pauli}}
\end{equation*}
The recursion and approximation conditions are in
Section~\ref{sec:low-energy-limits}; these are the same leading order
for scalar and spinor amplitudes respectively.

\Needspace{10\baselineskip}
\paragraph{Spin and probabilities.}
Recall a nonzero complex paravector amplitude $\psi$, real norm density
$\rho>0$, real state paravector $B$, and
\hyperref[def:real-unit-vector]{real unit measurement axis} $\bu$.
The real observable $Z$ acts from the left.
\begin{align*}
 \langle\psi'\mid\psi\rangle&=\scpart{\psi'^{\mathsf H}\psi}
   \tag*{\eqref{eq:paravector-inner-product}}\\
 \rho&=\scpart{\psi^{\mathsf H}\psi}
   \tag*{\eqref{eq:probability-density}}\\
 B&=\psi\psi^{\mathsf H}/\rho=1+\bm s\cdot\bsig
   \tag*{\eqref{eq:spin-state}}\\
 p_\pm(\bu)&=\tfrac12(1\pm\bu\cdot\bm s)
   \tag*{\eqref{eq:spin-probability-intrinsic}}\\
 \langle Z\rangle_\psi&=\scpart{\psi^{\mathsf H}Z\psi}/\rho.
   \tag*{\eqref{eq:paravector-expectation}}
\end{align*}
Here $\psi'$ is a second paravector amplitude and $\bm s$ the real polarization.
The matrix verification is \eqref{eq:matrix-state}--\eqref{eq:matrix-born};
spin uncertainty is \eqref{eq:spin-uncertainty}.

\endgroup
}

\appendix

\section{Pauli-matrix representation}
\label{sec:matrix}

This appendix translates the established paravector operations into
matrices.  Let $Z=S+\bV\cdot\bsig$, with $S\in\C$ and $\bV\in\C^3$.
Write $\bM=\operatorname{mat}(Z)$, a $2\times2$ matrix with entries in $\C$.
The identity matrices are $\bI_2$ and $\bI_4$ in their stated dimensions.

\subsection{Faithful representation}

The Pauli images of the generators are
\[
 \operatorname{mat}(\sigma_1)=\begin{Bmatrix}0&1\\1&0\end{Bmatrix},\quad
 \operatorname{mat}(\sigma_2)=\begin{Bmatrix}0&-\ii\\\ii&0\end{Bmatrix},\quad
 \operatorname{mat}(\sigma_3)=\begin{Bmatrix}1&0\\0&-1\end{Bmatrix}.
\]
Consequently,
\begin{equation}
 \operatorname{mat}(S+\bV\cdot\bsig)=
 \begin{Bmatrix}S+V_3&V_1-\ii V_2\\V_1+\ii V_2&S-V_3\end{Bmatrix}.
 \label{eq:matrix}
\end{equation}
The generator images obey the sigma product rule and are linearly independent over
$\C$ together with $\bI_2$.  The map is therefore faithful, preserves
addition, sends $1$ to $\bI_2$, and satisfies
\begin{equation}
 \operatorname{mat}(Z'Z)=\operatorname{mat}(Z')\operatorname{mat}(Z).
 \label{eq:matrix-product}
\end{equation}
Here $Z'$ is any second complex paravector.  Thus algebra multiplication
is ordinary matrix multiplication.

\subsection{Conjugations and product order}

For complex matrices $\bM,\bN$, transpose $^{\mathsf T}$ exchanges row
and column indices; entrywise complex conjugation $^{\mathsf C}$ sends each
$a+\ii b$ to $a-\ii b$ for real $a,b$.  Their composition is
$\bM^{\mathsf H}=(\bM^{\mathsf C})^{\mathsf T}$.
Introduce the fixed real matrix
\[
 \bm S=\begin{Bmatrix}0&1\\-1&0\end{Bmatrix},\qquad
 \bm S^{-1}=-\bm S.
\]
The paravector sigma conjugation has the matrix image
\begin{equation}
 \operatorname{mat}(Z^*)=\bm S^{-1}\bM^{\mathsf C}\bm S=: \bM^*.
 \label{eq:matrix-sigma-conjugation}
\end{equation}
Substituting the basis in \eqref{eq:matrix-sigma-conjugation} gives
\begin{center}
\begin{tabular}{c c c c c}
\toprule
Algebra basis & $1$ & $\ii$ & $\sigma_k$ & $\ii\sigma_k$\\
\midrule
Hermitian image & $\bI_2$ & $-\ii\bI_2$
 & $\operatorname{mat}(\sigma_k)$ & $-\ii\operatorname{mat}(\sigma_k)$\\
Sigma-conjugate image & $\bI_2$ & $-\ii\bI_2$
 & $-\operatorname{mat}(\sigma_k)$ & $\ii\operatorname{mat}(\sigma_k)$\\
\bottomrule
\end{tabular}
\end{center}
Here $k=1,2,3$.  These are exactly the operations of
Section~\ref{sec:conjugate-operators}, so
$\operatorname{mat}(Z^{\mathsf H})=\bM^{\mathsf H}$.
Their product-order rules are
\begin{align*}
 (\bM\bN)^{\mathsf T}&=\bN^{\mathsf T}\bM^{\mathsf T},\\
 (\bM\bN)^{\mathsf C}&=\bM^{\mathsf C}\bN^{\mathsf C},\\
 (\bM\bN)^{\mathsf H}&=\bN^{\mathsf H}\bM^{\mathsf H}.
\end{align*}
For sigma conjugation the middle factors cancel:
\[
 (\bM\bN)^*=\bm S^{-1}\bM^{\mathsf C}\bm S
                   \bm S^{-1}\bN^{\mathsf C}\bm S.
\]
Thus sigma conjugation preserves product order:
\begin{equation}
 (\bM\bN)^*=\bM^*\bN^*.
 \label{eq:matrix-star-product}
\end{equation}
This verifies \eqref{eq:sigma-conjugate-product} in matrices.
Hermitian and sigma conjugation commute, and applying either twice
restores the original matrix.
The matrix-only operations can be recovered as
\begin{align*}
 \bM^{\mathsf C}&=\bm S\bM^*\bm S^{-1},\\
 \bM^{\mathsf T}&=\bm S(\bM^*)^{\mathsf H}\bm S^{-1}.
\end{align*}

\subsection{Scalar part, determinant, and norm}

For $\bM=(m_{jk})$, its trace is $m_{11}+m_{22}$ and its determinant
is $m_{11}m_{22}-m_{12}m_{21}$.  Substituting \eqref{eq:matrix} gives
\begin{align}
 \scpart Z&=\tfrac12\operatorname{tr}\bM=S,
 \label{eq:matrix-scalar-part}\\
 (\vecpart Z)_k&=\tfrac12\operatorname{tr}
                    (\operatorname{mat}(\sigma_k)\bM)=V_k,\notag\\
 \det Z&=\det\bM=S^2-\bV^2.
 \label{eq:matrix-determinant}
\end{align}
By \eqref{eq:matrix-scalar-part}, matrix-trace cyclicity is precisely the scalar-part cyclicity used
in the main text.  For $\bM'=\operatorname{mat}(Z')$,
\[
 \scpart{Z'Z}=\tfrac12\operatorname{tr}(\bM'\bM)
             =\scpart{ZZ'}.
\]
The remaining extraction and norm formulas are
\begin{align}
 \operatorname{mat}(\vecpart Z\cdot\bsig)&=\bM-\scpart Z\,\bI_2,\notag\\
 \operatorname{mat}(\repart Z)&=\tfrac12(\bM+\bM^{\mathsf H}),
 \label{eq:matrix-real-part}\\
 \operatorname{mat}(\impart Z)&=(\bM-\bM^{\mathsf H})/(2\ii),
 \label{eq:matrix-imaginary-part}\\
 \lVert Z\rVert^2&=\tfrac12\operatorname{tr}(\bM^{\mathsf H}\bM)
                  =\lVert S\rVert^2+\lVert\bV\rVert^2.
 \label{eq:matrix-norm}
\end{align}
Equations~\eqref{eq:matrix-real-part}--\eqref{eq:matrix-imaginary-part}
extract Hermitian matrices, so a real paravector has a Hermitian image.
The norm in \eqref{eq:matrix-norm} is half the squared Frobenius norm;
the determinant in \eqref{eq:matrix-determinant} may be complex.

The matrix adjugate is
\begin{equation}
 \adj\bM=(\operatorname{tr}\bM)\bI_2-\bM
        =\operatorname{mat}(\adj Z)=(\bM^*)^{\mathsf H}.
 \label{eq:matrix-adjugate}
\end{equation}
It satisfies
\begin{equation*}
 \operatorname{mat}(\adj Z\,Z)=\det(\bM)\bI_2.
\end{equation*}
Consequently $\bM^{-1}=\adj\bM/\det\bM$ when the determinant is
nonzero.  The determinant product and adjugate order rules are the same
in both representations.  The quadratic identity~\eqref{eq:minimal} is
precisely Cayley--Hamilton for a $2\times2$ matrix.

\subsection{Matrix amplitudes and column spinors}
\label{sec:matrix-spinors}

Let $\psi,\psi'$ be complex paravector amplitudes, with matrix images
$\bM_\psi,\bM_{\psi'}$, and let $Z=Z^{\mathsf H}$ be a real observable.
For fields these depend on real $(t,\br)$.  The inner products
\eqref{eq:paravector-inner-product} and
\eqref{eq:paravector-matrix-element} become
\begin{align*}
 \langle\psi'\mid\psi\rangle
 &=\tfrac12\operatorname{tr}(\bM_{\psi'}^{\mathsf H}\bM_\psi),\\
 \langle\psi'\mid Z\mid\psi\rangle
 &=\tfrac12\operatorname{tr}
    (\bM_{\psi'}^{\mathsf H}\operatorname{mat}(Z)\bM_\psi).
\end{align*}

\paragraph{An arbitrary spin basis.}
Recall the arbitrary constant perpendicular
\hyperref[def:real-unit-vector]{real unit vectors}
$\bu,\ba\in\R^3$ of \eqref{eq:axis-ideal-basis}.
Set $\Pi_+=(1+\bu\cdot\bsig)/2$ and choose a constant complex
paravector $\psi'$ with $\Pi_+\psi'\ne0$.  Then
\begin{equation}
 \begin{Bmatrix}\psi_+\\\psi_-\end{Bmatrix}
 :=\frac1{\lVert\Pi_+\psi'\rVert}
 \begin{Bmatrix}\Pi_+\psi'\\(\ba\cdot\bsig)\Pi_+\psi'\end{Bmatrix}.
 \label{eq:ideal-spin-basis}
\end{equation}
This constructs \eqref{eq:axis-ideal-basis} with
$(\bu\cdot\bsig)\psi_\pm=\pm\psi_\pm$ and
$\langle\psi_s\mid\psi_{s'}\rangle=\delta_{ss'}$ for $s,s'=+,-$.
The two complex scalar components give
\begin{equation}
 \psi=c_+\psi_++c_-\psi_-=(c_++c_-\ba\cdot\bsig)\psi_+.
 \label{eq:ideal-components}
\end{equation}
For fields, $c_\pm:\R^4\to\C$, and
\eqref{eq:probability-density} is
$\rho=\lVert c_+\rVert^2+\lVert c_-\rVert^2$.

\paragraph{The two-component column.}
Every column of $\bM_{\psi_+}$ has eigenvalue $+1$ under
$\operatorname{mat}(\bu\cdot\bsig)$, whose $+1$ eigenspace is
one-dimensional.  Since $\lVert\psi_+\rVert=1$, its matrix therefore
factors as $\bM_{\psi_+}=\sqrt2\,\bpsi_+\bm v^{\mathsf H}$ for
constant columns $\bpsi_+,\bm v\in\C^2$ satisfying
$\bpsi_+^{\mathsf H}\bpsi_+=\bm v^{\mathsf H}\bm v=1$.
Left multiplication gives
$\bM_{\psi_-}=\sqrt2\,\bpsi_-\bm v^{\mathsf H}$, where
$\bpsi_-=\operatorname{mat}(\ba\cdot\bsig)\bpsi_+$.
The columns $\bpsi_+,\bpsi_-$ are orthonormal.  Thus every amplitude
in \eqref{eq:ideal-components} has the conventional Pauli column
\begin{equation}
 \bpsi:=\frac1{\sqrt2}\bM_\psi\bm v
 =c_+\bpsi_++c_-\bpsi_-.
 \label{eq:matrix-ideal-column}
\end{equation}
Conversely, $\bM_\psi=\sqrt2\,\bpsi\bm v^{\mathsf H}$, so this
correspondence is one-to-one.  For two amplitudes in this fixed basis,
substitution cancels the factor of two from the trace pairing:
\begin{align*}
 \langle\psi'\mid\psi\rangle&=\bpsi'^{\mathsf H}\bpsi,\\
 \langle\psi'\mid Z\mid\psi\rangle
 &=\bpsi'^{\mathsf H}\operatorname{mat}(Z)\bpsi.
\end{align*}
Left action $T\psi$ maps to $\operatorname{mat}(T)\bpsi$ for any complex paravector $T$.

Let $\bm C$ be the $2\times2$ matrix with columns
$\bpsi_+,\bpsi_-$.  Then $\bm C^{\mathsf H}\bm C=\bI_2$ and
$\bpsi=\bm C\{c_+,c_-\}^{\mathsf T}$.
For a complex paravector $Z=S+\bV\cdot\bsig$, its matrix in this basis is
\begin{equation}
 \operatorname{mat}'(Z):=\bm C^{\mathsf H}\operatorname{mat}(Z)\bm C
 =\begin{Bmatrix}
 S+\bu\cdot\bV&(\ba-\ii\bu\times\ba)\cdot\bV\\
 (\ba+\ii\bu\times\ba)\cdot\bV&S-\bu\cdot\bV
 \end{Bmatrix}.
 \label{eq:matrix-arbitrary-axis}
\end{equation}
Here the prime denotes only this change of matrix basis.
For nonzero norm $\rho=\bpsi^{\mathsf H}\bpsi$, the real state
paravector $B$ and projectors $\Pi_\pm=(1\pm\bu\cdot\bsig)/2$
translate to
\begin{align}
 \tfrac12\operatorname{mat}(B)&=\bpsi\bpsi^{\mathsf H}/\rho,
 \label{eq:matrix-state}\\
 p_\pm(\bu)&=\bpsi^{\mathsf H}\operatorname{mat}(\Pi_\pm)\bpsi/\rho.
 \label{eq:matrix-born}
\end{align}
Equations~\eqref{eq:matrix-state} and \eqref{eq:matrix-born} are the ordinary
density matrix and Born rule.

\subsection{Dirac columns}
\label{sec:matrix-dirac-columns}

For a Dirac block $\bPsi=\{\Psi_1,\Psi_2\}^{\mathsf T}$, let
$\bpsi_1,\bpsi_2\in\C^2$ be the corresponding columns.  Then
\[
 \Wop(Z)\bPsi\quad\longleftrightarrow\quad
 \begin{Bmatrix}0&\operatorname{mat}(Z)\\
 \operatorname{mat}(\adj Z)&0\end{Bmatrix}
 \begin{Bmatrix}\bpsi_1\\\bpsi_2\end{Bmatrix}.
\]
The block norm in \eqref{eq:dirac-norm} becomes the ordinary four-component norm:
\begin{equation}
 \rho=\bpsi_1^{\mathsf H}\bpsi_1+\bpsi_2^{\mathsf H}\bpsi_2.
 \label{eq:matrix-dirac-density}
\end{equation}

\paragraph{Characteristic polynomial.}
On one Fourier mode, $s\in\C$ is the spectral parameter and $\Box$
its scalar eigenvalue.  The two off-diagonal blocks multiply to
$-\Box\bI_2$, so block elimination gives the ordinary matrix determinant
\begin{equation}
 \begin{aligned}
 \det\begin{Bmatrix}
 -(s+m)\bI_2&\operatorname{mat}(\ii\partial)\\
 \operatorname{mat}(\ii\partial^*)&-(s+m)\bI_2
 \end{Bmatrix}
 &=\det\!\left(\bigl((s+m)^2+\Box\bigr)\bI_2\right)\\
 &=\bigl((s+m)^2+\Box\bigr)^2.
 \end{aligned}
 \label{eq:matrix-dirac-characteristic}
\end{equation}
This verifies \eqref{eq:dirac-characteristic}, including multiplicities.
The polynomial identity also holds at $s=-m$.

\Needspace{14\baselineskip}
\subsection{Four components and the standard Dirac basis}
\label{sec:matrix-dirac-components}

In the arbitrary-axis basis \eqref{eq:ideal-spin-basis}, write
$\Psi_1=a\psi_++b\psi_-$ and $\Psi_2=c\psi_++d\psi_-$,
where $a,b,c,d:\R^4\to\C$ are the four scalar components.
Their columns are $\bpsi_1'=\{a,b\}^{\mathsf T}$ and
$\bpsi_2'=\{c,d\}^{\mathsf T}$; here the prime means
$\bpsi_k'=\bm C^{\mathsf H}\bpsi_k$.
Recall real mass $m>0$, real charge $q$, and
real potentials $V:\R^4\to\R$ and $\bA:\R^4\to\R^3$.
The mechanical energy and momentum operators are $E=\ii\partial_t-qV$ and
$\bp=-\ii\pdr-q\bA$.  Their momentum components in this frame are
\[
 p_1=\ba\cdot\bp,\qquad p_2=(\bu\times\ba)\cdot\bp,
 \qquad p_3=\bu\cdot\bp.
\]
Using \eqref{eq:matrix-arbitrary-axis}, the coupled paravector equations give
\begin{align}
 (E-p_3)c-(p_1-\ii p_2)d&=ma,
 \label{eq:matrix-dirac-a}\\
 -(p_1+\ii p_2)c+(E+p_3)d&=mb,
 \label{eq:matrix-dirac-b}\\
 (E+p_3)a+(p_1-\ii p_2)b&=mc,
 \label{eq:matrix-dirac-c}\\
 (p_1+\ii p_2)a+(E-p_3)b&=md.
 \label{eq:matrix-dirac-d}
\end{align}
Operators in \eqref{eq:matrix-dirac-a}--\eqref{eq:matrix-dirac-d}
act on the scalar immediately to their right.  Their Hamiltonian
on $\{a,b,c,d\}^{\mathsf T}$ is
\begin{equation}
 H=\begin{Bmatrix}
 qV-p_3&-p_1+\ii p_2&m&0\\
 -p_1-\ii p_2&qV+p_3&0&m\\
 m&0&qV+p_3&p_1-\ii p_2\\
 0&m&p_1+\ii p_2&qV-p_3
 \end{Bmatrix}.
 \label{eq:matrix-dirac-hamiltonian}
\end{equation}
This is the component image of \eqref{eq:dirac-hamiltonian}.

\Needspace{10\baselineskip}
To identify conventional gamma matrices, put
$\bm\sigma_k=\operatorname{mat}(\sigma_k)$ for $k=1,2,3$.
These are also the images under $\operatorname{mat}'$ of
$\ba\cdot\bsig$, $(\bu\times\ba)\cdot\bsig$, and $\bu\cdot\bsig$,
respectively.  Their numerical entries apply to every choice of $\bu,\ba$.
Define
\begin{align}
 \gamma^0&=\begin{Bmatrix}0&\bI_2\\\bI_2&0\end{Bmatrix},
 \label{eq:gamma-time}\\
 \gamma^k&=\begin{Bmatrix}0&\bm\sigma_k\\-\bm\sigma_k&0\end{Bmatrix}.
 \label{eq:gamma-space}
\end{align}
Equations~\eqref{eq:gamma-time}--\eqref{eq:gamma-space} obey
$(\gamma^0)^2=\bI_4$, $(\gamma^k)^2=-\bI_4$, and
$\gamma^j\gamma^k+\gamma^k\gamma^j=0$ for distinct
$j,k=0,1,2,3$.  Thus, in this basis,
\begin{equation}
 \Wop(\ii\partial-q\Phi^*)
 \quad\longleftrightarrow\quad
 \gamma^0E-\sum_{k=1}^3\gamma^kp_k,
 \label{eq:matrix-covariant-dirac}
\end{equation}
which is the standard chiral Dirac operator with signature
$(+,-,-,-)$~\cite[p.~83]{tong-qft}.  Here $\Phi=V+\bA\cdot\bsig$ is the
real potential paravector.

The unitary change to the standard Dirac basis commutes with the scalar
rest-phase factor explained in \eqref{eq:dirac-rest-phase}.
Together they define the envelope columns
\begin{equation}
 \begin{Bmatrix}\bm\phi\\\bm\chi\end{Bmatrix}
 =\frac{e^{\ii mt}}{\sqrt2}
 \begin{Bmatrix}\bI_2&\bI_2\\-\bI_2&\bI_2\end{Bmatrix}
 \begin{Bmatrix}\bpsi_1'\\\bpsi_2'\end{Bmatrix}.
 \label{eq:matrix-envelope-map}
\end{equation}
Here $\bm\phi,\bm\chi\in\C^2$ are the coefficient columns in the same
spin basis of the paravector fields
$\phi=e^{\ii mt}(\Psi_1+\Psi_2)/\sqrt2$ and
$\chi=e^{\ii mt}(\Psi_2-\Psi_1)/\sqrt2$.
Before removing the rest phase, the basis change in
\eqref{eq:matrix-envelope-map} sends \eqref{eq:matrix-dirac-hamiltonian} to
\begin{equation}
 H'=\begin{Bmatrix}
 (qV+m)\bI_2&\sum_kp_k\bm\sigma_k\\
 \sum_kp_k\bm\sigma_k&(qV-m)\bI_2
 \end{Bmatrix}.
 \label{eq:matrix-standard-dirac}
\end{equation}
Multiplying the full column by $e^{\ii mt}$ subtracts $m$ from its
Hamiltonian: $\ii\partial_t(e^{\ii mt}\bpsi)=e^{\ii mt}(\ii\partial_t-m)\bpsi$.
Thus \eqref{eq:matrix-standard-dirac} becomes $H'-mI_4$:
its diagonal entries are $qV$ and $qV-2m$.
Moving $qV$ into $E=\ii\partial_t-qV$ gives the exact envelope equations
\begin{align*}
 E\bm\phi&=\sum_kp_k\bm\sigma_k\bm\chi,\\
 (E+2m)\bm\chi&=\sum_kp_k\bm\sigma_k\bm\phi.
\end{align*}
These identify the large and small amplitudes used in the Pauli limit.

For real electric and magnetic components $E_k,B_k$ in the same frame,
the squared Dirac field coupling has the explicit image
\begin{align}
 \operatorname{mat}'(\ii qF)
 &=q\begin{Bmatrix}
 \ii E_3-B_3&\ii E_1+E_2-B_1+\ii B_2\\
 \ii E_1-E_2-B_1-\ii B_2&-\ii E_3+B_3
 \end{Bmatrix},
 \label{eq:matrix-field-coupling}\\
 \operatorname{mat}'(\ii qF^*)&=\operatorname{mat}'(\ii qF)^{\mathsf H}.
 \notag
\end{align}
The indexed $E_k$ are field components;
unindexed $E$ remains the mechanical energy operator.
The two couplings are $-q\bBf\cdot\bsig\pm\ii q\bE\cdot\bsig$.

\subsection{Quaternion subalgebra}
\label{sec:quaternions}

Define $q_k=-\ii\sigma_k$ for $k=1,2,3$, and
$\bq=\{q_1,q_2,q_3\}$.  These are bivectors, not electric charges.
Their product follows from the sigma rule:
\begin{equation}
 q_jq_k=-\delta_{jk}+\sum_{n=1}^3\epsilon_{jkn}q_n.
 \label{eq:quaternion-product}
\end{equation}
In particular,
$q_k^2=-1$, $q_1q_2=q_3$, and $q_1q_2q_3=-1$.
For real $a$ and $\bA\in\R^3$, the quaternion
\begin{equation}
 Q=a+\bA\cdot\bq=a-\ii\bA\cdot\bsig
 \label{eq:quaternion-paravector}
\end{equation}
is a complex paravector with real quaternion coordinates.
These elements form the real scalar--bivector
subalgebra.  For a second such element
$Q'=a'+\bA'\cdot\bq$, expansion gives
\[
 QQ'=(aa'-\bA\cdot\bA')
 +(a\bA'+a'\bA+\bA\times\bA')\cdot\bq.
\]
Because $q_k^{\mathsf H}=-q_k$ and $q_k^*=q_k$, quaternion conjugation
is the existing Hermitian conjugation:
\begin{align}
 Q^{\mathsf H}&=a-\bA\cdot\bq,\notag\\
 Q^*&=Q,\notag\\
 \adj Q&=Q^{\mathsf H},\notag\\
 QQ^{\mathsf H}&=Q^{\mathsf H}Q=a^2+\bA^2=\det Q.
 \label{eq:quaternion-norm}
\end{align}
For nonzero $Q$, \eqref{eq:quaternion-norm} gives the inverse
$Q^{\mathsf H}/(a^2+\bA^2)$.
The explicit matrix image is
\begin{equation}
 \operatorname{mat}(Q)=
 \begin{Bmatrix}a-\ii A_3&-A_2-\ii A_1\\
                 A_2-\ii A_1&a+\ii A_3\end{Bmatrix}.
 \label{eq:quaternion-matrix}
\end{equation}
Here
$\operatorname{mat}(q_k)=-\ii\operatorname{mat}(\sigma_k)$.
The conjugation and norm identities become
\begin{align*}
 \operatorname{mat}(Q^{\mathsf H})&=\operatorname{mat}(Q)^{\mathsf H},\\
 \operatorname{mat}(Q)^{\mathsf H}\operatorname{mat}(Q)
 &=(a^2+\bA^2)\bI_2,\\
 \det(\operatorname{mat}(Q))&=a^2+\bA^2.
\end{align*}
A unit quaternion means $a^2+\bA^2=1$: it is a
\hyperref[eq:unit-norm-paravector]{paravector of unit norm} and corresponds
to a unitary matrix of determinant one.  For
\hyperref[def:real-unit-vector]{real unit vector} $\bu$ and angle $\theta$, let
\[
 Q=\cos(\theta/2)+\bu\cdot\bq\sin(\theta/2).
\]
Then $Q=R^{\mathsf H}$ for $R=\exp(\ii\theta\bu\cdot\bsig/2)$.
For a real vector $\br$, the quaternion rotation
\begin{equation}
 \br'\cdot\bq=Q(\br\cdot\bq)Q^{\mathsf H}
 \label{eq:quaternion-rotation}
\end{equation}
is Rodrigues' formula~\eqref{eq:rodrigues} in quaternion form.

\begin{thebibliography}{99}
\addcontentsline{toc}{section}{References}
\bibitem{baylis2004}
W. E. Baylis and G. Sobczyk,
``Relativity in Clifford's geometric algebras of space and spacetime,''
\emph{Int. J. Theor. Phys.}, vol.~43, pp.~2061--2079, 2004,
doi: \href{https://doi.org/10.1023/B:IJTP.0000049010.53558.b7}{10.1023/B:IJTP.0000049010.53558.b7}.
[Online]. Available: \url{https://arxiv.org/pdf/math-ph/0405026}.

\bibitem{tong-relativity}
D. Tong, \emph{Dynamics and Relativity}, University of Cambridge,
Cambridge, U.K., lecture notes, 2013.
[Online]. Available: \url{https://www.damtp.cam.ac.uk/user/tong/relativity/dynrel.pdf}.

\bibitem{tong-em}
D. Tong, \emph{Electromagnetism}, University of Cambridge,
Cambridge, U.K., lecture notes, 2015.
[Online]. Available: \url{https://www.damtp.cam.ac.uk/user/tong/em/electro.pdf}.

\bibitem{debroglie1929}
L. de Broglie, ``The wave nature of the electron,'' Nobel Lecture,
Dec.~12, 1929, in \emph{Nobel Lectures, Physics 1922--1941}.
Amsterdam, The Netherlands: Elsevier, 1965, pp.~244--256.
[Online]. Available: \url{https://www.nobelprize.org/uploads/2016/04/broglie-lecture.pdf}.

\bibitem{millikan1924}
R. A. Millikan, ``The electron and the light-quant from the experimental
point of view,'' Nobel Lecture, May~23, 1924,
in \emph{Nobel Lectures, Physics 1922--1941}.
Amsterdam, The Netherlands: Elsevier, 1965.
[Online]. Available: \url{https://www.nobelprize.org/uploads/2018/06/millikan-lecture.pdf}.

\bibitem{tong-qft}
D. Tong, \emph{Quantum Field Theory}, University of Cambridge,
Cambridge, U.K., lecture notes, 2006--2007.
[Online]. Available: \url{https://www.damtp.cam.ac.uk/user/tong/qft/qft.pdf}.

\bibitem{foldy-wouthuysen}
L. L. Foldy and S. A. Wouthuysen,
``On the Dirac theory of spin 1/2 particles and its non-relativistic limit,''
\emph{Phys. Rev.}, vol.~78, no.~1, pp.~29--36, Apr.~1950,
doi: \href{https://doi.org/10.1103/PhysRev.78.29}{10.1103/PhysRev.78.29}.

\bibitem{feynmanspin}
R. P. Feynman, R. B. Leighton, and M. Sands,
\emph{The Feynman Lectures on Physics}, vol.~III.
Reading, MA, USA: Addison-Wesley, 1965, ch.~6, ``Spin one-half.''
[Online]. Available: \url{https://www.feynmanlectures.caltech.edu/III_06.html}.

\end{thebibliography}

\clearpage
\referencesummary

\end{document}
